% La vie culturelle 1 (les arts et la littérature) — Allemand Tle A — Cours signé OnBuch+
% Programme officiel MINESEC. Compiler avec : tectonic AL22-la-vie-culturelle-1-les-arts-et-la-litterature.tex
\newcommand{\DOCMATIERE}{Allemand}
\newcommand{\DOCNIVEAU}{Tle A}
\newcommand{\DOCMODULE}{Chapitre 5 : Citoyenneté}
\newcommand{\DOCLECON}{AL22}
\newcommand{\DOCTITRE}{La vie culturelle 1 (les arts et la littérature)}
% OnBuch+ — préambule « cours signés » ENRICHI (compilé avec tectonic / XeLaTeX)
% Système visuel « Pop » d'OnBuch+ : fond crème (jamais blanc), bordures encre
% épaisses, ombres « dures » décalées, titres Archivo Black en capitales.
% Version MATHÉMATIQUES : graphes (pgfplots/tikz), boîtes pédagogiques
% (définition, propriété, méthode, attention, le savais-tu, exercice résolu,
% exercice, TP, bilan), page de garde et signature OnBuch+ ancrée en bas.
%
% Macros attendues dans main.tex AVANT \input{preamble} :
%   \DOCMATIERE  \DOCNIVEAU  \DOCMODULE  \DOCLECON (numéro)  \DOCTITRE
\documentclass[11pt]{article}

\usepackage{fontspec}
\usepackage[ngerman,french]{babel} % français = langue principale ; l'allemand via \all{..} / textebox
\usepackage[a4paper,margin=2cm,top=3.4cm,bottom=2.8cm,headheight=1.4cm,footskip=1.3cm]{geometry}
\usepackage{amsmath,amssymb,amsfonts}
\usepackage{mathtools} % \xmapsto, \coloneqq... souvent utilisés par les modèles
\usepackage{cancel} % simplifications barrées (bilans de forces)
% \abs{z} (module, valeur absolue) et \norm{u} (norme), utilisés par les modèles
\providecommand{\abs}{}\let\abs\relax\DeclarePairedDelimiter{\abs}{\lvert}{\rvert}
\providecommand{\norm}{}\let\norm\relax\DeclarePairedDelimiter{\norm}{\lVert}{\rVert}
\usepackage[table]{xcolor} % [table] : fournit aussi \rowcolors (alternance de lignes)
\usepackage{pagecolor}
\usepackage{tikz}
\usetikzlibrary{arrows.meta,positioning,calc,shapes.geometric,shapes.misc,shapes.arrows,decorations.pathmorphing,decorations.pathreplacing,decorations.markings,patterns,shadows,mindmap,trees,fit,backgrounds,matrix,intersections,angles,quotes}
\usepackage{pgfplots}
\usepackage[version=4]{mhchem} % équations nucléaires \ce{^{235}_{92}U}
\usepackage[european,siunitx]{circuitikz} % schémas électriques
\pgfplotsset{compat=1.18}
\usepgfplotslibrary{fillbetween,groupplots} % aires entre courbes (calcul intégral)
\usepackage{enumitem}
\usepackage{fancyhdr}
% [most] charge toutes les bibliothèques tcolorbox (listings, minted, raster,
% documentation...) et gonfle inutilement la mémoire TeX sur un document long
% (~300 boîtes) : on ne charge que ce qui est réellement utilisé.
\usepackage[skins,breakable]{tcolorbox}
\usepackage{graphicx}
\usepackage{colortbl}
\usepackage{booktabs}
\usepackage{tabularx}
\usepackage{diagbox} % case d'en-tête diagonale des tableaux à double entrée
% Colonnes extensibles centrée / gauche / droite (C, L, R), souvent utilisées par les modèles
\newcolumntype{C}{>{\centering\arraybackslash}X}
\newcolumntype{L}{>{\raggedright\arraybackslash}X}
\newcolumntype{R}{>{\raggedleft\arraybackslash}X}
\usepackage{array}
\usepackage{multicol}
\usepackage{multirow}
\usepackage{siunitx}
% parse-numbers=false : accepte aussi \SI{2,5 \times 10^{-3}}{...}
\sisetup{locale=FR,output-decimal-marker={,},group-minimum-digits=4,parse-numbers=false}
\DeclareSIUnit\ohm{\ensuremath{\symup{\Omega}}} % U+2126 absent de la police
\DeclareSIUnit\division{div} % réglages d'oscilloscope (ms/div, V/div)
\DeclareSIUnit\tr{tr} % tours (vitesse de rotation, ex. \SI{1500}{\tr\per\minute})
\DeclareSIUnit\molar{M} % concentration molaire (ex. \SI{3}{\molar}, \SI{1}{\micro\molar})
\DeclareSIUnit\an{an} % année (le modèle confond parfois avec \year, un compteur TeX) — ex. \SI{5}{\kilo\meter\per\an}
\DeclareSIUnit\FCFA{FCFA} % franc CFA (ex. \SI{500}{\FCFA\per\kilo\gram})
\DeclareSIUnit\degreeBrix{\textdegree Brix} % teneur en sucre d'un sirop/jus (réfractométrie)
\DeclareSIUnit\month{mois} % le modèle utilise parfois \month, un compteur TeX, comme unité de durée
\usepackage{qrcode}
% \degree : commande fréquemment utilisée par les modèles pour un angle ou une
% température en dehors de \SI{}{} (non fournie par les paquets chargés).
\providecommand{\degree}{^{\circ}}
% \qed (fin de démonstration), utilisé par les modèles sans amsthm
\providecommand{\qed}{\hfill$\square$}
% \barre : double barre des valeurs interdites dans un tableau de variations (tkz-tab)
\providecommand{\barre}{\ensuremath{\|}}
% environnement proof (amsthm non chargé) utilisé par les modèles
\ifdefined\proof\else\newenvironment{proof}[1][Démonstration]{\par\noindent\textit{#1.}\ }{\qed\par}\fi
\usepackage{lastpage}
\usepackage{hyperref}
\hypersetup{hidelinks}

% ============================================================
% Palette « Pop » OnBuch+ — mêmes hex que la classe P (plus_ui.dart)
% ============================================================
\definecolor{poporange}{HTML}{F59321}
\definecolor{popdark}{HTML}{E07A0C}
\definecolor{popcream}{HTML}{FFF6E8}
\definecolor{popink}{HTML}{1C1714}
\definecolor{poppurple}{HTML}{7A5AE0}
\definecolor{popgreen}{HTML}{2E9E6A}
\definecolor{poppink}{HTML}{E0457B}
\definecolor{popblue}{HTML}{2F7DE1}
\definecolor{popgold}{HTML}{C98A12}
\definecolor{popmuted}{HTML}{8A7F76}
\definecolor{popline}{HTML}{EADFD2}
\definecolor{popgray}{HTML}{8A7F76}
\colorlet{popgrey}{popgray}
% Alias tolérants (le modèle nomme parfois les couleurs différemment)
\colorlet{popwhite}{white}
\colorlet{popblack}{popink}
\colorlet{popred}{poppink}
\colorlet{popyellow}{popgold}
% Teintes claires (fonds de tableaux/encadrés) — le modèle nomme parfois
% les couleurs avec un suffixe L (ex. popblueL) sans les avoir définies.
\colorlet{poporangeL}{poporange!20}
\colorlet{popdarkL}{popdark!20}
\colorlet{poppurpleL}{poppurple!20}
\colorlet{popgreenL}{popgreen!20}
\colorlet{poppinkL}{poppink!20}
\colorlet{popblueL}{popblue!20}
\colorlet{popgoldL}{popgold!20}
\colorlet{popredL}{popred!20}
\colorlet{popgrayL}{popgray!20}
% Teintes claires pour fonds de tableaux / zones
\colorlet{popblueL}{popblue!10!white}
\colorlet{poporangeL}{poporange!14!white}
\colorlet{popgreenL}{popgreen!12!white}
\colorlet{poppurpleL}{poppurple!12!white}
\colorlet{poppinkL}{poppink!12!white}
\colorlet{popgoldL}{popgold!14!white}

\pagecolor{popcream}

% ============================================================
% Typographie
% ============================================================
\defaultfontfeatures{Ligatures=TeX}
\setmainfont{PlusJakartaSans-Medium.ttf}[Path=../../../fonts/,BoldFont=PlusJakartaSans-Bold.ttf]
% Symboles, API et emojis absents de la police principale : repli sur DejaVu Sans (caractères actifs)
\newfontface\symfont{DejaVuSans.ttf}[Path=../../../fonts/]
\makeatletter
\newcommand\symactive[1]{\catcode#1=13 \begingroup\lccode`\~=#1 \lowercase{\endgroup\def~}{{\symfont\char#1}}}
\newcommand\symblank[1]{\catcode#1=13 \begingroup\lccode`\~=#1 \lowercase{\endgroup\def~}{}}
\newcommand\symhyph[1]{\catcode#1=13 \begingroup\lccode`\~=#1 \lowercase{\endgroup\def~}{\mbox{-}}}
\newcount\symc
\newcommand\symrange[2]{\symc=#1 \@whilenum\symc<#2 \do{\expandafter\symactive\expandafter{\the\symc}\advance\symc by 1 }}
\newcommand\symblankrange[2]{\symc=#1 \@whilenum\symc<#2 \do{\expandafter\symblank\expandafter{\the\symc}\advance\symc by 1 }}
\makeatother
\symrange{"0250}{"02FF}\symactive{"2713}\symactive{"2714}\symactive{"2717}\symactive{"2718}\symactive{"2610}\symactive{"2611}\symactive{"2612}
\symrange{"2600}{"27BF}
\symhyph{"2011}\symhyph{"2E1A}\symblankrange{"1F300}{"1FAFF}\symblank{"FE0F}
% Maths en sans-serif (Fira Math) avec chiffres et lettres droites de Plus
% Jakarta, pour que les formules se fondent dans le texte.
\usepackage[math-style=ISO,bold-style=ISO]{unicode-math}
\setmathfont{FiraMath-Regular.otf}[Path=../../../fonts/]
\setmathfont{PlusJakartaSans-Medium.ttf}[Path=../../../fonts/,range={up/num,up/{Latin,latin},\mathup}]
\usepackage{icomma} % virgule décimale sans espace en mode math
% Caractères mathématiques Unicode absents des polices de texte (Plus Jakarta,
% Archivo Black) : rendus par leur équivalent mathématique, en texte comme en
% formule — sinon ils s'affichent en carré vide (ex. « Arithmétique dans ℤ »).
\usepackage{newunicodechar}
\newunicodechar{ℝ}{\ensuremath{\mathbb{R}}}
\newunicodechar{ℤ}{\ensuremath{\mathbb{Z}}}
\newunicodechar{ℂ}{\ensuremath{\mathbb{C}}}
\newunicodechar{ℕ}{\ensuremath{\mathbb{N}}}
\newunicodechar{ℚ}{\ensuremath{\mathbb{Q}}}
\newunicodechar{𝔻}{\ensuremath{\mathbb{D}}}
\newunicodechar{α}{\ensuremath{\alpha}}
\newunicodechar{β}{\ensuremath{\beta}}
\newunicodechar{γ}{\ensuremath{\gamma}}
\newunicodechar{δ}{\ensuremath{\delta}}
\newunicodechar{ε}{\ensuremath{\varepsilon}}
\newunicodechar{θ}{\ensuremath{\theta}}
\newunicodechar{λ}{\ensuremath{\lambda}}
\newunicodechar{μ}{\ensuremath{\mu}}
\newunicodechar{σ}{\ensuremath{\sigma}}
\newunicodechar{τ}{\ensuremath{\tau}}
\newunicodechar{φ}{\ensuremath{\varphi}}
\newunicodechar{ψ}{\ensuremath{\psi}}
\newunicodechar{ω}{\ensuremath{\omega}}
\newunicodechar{Σ}{\ensuremath{\Sigma}}
% majuscules produites par \MakeUppercase dans les titres (α-aminés -> Α)
\newunicodechar{Α}{\ensuremath{\alpha}}
\newunicodechar{Β}{\ensuremath{\beta}}
\newunicodechar{Γ}{\ensuremath{\Gamma}}
\newunicodechar{Δ}{\ensuremath{\Delta}}
\newunicodechar{Ε}{\ensuremath{\varepsilon}}
\newunicodechar{Θ}{\ensuremath{\Theta}}
\newunicodechar{Λ}{\ensuremath{\Lambda}}
\newunicodechar{Μ}{\ensuremath{\mu}}
\newunicodechar{Τ}{\ensuremath{\tau}}
\newunicodechar{Φ}{\ensuremath{\Phi}}
\newunicodechar{Ψ}{\ensuremath{\Psi}}
\newunicodechar{Ω}{\ensuremath{\Omega}}
\newunicodechar{↦}{\ensuremath{\mapsto}}
\newunicodechar{∈}{\ensuremath{\in}}
\newunicodechar{∉}{\ensuremath{\notin}}
\newunicodechar{∧}{\ensuremath{\wedge}}
\newunicodechar{⇔}{\ensuremath{\Leftrightarrow}}
\newunicodechar{⟺}{\ensuremath{\Longleftrightarrow}}
\newunicodechar{⇒}{\ensuremath{\Rightarrow}}
\newunicodechar{⟹}{\ensuremath{\Longrightarrow}}
\newunicodechar{∘}{\ensuremath{\circ}}
\newunicodechar{∪}{\ensuremath{\cup}}
\newunicodechar{∩}{\ensuremath{\cap}}
\newunicodechar{≡}{\ensuremath{\equiv}}
\newunicodechar{⊂}{\ensuremath{\subset}}
\newunicodechar{⊥}{\ensuremath{\perp}}
\newunicodechar{∃}{\ensuremath{\exists}}
\newunicodechar{∀}{\ensuremath{\forall}}
\newunicodechar{∛}{\ensuremath{\sqrt[3]{\,}}}
\newunicodechar{∜}{\ensuremath{\sqrt[4]{\,}}}
\newunicodechar{⃗}{\ensuremath{{}^{\to}}}
\newunicodechar{ⁿ}{\ifmmode^{n}\else\textsuperscript{\ensuremath{n}}\fi}
\newunicodechar{ˣ}{\ifmmode^{x}\else\textsuperscript{\ensuremath{x}}\fi}
\newunicodechar{ᵇ}{\ifmmode^{b}\else\textsuperscript{\ensuremath{b}}\fi}
\newunicodechar{ʳ}{\ifmmode^{r}\else\textsuperscript{\ensuremath{r}}\fi}
\newunicodechar{ᵘ}{\ifmmode^{u}\else\textsuperscript{\ensuremath{u}}\fi}
\newunicodechar{ʸ}{\ifmmode^{y}\else\textsuperscript{\ensuremath{y}}\fi}
\newunicodechar{ᵃ}{\ifmmode^{a}\else\textsuperscript{\ensuremath{a}}\fi}
\newunicodechar{ᵉ}{\ifmmode^{e}\else\textsuperscript{\ensuremath{e}}\fi}
\newunicodechar{ᵏ}{\ifmmode^{k}\else\textsuperscript{\ensuremath{k}}\fi}
\newunicodechar{ᵖ}{\ifmmode^{p}\else\textsuperscript{\ensuremath{p}}\fi}
\newunicodechar{ᴺ}{\ifmmode^{N}\else\textsuperscript{\ensuremath{N}}\fi}
\newunicodechar{ᶜ}{\ifmmode^{c}\else\textsuperscript{\ensuremath{c}}\fi}
\newunicodechar{ʰ}{\ifmmode^{h}\else\textsuperscript{\ensuremath{h}}\fi}
\newunicodechar{ᵠ}{\ifmmode^{q}\else\textsuperscript{\ensuremath{q}}\fi}
\newunicodechar{ₙ}{\ifmmode_{n}\else\textsubscript{\ensuremath{n}}\fi}
\newunicodechar{ₐ}{\ifmmode_{a}\else\textsubscript{\ensuremath{a}}\fi}
\newunicodechar{ᵢ}{\ifmmode_{i}\else\textsubscript{\ensuremath{i}}\fi}
\newunicodechar{ᵣ}{\ifmmode_{r}\else\textsubscript{\ensuremath{r}}\fi}
\newunicodechar{ₖ}{\ifmmode_{k}\else\textsubscript{\ensuremath{k}}\fi}
\newunicodechar{ᵦ}{\ifmmode_{\beta}\else\textsubscript{\ensuremath{\beta}}\fi}
\newunicodechar{ₓ}{\ifmmode_{x}\else\textsubscript{\ensuremath{x}}\fi}
\newfontfamily\popdisplay{ArchivoBlack-Regular.ttf}[Path=../../../fonts/]
\newfontfamily\poplogo{SpaceGrotesk-Bold.ttf}[Path=../../../fonts/]
\newfontfamily\popsemi{PlusJakartaSans-SemiBold.ttf}[Path=../../../fonts/]
\newfontfamily\popxbold{PlusJakartaSans-ExtraBold.ttf}[Path=../../../fonts/]
\newfontfamily\popxboldls{PlusJakartaSans-ExtraBold.ttf}[Path=../../../fonts/,LetterSpace=3.0]

\setlist[enumerate]{leftmargin=1.7em,itemsep=5pt,topsep=4pt}
\setlist[itemize]{leftmargin=1.7em,itemsep=4pt,topsep=4pt,label={\color{poporange}\textbullet}}
\setlength{\parindent}{0pt}
\setlength{\parskip}{5pt}
\renewcommand{\arraystretch}{1.25}

% pgfplots : style Pop par défaut
\pgfplotsset{
  every axis/.append style={axis line style={popink,line width=1pt},tick style={popink},
    grid style={popline,line width=0.5pt},label style={font=\small},tick label style={font=\footnotesize}},
  popcurve/.style={poporange,line width=1.8pt,no markers,smooth},
}
% Même style disponible dans un \draw TikZ simple (hors pgfplots)
\tikzset{popcurve/.style={poporange,line width=1.8pt}}
\tikzset{popgrid/.style={popline,very thin}}
\colorlet{popcurve}{poporange} % le nom du style est parfois utilisé comme couleur

% ============================================================
% Pill / squiggle
% ============================================================
\newcommand{\poppill}[3]{%
  \tikz[baseline=-0.55ex]\node[draw=popink,line width=1.2pt,rounded corners=2.6pt,%
    fill=#2,inner xsep=6.5pt,inner ysep=3pt]{%
    \popxboldls\fontsize{7}{7}\selectfont\color{#3}\MakeUppercase{#1}};%
}
\newcommand{\popsquiggle}[1]{%
  \tikz[baseline=-0.4ex]{
    \draw[#1,line width=2.6pt,line cap=round]
      plot [smooth] coordinates {(0,0.09) (0.34,0.01) (0.68,0.21) (1.02,0.01) (1.36,0.10)};
  }%
}
% Mot-clé surligné (marqueur orange sous le texte)
\newcommand{\cle}[1]{{\popxbold\color{popdark}#1}}

% ============================================================
% En-tête / pied
% ============================================================
\pagestyle{fancy}
\fancyhf{}
\renewcommand{\headrulewidth}{0pt}
\renewcommand{\footrulewidth}{0pt}
\lhead{\raisebox{-0.15ex}{\poplogo\fontsize{13}{13}\selectfont\color{popink}OnBuch\color{poporange}+}}
\rhead{\poppill{\DOCMATIERE\ \textbullet\ \DOCNIVEAU\ \textbullet\ Leçon \DOCLECON}{white}{popink}}
\lfoot{\popxbold\fontsize{7.6}{7.6}\selectfont\color{popmuted}\MakeUppercase{Cours signé OnBuch+}\ \textbullet\ {\popsemi\DOCTITRE}}
\rfoot{\tikz[baseline=-0.6ex]\node[rounded corners=8.5pt,fill=popink,minimum height=17pt,minimum width=17pt,inner xsep=5pt,inner sep=0]{\popxbold\fontsize{8}{8}\selectfont\color{popcream}\thepage\,/\,\pageref*{LastPage}};}

% ============================================================
% Titres
% ============================================================
\newcommand{\chaptitre}[2]{%
  \begin{tcolorbox}[enhanced,colback=poporange,colframe=popink,boxrule=2.6pt,arc=11pt,
      drop shadow=popink,left=15pt,right=15pt,top=12pt,bottom=11pt,before skip=3pt,after skip=18pt]
    \poppill{#2}{popink}{popcream}\par\vspace{9pt}
    {\raggedright\hyphenpenalty=10000\exhyphenpenalty=10000\popdisplay\fontsize{22}{24}\selectfont\color{popcream}\MakeUppercase{#1}\par}
  \end{tcolorbox}
}
% Section numérotée : pastille orange avec le numéro + titre Archivo
\newcounter{popsec}
\newcommand{\coursec}[1]{%
  \stepcounter{popsec}\setcounter{popsub}{0}%
  \par\vspace{16pt}\noindent
  \begin{minipage}{\linewidth}
  \tikz[baseline=-0.7ex]\node[draw=popink,line width=1.6pt,fill=poporange,rounded corners=4pt,minimum size=22pt,inner sep=0]{\popdisplay\fontsize{12}{12}\selectfont\color{popcream}\thepopsec};%
  \hspace{8pt}{\popdisplay\fontsize{15}{17}\selectfont\color{popink}\MakeUppercase{#1}}\par
  \vspace{4pt}\noindent{\color{poporange}\rule{36pt}{3.6pt}}
  \end{minipage}\par\nopagebreak\vspace{8pt}
}
\newcounter{popsub}
\newcommand{\courssub}[1]{%
  \stepcounter{popsub}%
  \par\vspace{9pt}\noindent{\popxbold\fontsize{11.5}{13}\selectfont\color{popink}{\color{poporange}\thepopsec.\thepopsub}\hspace{6pt}#1}\par\nopagebreak\vspace{4pt}
}

% ============================================================
% Boîtes pédagogiques — même recette : fond blanc, bordure épaisse colorée,
% ombre dure encre, badge plein. Titre optionnel [#1].
% ============================================================
\tcbset{popbase/.style={breakable,enhanced,colback=white,boxrule=2.3pt,arc=9pt,
  shadow={3pt}{-3pt}{0pt}{popink},left=14pt,right=14pt,top=11pt,bottom=11pt,before skip=11pt,after skip=15pt,
  fontupper=\popsemi\fontsize{10.6}{14.5}\selectfont\color{popink},
  fontlower=\popsemi\fontsize{10.6}{14.5}\selectfont\color{popink}}}
% Coche dessinée (le glyphe ✓ n'existe pas dans les polices du document)
\newcommand{\popcheck}[1]{\tikz[baseline=-0.5ex]\draw[#1,line width=1.6pt,line cap=round,line join=round](-0.13,0.0)--(-0.04,-0.09)--(0.13,0.1);}
\providecommand{\coche}{\popcheck{popgreen}}
\providecommand{\resultat}[1]{\par\smallskip\noindent\fcolorbox{popgreen}{popgreenL}{\parbox{\dimexpr\linewidth-2\fboxsep-2\fboxrule\relax}{\textbf{Résultat.} #1}}\par\smallskip}
\newcommand{\popboxhead}[3]{\poppill{#1}{#2}{white}\ifx\relax#3\relax\else\hspace{6pt}{\popxbold\fontsize{10.6}{12}\selectfont\color{popink}#3}\fi\par\vspace{7pt}}

\newtcolorbox{aretenir}[1][]{popbase,colframe=popgreen,before upper={\popboxhead{À retenir}{popgreen}{#1}}}
% Texte en allemand (document de lecture, dialogue, extrait) : cadre
% d'étude. Un titre optionnel (source, auteur) s'affiche dans l'en-tête.
\newtcolorbox{textebox}[1][]{popbase,colframe=popblue,before upper={\popboxhead{Text}{popblue}{#1}\begin{otherlanguage}{ngerman}},after upper={\end{otherlanguage}}}
% Marques ✓ / ✗ (forme correcte / fautive) souvent utilisées par les modèles
\providecommand{\cross}{\textcolor{poppink}{\ensuremath{\times}}}
\providecommand{\xmark}{\cross}\providecommand{\crossmark}{\cross}\providecommand{\cmark}{\ensuremath{\checkmark}}
% Ligne de réponse / justification à compléter (inventée par les modèles)
\providecommand{\justification}[1]{\par\noindent\textit{Justification~:} \dotfill\par}
% \textipa (package tipa non chargé) : la police ne couvre pas l'API -> simple passage
\providecommand{\textipa}[1]{#1}
% Soulignement (ulem non chargé) : \uline, \ul
\providecommand{\uline}[1]{\underline{#1}}\providecommand{\ul}[1]{\underline{#1}}
% \glossaire{\item ...} inventée par les modèles : liste de glossaire
\providecommand{\glossaire}[1]{\begin{itemize}#1\end{itemize}}
% \alert (beamer) inventée par les modèles : texte mis en évidence
\providecommand{\alert}[1]{\textbf{\textcolor{poppink}{#1}}}
% variantes ASCII d'ordinaux écrites en macro
\providecommand{\iere}{\textsuperscript{re}}\providecommand{\ieme}{\textsuperscript{e}}\providecommand{\ere}{\textsuperscript{re}}\providecommand{\eme}{\textsuperscript{e}}\providecommand{\er}{\textsuperscript{er}}\providecommand{\ie}{\textsuperscript{e}}
% Ordinaux français écrits en macro par les modèles : 1\ière, 3\ième
\newcommand{\ière}{\textsuperscript{re}}\newcommand{\ième}{\textsuperscript{e}}
% \exemple (inventée par les modèles) : étiquette « Exemple : »
\providecommand{\exemple}{\textbf{Exemple~:}\ }
% \square utilisable aussi hors mode math (cases à cocher)
\AtBeginDocument{\let\squareMath\square\renewcommand{\square}{\ensuremath{\squareMath}}}
% Mot ou phrase en allemand à l'intérieur d'un paragraphe français
\newcommand{\all}[1]{\ifmmode\text{\textit{#1}}\else\begingroup\selectlanguage{ngerman}\itshape #1\endgroup\fi}
\let\esp\all % alias toléré
\newtcolorbox{exemplebox}[1][]{popbase,colframe=poporange,before upper={\popboxhead{Exemple}{poporange}{#1}}}
\newtcolorbox{definition}[1][]{popbase,colframe=popblue,before upper={\popboxhead{Définition}{popblue}{#1}}}
\newtcolorbox{propriete}[1][]{popbase,colframe=poppurple,before upper={\popboxhead{Propriété}{poppurple}{#1}}}
\newtcolorbox{methode}[1][]{popbase,colframe=popgold,before upper={\popboxhead{Méthode}{popgold}{#1}}}
\newtcolorbox{attention}[1][]{popbase,colframe=poppink,before upper={\popboxhead{Attention}{poppink}{#1}}}
\newtcolorbox{experience}[1][]{popbase,colframe=popblue,colback=popblueL,before upper={\popboxhead{Expérience}{popblue}{#1}}}
% « Le savais-tu ? » : carte violette pleine (contexte, culture, Cameroun)
\newtcolorbox{savaistu}[1][]{popbase,colback=poppurple,colframe=popink,
  fontupper=\popsemi\fontsize{10.4}{14.2}\selectfont\color{white},
  before upper={\poppill{Le savais-tu\,?}{white}{poppurple}\ifx\relax#1\relax\else\hspace{6pt}{\popxbold\color{white}#1}\fi\par\vspace{7pt}}}
% Exercice résolu : énoncé en haut, \tcblower puis la solution sur fond crème
\newcounter{popexo}\newcounter{popexr}
\newtcolorbox{exoresolu}[1][]{popbase,colframe=poporange,colbacklower=popcream,
  segmentation style={popink,line width=1.2pt,dashed},
  before upper={\stepcounter{popexr}\popboxhead{Exercice résolu \thepopexr}{poporange}{#1}},
  before lower={\poppill{Solution}{popink}{popcream}\par\vspace{6pt}}}
% Exercice (non résolu en place) — la correction est dans la partie Corrigés
\newtcolorbox{exercice}[1][]{popbase,colframe=popink,
  before upper={\stepcounter{popexo}\popboxhead{Exercice \thepopexo}{popink}{#1}}}
% Corrigé
\newtcolorbox{corrige}[1][]{popbase,colframe=popgreen,colback=popgreenL,
  before upper={\popboxhead{Corrigé}{popgreen}{#1}}}
% Objectifs / prérequis / bilan
\newtcolorbox{objectifs}{popbase,colframe=popink,colback=poporangeL,
  before upper={\popboxhead{Objectifs de la leçon}{popink}{}}}
\newtcolorbox{prerequis}{popbase,colframe=popink,colback=popblueL,
  before upper={\popboxhead{Prérequis}{popblue}{}}}
\newtcolorbox{bilan}{popbase,colframe=popink,colback=white,boxrule=2.8pt,
  before upper={\popboxhead{Fiche bilan}{poporange}{L'essentiel à connaître pour l'examen}}}
% Figure encadrée avec légende
\newtcolorbox{popfigure}[1][]{enhanced,breakable=false,colback=white,colframe=popline,boxrule=1.4pt,arc=8pt,
  left=10pt,right=10pt,top=10pt,bottom=8pt,before skip=10pt,after skip=12pt,halign=center,
  lower separated=false,halign lower=center,
  fontlower=\popsemi\fontsize{9}{11.5}\selectfont\color{popmuted},
  #1}
\newcommand{\legende}[1]{\par\vspace{6pt}{\popsemi\fontsize{9}{11.5}\selectfont\color{popmuted}#1}}

% ============================================================
% Signature OnBuch+
% ============================================================
\newcommand{\poplogoinline}[1]{{\poplogo\fontsize{#1}{#1}\selectfont\color{popink}OnBuch\color{poporange}+}}
\newcommand{\signatureonbuch}{%
  \begin{tcolorbox}[enhanced,colback=popink,colframe=popink,boxrule=2.2pt,arc=10pt,
      drop shadow=poporange,left=14pt,right=14pt,top=10pt,bottom=10pt,before skip=0pt,after skip=0pt]
    \tikz[baseline=-0.7ex]\node[circle,fill=popgreen,draw=popcream,line width=1.3pt,minimum size=22pt,inner sep=0]{\popcheck{white}};%
    \hspace{9pt}\begin{minipage}[c]{0.62\linewidth}
      {\popxbold\fontsize{12}{13}\selectfont\color{popcream}Signé\ \textbullet\ {\poplogo OnBuch\color{poporange}+}}\par\vspace{2pt}
      {\raggedright\popsemi\fontsize{8}{10.5}\selectfont\color{popline}\MakeUppercase{Équipe pédagogique OnBuch+}\\\MakeUppercase{Conforme au programme officiel MINESEC}\par}
    \end{minipage}\hfill
    {\poplogo\fontsize{22}{22}\selectfont\color{popcream}OnBuch\color{poporange}+}
  \end{tcolorbox}%
}

% ============================================================
% Page de garde
% #1 titre · #2 matière · #3 niveau · #4 module · #5 n° leçon · #6 durée ·
% #7 description / objectifs (texte ou itemize)
% ============================================================
\newcommand{\pagegarde}[7]{%
  \thispagestyle{empty}
  \newgeometry{a4paper,margin=1.7cm,top=1.6cm,bottom=1.4cm}%
  \begin{tikzpicture}[remember picture,overlay]
    \fill[poporange!18] ([xshift=-3.2cm,yshift=-3.6cm]current page.north east) circle (4.6cm);
    \fill[poppurple!14] ([xshift=2.4cm,yshift=7.6cm]current page.south west) circle (3.8cm);
  \end{tikzpicture}%
  {\poplogo\fontsize{30}{30}\selectfont\color{popink}OnBuch\color{poporange}+}\hfill
  \raisebox{0.6ex}{\poppill{Cours signé \textbullet\ Premium}{popink}{popcream}}\par
  \vspace{1pt}\popsquiggle{poppurple}\par
  \vspace{8pt}
  \begin{tcolorbox}[enhanced,colback=poporange,colframe=popink,boxrule=2.8pt,arc=13pt,
      drop shadow=popink,left=18pt,right=18pt,top=12pt,bottom=13pt,after skip=10pt]
    \poppill{#2 \textbullet\ #3}{popink}{popcream}\hspace{5pt}\poppill{#4}{white}{popink}\par\vspace{8pt}
    {\popdisplay\fontsize{14}{14}\selectfont\color{popink}LEÇON #5}\par\vspace{4pt}
    {\raggedright\hyphenpenalty=10000\exhyphenpenalty=10000\popdisplay\fontsize{25}{27}\selectfont\color{popcream}\MakeUppercase{#1}\par}
    \vspace{8pt}
    {\popxbold\fontsize{9.5}{9.5}\selectfont\color{popink}\MakeUppercase{Durée conseillée : #6}}
  \end{tcolorbox}
  \begin{tcolorbox}[enhanced,colback=white,colframe=popink,boxrule=2.2pt,arc=10pt,
      drop shadow=popink,left=16pt,right=16pt,top=10pt,bottom=10pt,after skip=8pt]
    \poppill{Dans cette leçon}{poppurple}{white}\par\vspace{5pt}
    \popsemi\fontsize{9.3}{12.6}\selectfont\color{popink}#7
  \end{tcolorbox}
  \noindent\begin{minipage}[t]{0.487\linewidth}
    \begin{tcolorbox}[enhanced,colback=popgreen,colframe=popink,boxrule=2.2pt,arc=10pt,
        drop shadow=popink,left=11pt,right=11pt,top=10pt,bottom=10pt,height=3.1cm,valign=center]
      \tcbox[colback=white,colframe=popink,boxrule=1.3pt,arc=5pt,boxsep=0pt,nobeforeafter,
        left=3pt,right=3pt,top=3pt,bottom=3pt,valign=center]{\qrcode[height=1.9cm]{https://play.google.com/store/apps/details?id=com.luvvix.onbuch&hl=fr}}%
      \hspace{8pt}\begin{minipage}[c]{0.5\linewidth}
        {\popdisplay\fontsize{11}{12.5}\selectfont\color{white}\MakeUppercase{Télécharge OnBuch}}\par\vspace{4pt}
        {\raggedright\popsemi\fontsize{8.4}{11}\selectfont\color{white}Gratuit \textbullet\ Google Play\\Tuteur IA, annales, concours, résultats\par}
      \end{minipage}
    \end{tcolorbox}
  \end{minipage}\hfill
  \begin{minipage}[t]{0.487\linewidth}
    \begin{tcolorbox}[enhanced,colback=poppurple,colframe=popink,boxrule=2.2pt,arc=10pt,
        drop shadow=popink,left=11pt,right=11pt,top=10pt,bottom=10pt,height=3.1cm,valign=center]
      \tikz[baseline=-0.7ex]\node[circle,fill=white,draw=popink,line width=1.3pt,minimum size=1.6cm,inner sep=0]{\popdisplay\fontsize{24}{24}\selectfont\color{poppurple}+};%
      \hspace{8pt}\begin{minipage}[c]{0.6\linewidth}
        {\popdisplay\fontsize{11}{12.5}\selectfont\color{white}\MakeUppercase{Passe à OnBuch+}}\par\vspace{4pt}
        {\raggedright\popsemi\fontsize{8.4}{11}\selectfont\color{white}Tous les cours signés, Tuteur IA illimité, fiches PDF — onglet OnBuch+ de l'app\par}
      \end{minipage}
    \end{tcolorbox}
  \end{minipage}
  \vspace{8pt}
  \popcontenu
  \vfill
  \signatureonbuch
  \restoregeometry
  \newpage
}

% Rangée « Ce cours contient » (page de garde)
\newcommand{\poptile}[3]{% #1 couleur #2 gros texte #3 légende
  \begin{tcolorbox}[enhanced,colback=#1,colframe=popink,boxrule=2pt,arc=9pt,shadow={2.5pt}{-2.5pt}{0pt}{popink},
      left=6pt,right=6pt,top=9pt,bottom=9pt,halign=center,valign=center,height=1.75cm,width=\linewidth,nobeforeafter]
    {\popdisplay\fontsize{12.5}{13}\selectfont\color{white}\MakeUppercase{#2}}\par\vspace{3pt}
    {\centering\popsemi\fontsize{7.8}{9.5}\selectfont\color{white}#3\par}
  \end{tcolorbox}}
\newcommand{\popcontenu}{%
  \par\noindent{\popxboldls\fontsize{7.5}{7.5}\selectfont\color{popmuted}CE COURS SIGNÉ CONTIENT}\par\vspace{7pt}
  \noindent\begin{minipage}[t]{0.235\linewidth}\poptile{popblue}{Cours}{complet, illustré, pas à pas}\end{minipage}\hfill
  \begin{minipage}[t]{0.235\linewidth}\poptile{poporange}{Méthodes}{et exercices résolus}\end{minipage}\hfill
  \begin{minipage}[t]{0.235\linewidth}\poptile{poppink}{Exercices}{type examen + corrigés}\end{minipage}\hfill
  \begin{minipage}[t]{0.235\linewidth}\poptile{popgreen}{Fiche bilan}{l'essentiel à retenir}\end{minipage}\par}

% Sommaire
\newtcolorbox{sommaire}{enhanced,colback=white,colframe=popink,boxrule=2.2pt,arc=10pt,
  drop shadow=popink,left=18pt,right=18pt,top=13pt,bottom=13pt,before skip=6pt,after skip=20pt,
  before upper={\poppill{Sommaire}{poppurple}{white}\par\vspace{9pt}\popsemi\fontsize{10.6}{14.5}\selectfont\color{popink}}}

% Signature de fin de cours — poussée tout en bas de la dernière page
\newcommand{\findecours}{%
  \par\vfill\nopagebreak
  \begin{center}\popsquiggle{poporange}\end{center}\vspace{4pt}
  \signatureonbuch
}
\AtBeginDocument{\def\llbracket{\mathopen{[\mkern-3.5mu[}}\def\rrbracket{\mathclose{]\mkern-3.5mu]}}} % intervalles d'entiers (glyphe ⟦ absent de la police)
% Glyphes absents de Fira Math : redéfinis à partir de glyphes disponibles
\AtBeginDocument{%
  \def\setminus{\mathbin{\backslash}}\let\smallsetminus\setminus
  \def\vdots{\mathord{\vbox{\baselineskip3.2pt\lineskiplimit0pt\kern4pt\hbox{.}\hbox{.}\hbox{.}}}}%
  \def\ddots{\mathinner{\mkern1mu\raise6.5pt\hbox{.}\mkern2mu\raise3.6pt\hbox{.}\mkern2mu\raise0.7pt\hbox{.}\mkern1mu}}%
  \def\triangle{\mathord{\scalebox{0.9}{$\symup{\Delta}$}}}%
  \def\nabla{\mathord{\raisebox{\depth}{\scalebox{1}[-1]{$\symup{\Delta}$}}}}%
  \def\checkmark{\popcheck{popdark}}%
  \def\star{\mathbin{\ast}}\def\bigstar{\mathord{\ast}}%
  \def\diamond{\mathbin{\scalebox{0.6}{$\Diamond$}}}%
  \def\bigcup{\mathop{\mathchoice{\raisebox{-0.3em}{\scalebox{1.7}{$\cup$}}}{\scalebox{1.25}{$\cup$}}{\cup}{\cup}}\displaylimits}%
  \def\bigcap{\mathop{\mathchoice{\raisebox{-0.3em}{\scalebox{1.7}{$\cap$}}}{\scalebox{1.25}{$\cap$}}{\cap}{\cap}}\displaylimits}%
  \def\lesseqqgtr{\mathrel{\lessgtr}}\def\bigoplus{\mathop{\oplus}\displaylimits}%
}
\providecommand{\ssi}{si et seulement si} % abréviation fréquente
\AtBeginDocument{\providecommand{\wideparen}{\overparen}} % arc de cercle

%% FIGFIX-BEGIN v5 — correctifs globaux des figures (docs/onbuch-generation/tools/figfix)
\usepackage{adjustbox}\usepackage{etoolbox}\tcbuselibrary{raster}
% toute figure TikZ/pgfplots plus large que la ligne est réduite à la largeur disponible
\BeforeBeginEnvironment{tikzpicture}{\begin{adjustbox}{max width=\linewidth}}
\AfterEndEnvironment{tikzpicture}{\end{adjustbox}}
% légendes pgfplots lisibles par-dessus les courbes
\pgfplotsset{every axis/.append style={legend style={fill=white,fill opacity=0.92,text opacity=1,draw=black!15,inner sep=3pt}}}
% étiquettes d'arêtes (syntaxe edge["..."]) avec fond blanc
\tikzset{every edge quotes/.append style={fill=white,inner sep=1.2pt}}
%% FIGFIX-END
\begin{document}

\pagegarde{La vie culturelle 1 (les arts et la littérature)}{Allemand}{Tle A}{Chapitre 5 : Citoyenneté}{AL22}{2 h}{Cette leçon de type texte fait découvrir la vie culturelle germanophone à travers les arts et la littérature. Elle combine la compréhension d'un texte authentique, l'enrichissement du lexique thématique et l'étude de trois points de grammaire fondamentaux : les participes présent et passé en épithète, le génitif et le prétérit des verbes forts.
\vspace{4pt}
\begin{multicols}{2}\small
\begin{itemize}
\item À la fin de cette leçon, je sais comprendre globalement et en détail un texte allemand sur les arts et la littérature.
\item À la fin de cette leçon, je sais employer correctement le vocabulaire des arts et de la littérature avec articles et pluriels.
\item À la fin de cette leçon, je sais former et accorder les participes présent et passé employés comme épithètes.
\item À la fin de cette leçon, je sais former et utiliser le génitif pour exprimer la possession et la complémentation.
\item À la fin de cette leçon, je sais conjuguer au prétérit les verbes forts les plus fréquents.
\item À la fin de cette leçon, je sais répondre par écrit et oralement à des questions de compréhension.
\end{itemize}\end{multicols}}

\begin{sommaire}
\begin{multicols}{2}
\begin{enumerate}
\item Texte support : Kunst und Literatur im deutschsprachigen Raum
\item Lexique thématique : die Kunst und die Literatur
\item Grammaire 1 : les participes présent et passé en épithète
\item Grammaire 2 : le génitif
\item Conjugaison : le prétérit des verbes forts
\item Méthode du commentaire, commentaire modèle et production guidée
\item Activité d'intégration
\item Méthodes et exercices résolus
\item Exercices
\item Corrigés des exercices
\item Fiche bilan
\end{enumerate}
\end{multicols}
\end{sommaire}

\begin{objectifs}
\begin{itemize}
\item À la fin de cette leçon, je sais comprendre globalement et en détail un texte allemand sur les arts et la littérature.
\item À la fin de cette leçon, je sais employer correctement le vocabulaire des arts et de la littérature avec articles et pluriels.
\item À la fin de cette leçon, je sais former et accorder les participes présent et passé employés comme épithètes.
\item À la fin de cette leçon, je sais former et utiliser le génitif pour exprimer la possession et la complémentation.
\item À la fin de cette leçon, je sais conjuguer au prétérit les verbes forts les plus fréquents.
\item À la fin de cette leçon, je sais répondre par écrit et oralement à des questions de compréhension.
\item À la fin de cette leçon, je sais rédiger un résumé et un court paragraphe sur une œuvre ou un artiste.
\item À la fin de cette leçon, je sais commenter un texte littéraire ou artistique selon une méthode claire.
\end{itemize}
\end{objectifs}

\begin{prerequis}
\begin{itemize}
\item Les articles définis et indéfinis et la déclinaison de l'adjectif (Première)
\item Le prétérit des verbes faibles et des auxiliaires haben/sein (Première)
\item Le nominatif, l'accusatif et le datif (Seconde/Première)
\item La formation du participe passé (ge-...-t / ge-...-en) vue avec le parfait
\item Le vocabulaire de base des loisirs et des sorties culturelles
\end{itemize}
\end{prerequis}

\newpage

\coursec{Texte support : Kunst und Literatur im deutschsprachigen Raum}

\courssub{Situation d'accroche et problématique}

Au lycée de Yaoundé, le club culturel prépare une grande exposition sur les artistes camerounais et allemands. En feuilletant un magazine allemand, les élèves tombent sur un article consacré à Goethe et à Caspar David Friedrich. Ils lisent : \all{ein berühmter Schriftsteller}, \all{ein beeindruckendes Gemälde}, \all{die Museen Berlins}... Immédiatement, les questions fusent : pourquoi ces formes en \all{-d} et \all{-end} placées devant les noms ? Comment dit-on « le roman de l'écrivain » ? Et pourquoi le verbe \all{schrieb} ne ressemble-t-il pas à \all{schreiben} ?

\textbf{Problématique :} comment comprendre et exploiter un texte sur les arts et la littérature dans l'espace germanophone, en maîtrisant le vocabulaire du thème, les participes épithètes, le génitif et le prétérit des verbes forts ? Cette leçon ouvre les portes du monde culturel de l'Allemagne, de l'Autriche et de la Suisse, de Goethe à Thomas Mann.

\courssub{Lecture globale : premiers repères}

Avant toute lecture détaillée, il faut se poser trois questions simples : \emph{De quoi parle le texte ? Qui sont les personnes mentionnées ? Où et quand se situe l'action ?} Une première lecture rapide du texte ci-dessous permet de répondre : le thème est la \cle{culture} (arts et littérature) ; les personnages sont Goethe, Caspar David Friedrich et Thomas Mann ; le cadre spatial est l'Allemagne (Berlin, Weimar), l'Autriche et la Suisse ; le cadre temporel va du XVIII\textsuperscript{e} siècle à aujourd'hui.

\begin{methode}[Lire un texte au Bac : les quatre étapes]
\begin{enumerate}
\item Lire le \cle{titre} et la \cle{source} : ils annoncent le thème et le type de texte.
\item Faire une première lecture rapide pour saisir l'\cle{idée générale}, sans s'arrêter sur les mots inconnus.
\item Faire une deuxième lecture avec le \cle{glossaire} : entourer les mots-clés, les noms propres, les dates.
\item Souligner les \cle{mots-clés} de chaque question avant de rédiger la réponse en phrase complète.
\end{enumerate}
\end{methode}

\begin{attention}[Le piège de la traduction mot à mot]
Ne traduis jamais un texte mot à mot : l'allemand et le français n'ont ni le même ordre des mots ni les mêmes constructions. Cherche d'abord le \cle{sens global} de la phrase : qui fait quoi, quand, où ? Exemple : \all{Die Gemälde des deutschen Malers zeigen die Natur.} se comprend globalement « Les tableaux du peintre allemand montrent la nature », et non en calquant mot à mot. Le sens d'abord, la précision ensuite !
\end{attention}

\courssub{Le texte : Kultur in Deutschland, Österreich und der Schweiz}

\begin{textebox}[Kultur in Deutschland, Österreich und der Schweiz]
Deutschland ist ein Land mit einer reichen Kultur. Johann Wolfgang von Goethe, der berühmteste Schriftsteller des Landes, lebte von 1749 bis 1832 und schrieb den Roman „Die Leiden des jungen Werthers“. Dieser Roman machte den jungen Dichter in ganz Europa bekannt. Später arbeitete Goethe viele Jahre an seinem Drama „Faust“, dem meistgelesenen Werk der deutschen Literatur.

Die Gemälde des deutschen Malers Caspar David Friedrich zeigen die beeindruckende Natur der Romantik : einsame Menschen vor weiten Landschaften, neblige Berge und stille Wälder. Viele Touristen besuchen heute die Museen Berlins, zum Beispiel die Alte Nationalgalerie, wo man die bekannten Werke des Künstlers sehen kann.

Auch Österreich und die Schweiz haben eine große Kultur. Wien, die Hauptstadt Österreichs, war im 18. Jahrhundert das Zentrum der klassischen Musik : Mozart und Beethoven lebten und komponierten dort. Der österreichische Schriftsteller Stefan Zweig schrieb spannende Romane und Erzählungen, die man noch heute in vielen Sprachen liest.

Die moderne Literatur ist ebenfalls lebendig. Thomas Mann, der Träger des Nobelpreises von 1929, beschrieb in „Buddenbrooks“ die Geschichte einer reichen Familie. Heute schreiben viele junge Autoren über das Leben in den großen Städten. Zahlreiche Ausstellungen und Lesungen finden jedes Jahr in Berlin, Wien und Zürich statt. Die Kultur der deutschsprachigen Länder bleibt also ein faszinierendes Erbe für die ganze Welt.
\end{textebox}

\textbf{Glossaire (mots nouveaux) :}

\begin{center}
\begin{tabularx}{\linewidth}{|X|X|}
\hline
\rowcolor{popblueL} \textbf{Deutsch} & \textbf{Français} \\
\hline
der Schriftsteller, - / die Schriftstellerin, -nen & l'écrivain / la femme écrivain \\
\hline
der Dichter, - & le poète \\
\hline
das Gedicht, -e & le poème \\
\hline
das Gemälde, - & le tableau, la peinture \\
\hline
der Maler, - & le peintre \\
\hline
die Ausstellung, -en & l'exposition \\
\hline
die Epoche, -n & l'époque \\
\hline
berühmt & célèbre \\
\hline
beeindruckend & impressionnant \\
\hline
die Erzählung, -en & la nouvelle, le récit \\
\hline
der Träger des Nobelpreises & le lauréat du prix Nobel \\
\hline
das Erbe & l'héritage \\
\hline
\end{tabularx}
\end{center}

\begin{savaistu}[Goethe et le « Faust »]
Johann Wolfgang von Goethe (1749-1832) est l'auteur le plus célèbre de la langue allemande, un peu comme Victor Hugo pour le français. Son drame \all{„Faust“}, l'histoire d'un savant qui vend son âme au diable, est considéré comme le chef-d'œuvre de la littérature allemande. Goethe était aussi scientifique, ministre et amateur d'art : un vrai esprit universel ! La ville de Weimar, où il vécut, est aujourd'hui un haut lieu culturel de l'Allemagne.
\end{savaistu}

\courssub{Lecture détaillée : activités de repérage}

La lecture détaillée consiste à observer la langue du texte de près. Trois tâches guident ton œil :

\begin{enumerate}
\item \textbf{Souligne les participes épithètes} (adjectifs verbaux en \all{-d} ou \all{-end} devant un nom) : \all{der berühmteste Schriftsteller}, \all{das meistgelesene Werk}, \all{die beeindruckende Natur}, \all{die bekannten Werke}, \all{spannende Romane}, \all{ein faszinierendes Erbe}.
\item \textbf{Entoure les génitifs} (compléments du nom en \all{-s} ou avec article \all{des/der}) : \all{des Landes}, \all{des deutschen Malers}, \all{der Romantik}, \all{der deutschen Literatur}, \all{der klassischen Musik}, \all{des Nobelpreises}, \all{der deutschsprachigen Länder}.
\item \textbf{Relève les verbes forts au prétérit} : \all{schrieb} (schreiben), \all{machte} (machen, faible), \all{war} (sein), \all{lebte} (leben, faible), \all{beschrieb} (beschreiben), \all{fanden ... statt} (stattfinden).
\end{enumerate}

\begin{center}
\begin{tabularx}{\linewidth}{|X|X|X|}
\hline
\rowcolor{popblueL} \textbf{Forme du texte} & \textbf{Infinitif / base} & \textbf{Français} \\
\hline
schrieb & schreiben & il écrivit \\
\hline
war & sein & était \\
\hline
beschrieb & beschreiben & il décrivit \\
\hline
fanden statt & stattfinden & eurent lieu \\
\hline
beeindruckende (Natur) & beeindrucken (part. présent) & impressionnante \\
\hline
meistgelesene (Werk) & lesen (part. passé) & le plus lu \\
\hline
\end{tabularx}
\end{center}

\courssub{Compréhension orale : écoute du texte}

L'enseignant lit le texte à voix haute, les livres fermés. Écoute d'abord sans prendre de notes, puis une deuxième fois en notant les noms propres et les dates. Réponds ensuite par \all{richtig} ou \all{falsch} :

\begin{exercice}[Hörverstehen : richtig oder falsch ?]
\begin{enumerate}
\item Goethe schrieb den Roman „Die Leiden des jungen Werthers“.
\item Caspar David Friedrich war ein berühmter Musiker.
\item Wien war das Zentrum der klassischen Musik.
\item Thomas Mann bekam den Nobelpreis 1929.
\item Die Museen Berlins besuchen nur wenige Touristen.
\end{enumerate}
\end{exercice}

\begin{corrige}[Hörverstehen]
\begin{enumerate}
\item \all{richtig} — le texte dit : \all{... schrieb den Roman „Die Leiden des jungen Werthers“}.
\item \all{falsch} — Friedrich war ein \all{Maler} (peintre), pas un musicien.
\item \all{richtig} — \all{Wien ... war das Zentrum der klassischen Musik}.
\item \all{richtig} — \all{Thomas Mann, der Träger des Nobelpreises von 1929}.
\item \all{falsch} — \all{Viele Touristen besuchen heute die Museen Berlins}.
\end{enumerate}
\end{corrige}

\courssub{Questions de compréhension écrite}

\begin{exoresolu}[Fragen zum Text]
Réponds aux six questions suivantes en phrases complètes, en allemand, en t'appuyant sur le texte.
\begin{enumerate}
\item Worum geht es in diesem Text ?
\item Wer war Johann Wolfgang von Goethe ?
\item Was zeigen die Gemälde von Caspar David Friedrich ?
\item Welche Stadt war das Zentrum der klassischen Musik ?
\item Wer war Thomas Mann ?
\item Was findet jedes Jahr in Berlin, Wien und Zürich statt ?
\end{enumerate}
\tcblower
\textbf{Lösungen (réponses en phrases complètes) :}
\begin{enumerate}
\item \all{Der Text handelt von der Kultur, der Literatur und der Kunst in Deutschland, Österreich und der Schweiz.} (Le texte traite de la culture, de la littérature et de l'art dans les pays germanophones.)
\item \all{Goethe war der berühmteste Schriftsteller Deutschlands. Er schrieb den Roman „Die Leiden des jungen Werthers“ und das Drama „Faust“.}
\item \all{Seine Gemälde zeigen die beeindruckende Natur der Romantik : einsame Menschen, neblige Berge und stille Wälder.}
\item \all{Wien, die Hauptstadt Österreichs, war das Zentrum der klassischen Musik.}
\item \all{Thomas Mann war ein deutscher Schriftsteller und der Träger des Nobelpreises von 1929. Er schrieb „Buddenbrooks“.}
\item \all{Jedes Jahr finden dort zahlreiche Ausstellungen und Lesungen statt.}
\end{enumerate}
Méthode : repère d'abord le mot-clé de la question (\all{wer}, \all{was}, \all{welche Stadt}), localise la phrase du texte, puis reformule en phrase complète avec sujet + verbe en deuxième position.
\end{exoresolu}

\courssub{Carte mentale et frise chronologique}

\begin{popfigure}
\begin{center}\tcbox[colback=popblue,colframe=popblue,coltext=white,arc=9pt,boxsep=4pt,left=6pt,right=6pt]{\bfseries\small Kultur}\end{center}
\vspace{-2pt}
\begin{tcbraster}[raster columns=2,raster equal height=rows,raster column skip=6pt,raster row skip=6pt,enhanced,blankest]
\begin{tcolorbox}[enhanced,colback=poporange,colframe=poporange,coltext=white,boxrule=0.9pt,arc=3pt,left=4pt,right=4pt,top=3pt,bottom=3pt,boxsep=1pt,halign=left]\bfseries\scriptsize Malerei \tiny das Gemälde, der Maler\end{tcolorbox}
\begin{tcolorbox}[enhanced,colback=popgreen,colframe=popgreen,coltext=white,boxrule=0.9pt,arc=3pt,left=4pt,right=4pt,top=3pt,bottom=3pt,boxsep=1pt,halign=left]\bfseries\scriptsize Literatur \tiny der Roman, das Gedicht\end{tcolorbox}
\begin{tcolorbox}[enhanced,colback=poppurple,colframe=poppurple,coltext=white,boxrule=0.9pt,arc=3pt,left=4pt,right=4pt,top=3pt,bottom=3pt,boxsep=1pt,halign=left]\bfseries\scriptsize Museen \tiny die Ausstellung, Berlin\end{tcolorbox}
\begin{tcolorbox}[enhanced,colback=poppink,colframe=poppink,coltext=white,boxrule=0.9pt,arc=3pt,left=4pt,right=4pt,top=3pt,bottom=3pt,boxsep=1pt,halign=left]\bfseries\scriptsize Personen \tiny Goethe, Friedrich, Mann\end{tcolorbox}
\end{tcbraster}

\legende{Carte mentale du vocabulaire : la culture et ses quatre branches.}
\end{popfigure}

\begin{popfigure}
\begin{tikzpicture}[scale=1]
\draw[-{Stealth}, thick, popdark] (0,0) -- (10.5,0);
\foreach \x/\jahr in {0.5/1749, 3.2/1800, 6/1929, 9.5/heute}
  \draw[thick, popdark] (\x,0.12) -- (\x,-0.12) node[below=2pt, font=\small]{\jahr};
\node[above=6pt, align=center, font=\small, text=popblue] at (0.5,0.1) {Goethe\\ geboren};
\node[above=6pt, align=center, font=\small, text=poporange] at (3.2,0.1) {Romantik\\ C. D. Friedrich};
\node[above=6pt, align=center, font=\small, text=popgreen] at (6,0.1) {Thomas Mann\\ Nobelpreis};
\node[above=6pt, align=center, font=\small, text=poppurple] at (9.5,0.1) {moderne\\ Literatur};
\end{tikzpicture}
\legende{Frise chronologique : de Goethe (1749-1832) à la littérature d'aujourd'hui.}
\end{popfigure}

\begin{aretenir}[Ce qu'il faut retenir de cette section]
\begin{itemize}
\item Un texte se lit en deux temps : lecture \cle{globale} (thème, personnes, lieu, époque) puis lecture \cle{détaillée} (mots-clés, structures grammaticales).
\item Le vocabulaire des arts : \all{die Kunst}, \all{der Maler}, \all{das Gemälde}, \all{der Schriftsteller}, \all{der Roman}, \all{das Gedicht}, \all{die Ausstellung}.
\item Les formes en \all{-d} et \all{-end} devant les noms sont des \cle{participes épithètes} ; les compléments en \all{des/der} sont des \cle{génitifs} : ils seront étudiés en détail dans les sections suivantes.
\item Les verbes forts changent de radical au prétérit : \all{schreiben -- schrieb}, \all{finden -- fand}, \all{sein -- war}.
\end{itemize}
\end{aretenir}

\coursec{Lexique thématique : die Kunst und die Literatur}

Dans la section précédente, nous avons lu et commenté le texte \all{Kunst und Literatur im deutschsprachigen Raum}. Ce texte regorgeait de noms comme \all{der Maler}, \all{das Gemälde}, \all{der Schriftsteller} ou \all{das Gedicht}. Pour bien comprendre un tel texte — et surtout pour pouvoir en parler et en écrire — il faut maintenant \cle{maîtriser le vocabulaire} des arts et de la littérature. C'est l'objet de cette section : nous allons observer ces mots, les classer par champs lexicaux, apprendre leurs articles et leurs pluriels, découvrir leurs familles de mots, puis nous entraîner à les employer dans des phrases complètes.

\begin{methode}[Comment apprendre un nom allemand ?]
En allemand, un nom ne s'apprend jamais seul. Pour chaque nom, retiens toujours trois informations :
\begin{enumerate}
\item son \cle{article} (der, die, das), car le genre est imprévisible ;
\item son \cle{pluriel}, noté entre parenthèses : \all{das Museum (Museen)} ;
\item sa \cle{traduction} française et, si possible, un exemple en contexte.
\end{enumerate}
Exemple : \all{der Roman (-e)} = le roman, pluriel \all{die Romane}. Un nom appris sans article est un nom à moitié appris !
\end{methode}

\courssub{Champ 1 : Les arts et la peinture}

Observons d'abord une phrase tirée de notre texte support :

\all{Der Maler zeigt sein neues Gemälde in einer Ausstellung.} « Le peintre présente son nouveau tableau dans une exposition. »

Trois noms du champ des arts y apparaissent : \all{der Maler} (la personne), \all{das Gemälde} (l'œuvre), \all{die Ausstellung} (le lieu d'exposition). Voici le champ lexical complet :

\begin{center}
\begin{tabularx}{\linewidth}{|l|X|X|}
\hline
\rowcolor{popblueL} \textbf{Deutsch} & \textbf{Pluriel} & \textbf{Français} \\
\hline
die Kunst & die Künste & l'art (les arts) \\
\hline
der Künstler & die Künstler & l'artiste (homme) \\
\hline
die Künstlerin & die Künstlerinnen & l'artiste (femme) \\
\hline
der Maler & die Maler & le peintre \\
\hline
die Malerin & die Malerinnen & la peintre (femme) \\
\hline
das Gemälde & die Gemälde & le tableau, la peinture (œuvre) \\
\hline
das Bild & die Bilder & l'image, le tableau (sens large) \\
\hline
die Ausstellung & die Ausstellungen & l'exposition \\
\hline
das Museum & die Museen & le musée \\
\hline
die Malerei & die Malereien & la peinture (l'art de peindre) \\
\hline
\end{tabularx}
\end{center}

\begin{exemplebox}[Le champ des arts en contexte]
\all{Das Museum in der Stadt zeigt eine Ausstellung moderner Kunst.} « Le musée de la ville présente une exposition d'art moderne. »

\all{Die Malerin hat viele Gemälde von der Landschaft Kameruns gemalt.} « La femme peintre a peint de nombreux tableaux des paysages du Cameroun. »

\all{Am Wochenende besuchen wir das Museum und bewundern die Bilder.} « Le week-end, nous visitons le musée et admirons les tableaux. »
\end{exemplebox}

\begin{attention}[das Gemälde ou das Bild ?]
Les deux se traduisent par « tableau », mais : \all{das Gemälde} désigne une \cle{œuvre peinte} (un tableau de peintre, souvent de valeur), tandis que \all{das Bild} a un sens plus large : image, photo, illustration. On dit \all{ein Gemälde von Picasso} mais \all{ein Foto / ein Bild von meiner Familie}. Au pluriel, attention : \all{das Bild → die Bilder} (pluriel en \all{-er}, fréquent chez les noms neutres).
\end{attention}

\courssub{Champ 2 : La littérature}

Observons maintenant cette phrase :

\all{Die Schriftstellerin hat einen berühmten Roman geschrieben.} « L'écrivaine a écrit un roman célèbre. »

\begin{center}
\begin{tabularx}{\linewidth}{|l|X|X|}
\hline
\rowcolor{popblueL} \textbf{Deutsch} & \textbf{Pluriel} & \textbf{Français} \\
\hline
die Literatur & die Literaturen & la littérature \\
\hline
der Schriftsteller & die Schriftsteller & l'écrivain (homme) \\
\hline
die Schriftstellerin & die Schriftstellerinnen & l'écrivaine \\
\hline
der Autor & die Autoren & l'auteur \\
\hline
der Dichter & die Dichter & le poète \\
\hline
der Roman & die Romane & le roman \\
\hline
das Gedicht & die Gedichte & le poème \\
\hline
die Erzählung & die Erzählungen & le récit, la nouvelle \\
\hline
das Theaterstück & die Theaterstücke & la pièce de théâtre \\
\hline
das Buch & die Bücher & le livre \\
\hline
\end{tabularx}
\end{center}

\begin{exemplebox}[Le champ de la littérature en contexte]
\all{Goethe war ein großer Dichter und Schriftsteller.} « Goethe était un grand poète et écrivain. »

\all{In der Schule lesen wir Gedichte und Erzählungen.} « À l'école, nous lisons des poèmes et des récits. »

\all{Der Autor hat ein neues Theaterstück veröffentlicht.} « L'auteur a publié une nouvelle pièce de théâtre. »
\end{exemplebox}

\begin{attention}[Faux ami : das Gedicht n'est pas « l'histoire » !]
\all{das Gedicht (-e)} signifie \cle{le poème}. Le mot français « histoire » se dit \all{die Geschichte (-n)}. Ne confonds jamais :

\all{Sie liest ein Gedicht.} « Elle lit un poème. »

\all{Sie erzählt eine Geschichte.} « Elle raconte une histoire. »

Astuce : \all{das Gedicht} vient du verbe \all{dichten} (composer des poèmes) ; \all{die Geschichte} vient de \all{geschehen} (se passer, arriver) — l'histoire, c'est ce qui s'est passé.
\end{attention}

\courssub{Champ 3 : Verbes et expressions}

\begin{center}
\begin{tabularx}{\linewidth}{|l|X|X|}
\hline
\rowcolor{popblueL} \textbf{Deutsch} & \textbf{Français} & \textbf{Exemple} \\
\hline
malen & peindre & \all{Er malt ein Bild.} \\
\hline
schreiben (schrieb, geschrieben) & écrire & \all{Sie schreibt einen Roman.} \\
\hline
dichten & composer des poèmes & \all{Der Dichter dichtet.} \\
\hline
lesen (las, gelesen) & lire & \all{Wir lesen ein Gedicht.} \\
\hline
ausstellen & exposer & \all{Das Museum stellt Gemälde aus.} \\
\hline
bewundern & admirer & \all{Wir bewundern die Kunst.} \\
\hline
veröffentlichen & publier & \all{Er veröffentlicht ein Buch.} \\
\hline
\end{tabularx}
\end{center}

\begin{definition}[Expressions utiles à retenir par cœur]
\all{ein Buch veröffentlichen} = publier un livre ; \all{eine Ausstellung besuchen} = visiter une exposition ; \all{ein Gedicht schreiben / dichten} = écrire un poème ; \all{ein Bild malen} = peindre un tableau.

Exemple : \all{Gestern haben wir eine Ausstellung besucht und viele Gemälde bewundert.} « Hier, nous avons visité une exposition et admiré beaucoup de tableaux. »
\end{definition}

\begin{attention}[Verbes forts et verbe à particule : deux pièges]
\all{schreiben} et \all{lesen} sont des \cle{verbes forts} : \all{schreiben – schrieb – hat geschrieben} ; \all{lesen – las – hat gelesen}. Au Präsens, \all{lesen} change de voyelle : \all{ich lese}, mais \all{er/sie liest}. Ne dis jamais \all{er lesen} ni \all{er leset} !

De plus, \all{ausstellen} est un \cle{verbe à particule séparable} : la particule \all{aus-} va en fin de phrase au Präsens : \all{Das Museum \textbf{stellt} Gemälde \textbf{aus}.} Au Perfekt, le \all{ge-} s'insère entre la particule et le verbe : \all{Das Museum hat Gemälde ausgestellt.}
\end{attention}

\courssub{Champ 4 : Les qualificatifs}

Pour parler d'une œuvre, il faut pouvoir la qualifier :

\begin{center}
\begin{tabularx}{\linewidth}{|X|X|X|}
\hline
\rowcolor{popblueL} \textbf{Deutsch} & \textbf{Français} & \textbf{Exemple} \\
\hline
berühmt & célèbre & \all{ein berühmter Maler} \\
\hline
bekannt & connu & \all{ein bekannter Autor} \\
\hline
beeindruckend & impressionnant & \all{ein beeindruckendes Gemälde} \\
\hline
kreativ & créatif & \all{eine kreative Künstlerin} \\
\hline
modern & moderne & \all{moderne Kunst} \\
\hline
klassisch & classique & \all{klassische Literatur} \\
\hline
romantisch & romantique & \all{ein romantisches Gedicht} \\
\hline
\end{tabularx}
\end{center}

\begin{attention}[berühmt ou bekannt ?]
\all{berühmt} = célèbre (connu de tous, dans le monde entier) ; \all{bekannt} = connu (dans un cercle plus restreint). \all{Goethe ist berühmt}, mais \all{mein Lehrer ist in der Schule bekannt}. Ne traduis pas « connu » par \all{gekannt} : \all{kennen} est le verbe « connaître », l'adjectif est \all{bekannt}.
\end{attention}

\courssub{Le féminin des noms de personnes : la terminaison -in}

\begin{propriete}[Le féminin en -in]
En allemand, les noms de métiers et de personnes forment leur \cle{féminin} en ajoutant \all{-in} au masculin ; l'article devient \all{die}. Au pluriel, le féminin prend \all{-innen}.

\begin{center}
\begin{tabular}{|l|l|l|}
\hline
\rowcolor{popblueL} \textbf{Masculin} & \textbf{Féminin} & \textbf{Pluriel féminin} \\
\hline
der Maler & die Malerin & die Malerinnen \\
\hline
der Künstler & die Künstlerin & die Künstlerinnen \\
\hline
der Schriftsteller & die Schriftstellerin & die Schriftstellerinnen \\
\hline
der Dichter & die Dichterin & die Dichterinnen \\
\hline
der Autor & die Autorin & die Autorinnen \\
\hline
der Lehrer & die Lehrerin & die Lehrerinnen \\
\hline
\end{tabular}
\end{center}

Exemple : \all{Die Schriftstellerin hat einen berühmten Roman geschrieben.} « L'écrivaine a écrit un roman célèbre. »
\end{propriete}

\begin{attention}[Erreur fréquente des francophones]
En français, « le peintre » peut désigner un homme ou une femme. En allemand, c'est impossible : il faut choisir \all{der Maler} ou \all{die Malerin}. Et n'oublie pas l'accord de l'adjectif et du possessif : \all{die Malerin zeigt \textbf{ihr} neues Gemälde} (et non \all{sein}).
\end{attention}

\courssub{Les familles de mots}

L'allemand construit des familles entières de mots autour d'une racine verbale. Connaître la racine, c'est deviner le sens des dérivés.

\begin{popfigure}
\begin{center}
\begin{tikzpicture}[
  grow=down,
  level distance=1.6cm,
  level 1/.style={sibling distance=4.6cm},
  level 2/.style={sibling distance=3.2cm},
  every node/.style={draw=popblue, rounded corners, fill=popblueL, align=center, font=\small, inner sep=4pt},
  edge from parent/.style={draw=popdark, thick, -{Stealth}}
]
\node[fill=poporangeL, draw=poporange, font=\bfseries, align=center]{malen\\ (peindre)}
  child { node[align=center]{der Maler\\ die Malerin\\ (le peintre)} }
  child { node[align=center]{das Gemälde\\ (le tableau)} }
  child { node[align=center]{die Malerei\\ (la peinture)} };
\end{tikzpicture}
\end{center}
\legende{Famille de mots autour du verbe \all{malen}}
\end{popfigure}

De la même façon :

\begin{center}
\begin{tabularx}{\linewidth}{|X|X|X|}
\hline
\rowcolor{popblueL} \textbf{Verbe} & \textbf{Personne} & \textbf{Œuvre / résultat} \\
\hline
malen (peindre) & der Maler / die Malerin & das Gemälde, die Malerei \\
\hline
schreiben (écrire) & der Schriftsteller / die Schriftstellerin & die Schrift (l'écriture), der Roman \\
\hline
dichten (composer) & der Dichter / die Dichterin & das Gedicht (le poème) \\
\hline
erzählen (raconter) & der Erzähler & die Erzählung (le récit) \\
\hline
ausstellen (exposer) & der Aussteller & die Ausstellung (l'exposition) \\
\hline
\end{tabularx}
\end{center}

\begin{aretenir}[Une logique à retenir]
Le suffixe \all{-er} désigne souvent \cle{la personne} qui fait l'action (\all{malen → der Maler}) ; le suffixe \all{-ung} désigne \cle{l'action ou son résultat} (\all{erzählen → die Erzählung} ; \all{ausstellen → die Ausstellung}) ; le préfixe \all{Ge-} avec un neutre désigne souvent \cle{le produit} de l'action (\all{dichten → das Gedicht} ; \all{malen → das Gemälde}).
\end{aretenir}

\begin{savaistu}[Dichter ou Schriftsteller ?]
Le mot \all{Dichter} désigne celui qui « dichtet », c'est-à-dire qui compose des poèmes : c'est le poète. \all{Schriftsteller} (littéralement « celui qui pose des écrits ») désigne l'écrivain en général, surtout l'auteur de prose. Johann Wolfgang von Goethe était les deux à la fois : \all{Dichter} (il a écrit de célèbres poèmes) et \all{Schriftsteller} (il a écrit des romans et des pièces de théâtre comme \all{Faust}). En allemand, on admire d'ailleurs le pays des \all{Dichter und Denker} — « les poètes et les penseurs ».
\end{savaistu}

\courssub{Texte d'application : Ein Besuch im Museum}

\begin{textebox}[Ein Besuch im Museum]
Am Samstag besucht die Klasse 12a das Nationalmuseum in Yaoundé. Dort gibt es eine große Ausstellung moderner Kunst aus Kamerun und Deutschland. Die Schüler bewundern die Gemälde einer bekannten Malerin aus Douala. Ihre Bilder zeigen das Leben in der Stadt: die Märkte, die Menschen, die Farben. „Diese Malerei ist wirklich beeindruckend“, sagt Amina. „Die Künstlerin ist sehr kreativ!“

Der Museumsführer erzählt auch von einem berühmten kamerunischen Schriftsteller. Er hat viele Romane und Erzählungen geschrieben und mehrere Bücher veröffentlicht. In einer Glasvitrine sehen die Schüler seine ersten Bücher. „Lest ihr gern Gedichte?“, fragt der Führer. „Ja!“, antwortet Paul. „Ich lese gern romantische Gedichte, aber ich schreibe auch selbst welche.“ Alle lachen.

Am Ende kaufen die Schüler Postkarten mit den schönsten Bildern der Ausstellung. Amina sagt: „Kunst und Literatur zeigen uns die Seele eines Landes.“ Der Lehrer nickt: „Das ist ein schöner Satz für euren Aufsatz am Montag!“
\end{textebox}

\textbf{Glossaire :} \all{der Museumsführer (-)} = le guide du musée ; \all{die Glasvitrine (-n)} = la vitrine ; \all{die Postkarte (-n)} = la carte postale ; \all{die Seele (-n)} = l'âme ; \all{der Aufsatz (Aufsätze)} = la rédaction ; \all{nicken} = hocher la tête ; \all{selbst} = soi-même ; \all{welche} (ici) = en (quelques-uns).

\textbf{Questions de compréhension :}
\begin{enumerate}
\item \all{Was besucht die Klasse am Samstag?}
\item \all{Wessen Gemälde bewundern die Schüler?}
\item \all{Was hat der Schriftsteller veröffentlicht?}
\item \all{Was liest Paul gern?}
\item \all{Was kaufen die Schüler am Ende?}
\end{enumerate}

\begin{corrige}[Questions de compréhension du texte]
1. \all{Die Klasse besucht am Samstag das Nationalmuseum in Yaoundé.} \quad 2. \all{Sie bewundern die Gemälde einer bekannten Malerin aus Douala.} \quad 3. \all{Er hat mehrere Bücher veröffentlicht (viele Romane und Erzählungen).} \quad 4. \all{Paul liest gern romantische Gedichte.} \quad 5. \all{Am Ende kaufen die Schüler Postkarten mit den schönsten Bildern der Ausstellung.}
\end{corrige}

\courssub{Entraînement}

\begin{exoresolu}[Complète avec le bon mot (article et forme corrects)]
1. \all{Goethe war ein berühmter \dots\ (poète).} \quad 2. \all{Die \dots\ (peintre, femme) zeigt ihr neues \dots\ (tableau).} \quad 3. \all{Wir besuchen eine \dots\ (exposition) moderner Kunst.} \quad 4. \all{Sie \dots\ (lit) gern Romane.} \quad 5. \all{Der Autor hat ein Buch \dots\ (publié).}
\tcblower
1. \all{Dichter} — masculin, nominatif après \all{ein}. \quad 2. \all{Malerin} (féminin en \all{-in}) et \all{Gemälde} (neutre, accusatif : \all{ihr neues Gemälde}). \quad 3. \all{Ausstellung} — \all{eine Ausstellung} (féminin). \quad 4. \all{liest} — verbe fort \all{lesen}, changement de voyelle \all{e → ie} à la 3\up{e} personne. \quad 5. \all{veröffentlicht} — participe passé de \all{veröffentlichen} (verbe faible, sans \all{ge-} car il commence par le préfixe inséparable \all{ver-}).
\end{exoresolu}

\begin{exercice}[À toi de jouer]
1. Donne le féminin : \all{der Künstler}, \all{der Autor}, \all{der Dichter}.
2. Donne le pluriel : \all{das Museum}, \all{das Bild}, \all{der Roman}, \all{das Gedicht}.
3. Traduis en allemand : « L'écrivaine a publié un roman célèbre. » / « Nous admirons les tableaux dans le musée. »
4. Construis la famille de mots de \all{erzählen} (verbe, personne, résultat).
\end{exercice}

\begin{corrige}[Exercice : À toi de jouer]
1. \all{die Künstlerin}, \all{die Autorin}, \all{die Dichterin}. \quad 2. \all{die Museen}, \all{die Bilder}, \all{die Romane}, \all{die Gedichte}. \quad 3. \all{Die Schriftstellerin hat einen berühmten Roman veröffentlicht.} / \all{Wir bewundern die Gemälde im Museum.} \quad 4. \all{erzählen → der Erzähler → die Erzählung}.
\end{corrige}

\begin{aretenir}[L'essentiel du lexique]
\begin{itemize}
\item Chaque nom s'apprend avec son \cle{article} et son \cle{pluriel} : \all{das Gemälde (-)}, \all{die Ausstellung (-en)}, \all{das Museum (Museen)}.
\item Le féminin des personnes se forme en \all{-in} : \all{der Maler → die Malerin}.
\item \all{das Gedicht} = le poème ; \all{die Geschichte} = l'histoire (faux ami !).
\item Les familles de mots tournent autour du verbe : \all{malen → der Maler → das Gemälde → die Malerei}.
\item Verbes forts à connaître : \all{schreiben – schrieb – geschrieben} ; \all{lesen – las – gelesen} (\all{er liest}).
\item Verbe à particule : \all{ausstellen} → \all{Das Museum stellt Gemälde aus.}
\end{itemize}
\end{aretenir}

\coursec{Grammaire 1 : les participes présent et passé en épithète}

Après avoir enrichi notre vocabulaire thématique autour des arts et de la littérature, nous allons maintenant examiner de près deux outils grammaticaux indispensables pour comprendre et produire des phrases complexes et élégantes en allemand : \cle{les participes présent et passé employés comme épithètes} (adjectifs verbaux placés devant le nom). Dans le texte de la section précédente, vous avez rencontré des tournures comme \all{der schreibende Dichter} (le poète qui écrit), \all{das beeindruckende Gemälde} (le tableau impressionnant) ou \all{der berühmte, viel gelesene Roman} (le roman célèbre, très lu). Ces formes condensent une proposition relative en un seul mot (ou groupe de mots) placé \emph{avant} le nom, ce qui est une caractéristique majeure de la syntaxe allemande écrite (presse, littérature, administration). Maîtriser cette structure est crucial pour réussir l'épreuve de thème/version au Baccalauréat et pour accéder à une expression écrite de niveau B2.

\courssub{Observation à partir du texte}

Reprenons les exemples repérés dans notre texte support \emph{Kunst und Literatur im deutschsprachigen Raum} pour en analyser la structure.

\begin{textebox}[Extrait du texte support — Observation grammaticale]
In der Ausstellung sieht man \textbf{den schreibenden Dichter} und \textbf{die malende Künstlerin}. An der Wand hängen \textbf{das beeindruckende Gemälde} und \textbf{die ausgestellten Skulpturen}. Besonders interessant ist \textbf{der berühmte, viel gelesene Roman} von Thomas Mann. Die \textbf{von Goethe geschriebenen} Gedichte sind Klassiker.
\end{textebox}

\begin{itemize}
    \item \textbf{den schreibenden Dichter} / \textbf{die malende Künstlerin} : le verbe \all{schreiben} / \all{malen} apparaît sous la forme \textbf{-nd} (infinitif + d) + désinence adjectivale (\textbf{-en}, \textbf{-e}). L'action est \emph{active} et \emph{simultanée} : le poète \emph{est en train d'écrire}.
    \item \textbf{das beeindruckende Gemälde} : \all{beeindrucken} (impressionner) $\to$ \all{beeindruckend}. Le tableau \emph{impressionne} (sujet actif), action simultanée.
    \item \textbf{die ausgestellten Skulpturen} : \all{ausstellen} (exposer) $\to$ Participe II \all{ausgestellt} + \textbf{-en}. Les sculptures \emph{ont été exposées} (sujet passif), action \emph{antérieure}.
    \item \textbf{viel gelesene Roman} : \all{lesen} (lire) $\to$ Participe II \all{gelesen} + \textbf{-e}. Le roman \emph{a été lu} (passif), modifié par l'adverbe \all{viel}.
    \item \textbf{von Goethe geschriebenen} : Participe II \all{geschrieben} + \textbf{-en}, précédé de son complément d'agent \all{von Goethe}.
\end{itemize}

\courssub{Le participe présent en épithète (Partizip I)}

\begin{propriete}[Formation et valeur du Partizip I épithète]
\textbf{Forme :} Infinitif du verbe + \textbf{-d} + \textbf{désinence adjectivale} (selon le genre, le nombre, le cas et le déterminant).
\all{lesen} $\to$ \all{lesend} + \textbf{-e/-er/-es/-en} $\to$ \all{der lesende Junge}, \all{ein lesender Junge}, \all{dem lesenden Jungen}.

\textbf{Valeur :} il exprime une action \textbf{active} (le nom est le sujet de l'action) et \textbf{simultanée} à l'action principale de la phrase. Il équivaut à une \textbf{proposition relative active} au présent (ou à l'imparfait) introduite par \emph{qui} + verbe actif.
\all{der lesende Junge} = \all{der Junge, der liest} (le garçon qui lit / qui lisait).
\end{propriete}

\begin{exemplebox}[Exemples traduits — Partizip I]
\begin{itemize}
    \item \all{Das schlafende Kind} — « l'enfant qui dort ».
    \item \all{Die singenden Vögel} — « les oiseaux qui chantent ».
    \item \all{Ein spannendes Buch} — « un livre passionnant » (qui passionne $\to$ adjectif lexicalisé).
    \item \all{Die wartenden Menschen} — « les gens qui attendent ».
    \item \all{Der steigende Preis} — « le prix qui augmente ».
\end{itemize}
\end{exemplebox}

\begin{center}
\begin{tabularx}{\linewidth}{|X|X|X|}
\hline
\rowcolor{popblueL} \textbf{Français (relative)} & \textbf{Allemand (relative)} & \textbf{Allemand (Partizip I épithète)} \\
\hline
le garçon qui lit & der Junge, der liest & \textbf{der lesende Junge} \\
les tableaux qui impressionnent & die Bilder, die beeindrucken & \textbf{die beeindruckenden Bilder} \\
une histoire qui passionne & eine Geschichte, die fesselt & \textbf{eine fesselnde Geschichte} \\
les artistes qui peignent & die Künstler, die malen & \textbf{die malenden Künstler} \\
\hline
\end{tabularx}
\end{center}

\begin{popfigure}
\begin{tikzpicture}[
    node distance=1.2cm and 1.5cm,
    box/.style={draw=popblue, thick, rounded corners, fill=popblueL, text width=3.5cm, align=center, font=\small},
    arrow/.style={-Stealth, thick, popdark},
    labelbox/.style={font=\footnotesize, text width=3cm, align=center}
]
\node[box] (inf) {Infinitif : \textbf{lesen}};
\node[box, right=of inf] (part) {+ \textbf{-d} $\to$ \textbf{lesend}};
\node[box, right=of part] (adj) {+ désinence adj. $\to$ \textbf{lesende}};
\node[labelbox, below=0.3cm of inf] (l1) {Valeur : \textbf{actif}};
\node[labelbox, below=0.3cm of part] (l2) {Simultané};
\node[labelbox, below=0.3cm of adj, align=center] (l3){\all{der lesende Mann} \\ (l'homme qui lit)};
\draw[arrow] (inf) -- (part);
\draw[arrow] (part) -- (adj);
\end{tikzpicture}
\legende{Schéma de formation du Partizip I épithète : actif / simultané}
\end{popfigure}

\courssub{Le participe passé en épithète (Partizip II)}

\begin{propriete}[Formation et valeur du Partizip II épithète]
\textbf{Forme :} \textbf{Participe II (forme du Perfekt/Passiv)} + \textbf{désinence adjectivale}.
Verbes faibles : \all{ge-} + radical + \textbf{-t} (ex. \all{gemalt}).
Verbes forts : \all{ge-} + radical modifié + \textbf{-en} (ex. \all{geschrieben}, \all{gelesen}).
Verbes en \textbf{-ieren} ou à \textbf{préfixe inséparable} : \textbf{sans ge-} (voir encadré ci-dessous).
Puis ajout de la désinence : \all{das gemalte Bild}, \all{der geschriebene Brief}.

\textbf{Valeur :} il exprime une action \textbf{passive} (le nom est l'objet/patient de l'action) et \textbf{antérieure} (achevée). Il équivaut à une \textbf{proposition relative au passif} (présent, passé composé ou imparfait).
\all{das geschriebene Buch} = \all{das Buch, das geschrieben wurde} (le livre qui a été écrit).
\end{propriete}

\begin{attention}[Cas particuliers : pas de \textit{ge-} au Participe II]
Deux groupes de verbes ne prennent \textbf{jamais} le préfixe \all{ge-} au Participe II (donc pas non plus en épithète) :
\begin{enumerate}
    \item \textbf{Verbes en \textit{-ieren}} (souvent d'origine latine ou française) : \all{studieren $\to$ studiert}, \all{illustrieren $\to$ illustriert}, \all{kopieren $\to$ kopiert}, \all{organisieren $\to$ organisiert}.
    \item \textbf{Verbes à préfixe inséparable} (\all{be-, emp-, ent-, er-, ge-, ver-, zer-, miss-}) : \all{besuchen $\to$ besucht}, \all{verstehen $\to$ verstanden}, \all{erzählen $\to$ erzählt}, \all{gefallen $\to$ gefallen}.
\end{enumerate}
\emph{Exemples :} \all{das illustrierte Buch} (le livre illustré), \all{der besuchte Künstler} (l'artiste que l'on visite), \all{die erzählte Geschichte} (l'histoire racontée).
\end{attention}

\begin{exemplebox}[Exemples traduits — Partizip II]
\begin{itemize}
    \item \all{das gemalte Bild} — « le tableau (qui a été) peint ».
    \item \all{der geschriebene Brief} — « la lettre (qui a été) écrite ».
    \item \all{die viel gelesene Zeitung} — « le journal très lu » (l'adverbe \all{viel} se place avant le participe).
    \item \all{die im Museum ausgestellten Werke} — « les œuvres exposées au musée » (le groupe prépositionnel \all{im Museum} s'insère avant le participe).
    \item \all{das organisierte Festival} — « le festival organisé » (verbe en \textit{-ieren}, pas de \textit{ge-}).
    \item \all{der verstandene Text} — « le texte compris » (préfixe inséparable \textit{ver-}, pas de \textit{ge-}).
\end{itemize}
\end{exemplebox}

\begin{popfigure}
\begin{tikzpicture}[
    node distance=1.2cm and 1.5cm,
    box/.style={draw=poporange, thick, rounded corners, fill=poporangeL, text width=3.5cm, align=center, font=\small},
    arrow/.style={-Stealth, thick, popdark},
    labelbox/.style={font=\footnotesize, text width=3cm, align=center}
]
\node[box] (inf) {Verbe fort : \textbf{schreiben}};
\node[box, right=of inf] (part) {Partizip II : \textbf{geschrieben}};
\node[box, right=of part] (adj) {+ désinence adj. $\to$ \textbf{geschriebene}};
\node[labelbox, below=0.3cm of inf] (l1) {Valeur : \textbf{passif}};
\node[labelbox, below=0.3cm of part] (l2) {Antérieur (achevé)};
\node[labelbox, below=0.3cm of adj, align=center] (l3){\all{das geschriebene Buch} \\ (le livre écrit)};
\draw[arrow] (inf) -- (part);
\draw[arrow] (part) -- (adj);
\end{tikzpicture}
\legende{Schéma de formation du Partizip II épithète : passif / antérieur}
\end{popfigure}

\courssub{Transformation : proposition relative $\rightarrow$ participe épithète}

C'est la compétence clé pour le thème (français $\to$ allemand) et pour la compression de texte.

\begin{center}
\begin{tabularx}{\linewidth}{|X|X|X|}
\hline
\rowcolor{popgreenL} \textbf{Relative française} & \textbf{Relative allemande} & \textbf{Participe épithète allemand} \\
\hline
le peintre \textbf{qui peint} & der Maler, \textbf{der malt} & \textbf{der malende Maler} (Part. I) \\
le tableau \textbf{qui est peint} & das Bild, \textbf{das gemalt wird} & \textbf{das gemalte Bild} (Part. II) \\
le roman \textbf{qui a été lu} & der Roman, \textbf{der gelesen wurde} & \textbf{der gelesene Roman} (Part. II) \\
l'artiste \textbf{qu'on visite} & der Künstler, \textbf{der besucht wird} & \textbf{der besuchte Künstler} (Part. II, pas de \textit{ge-}) \\
la fillette \textbf{qui rit} & das Mädchen, \textbf{das lacht} & \textbf{das lachende Mädchen} (Part. I) \\
\hline
\end{tabularx}
\end{center}

\begin{methode}[Traduire une relative en participe épithète (Bac / Thème)]
\emph{Étape 1 :} identifier le verbe de la relative et son antécédent (le nom auquel se rapporte \emph{qui/que}).
\emph{Étape 2 :} déterminer la \textbf{diathèse} (active ou passive) et la \textbf{relation temporelle}.
\begin{itemize}
    \item Relative \textbf{active} + action \textbf{simultanée} (présent/imparfait) $\to$ \textbf{Partizip I} (\textit{infinitif + d + désinence}).
    \item Relative \textbf{passive} (ou tournure française de sens passif : « qu'on a lu », « qu'on visite ») + action \textbf{antérieure ou subie} $\to$ \textbf{Partizip II} (\textit{Participe II + désinence}).
\end{itemize}
\emph{Étape 3 :} construire le Participe II correctement (attention : \textit{ge-} ou pas de \textit{ge-}, verbe fort ou faible).
\emph{Étape 4 :} ajouter la \textbf{désinence adjectivale} selon le déterminant, le genre, le nombre et le cas.
\emph{Étape 5 :} placer les compléments (adverbes, groupes prépositionnels, \textit{von} + datif pour l'agent) \textbf{avant} le participe.
\end{methode}

\courssub{La déclinaison obligatoire des participes épithètes}

C'est l'erreur n°1 des francophones : \textbf{oublier la désinence}. En allemand, un participe épithète se comporte \textbf{exactement comme un adjectif}. Il \textbf{doit} porter la marque du genre, du nombre et du cas.

\begin{propriete}[Déclinaison du participe épithète]
Le participe suit la déclinaison \textbf{faible} (après \all{der/die/das}), \textbf{mixte} (après \all{ein/kein/mein}...) ou \textbf{forte} (sans déterminant), exactement comme un adjectif épithète ordinaire.
\end{propriete}

\begin{center}
\begin{tabular}{|l|l|l|l|}
\hline
\rowcolor{poppurpleL} \textbf{Cas} & \textbf{Article défini} & \textbf{Article indéfini} & \textbf{Possessif} \\
\hline
Nom. masc. & der geschrieben\textbf{e} Brief & ein geschrieben\textbf{er} Brief & mein geschrieben\textbf{er} Brief \\
Acc. masc. & den geschrieben\textbf{en} Brief & einen geschrieben\textbf{en} Brief & meinen geschrieben\textbf{en} Brief \\
Dat. masc. & dem geschrieben\textbf{en} Brief & einem geschrieben\textbf{en} Brief & meinem geschrieben\textbf{en} Brief \\
Gén. masc. & des geschrieben\textbf{en} Briefes & eines geschrieben\textbf{en} Briefes & meines geschrieben\textbf{en} Briefes \\
\hline
Nom./Acc. fém. & die geschrieben\textbf{e} Schrift & eine geschrieben\textbf{e} Schrift & meine geschrieben\textbf{e} Schrift \\
Dat./Gén. fém. & der geschrieben\textbf{en} Schrift & einer geschrieben\textbf{en} Schrift & meiner geschrieben\textbf{en} Schrift \\
\hline
Nom./Acc. neutre & das geschrieben\textbf{e} Buch & ein geschrieben\textbf{es} Buch & mein geschrieben\textbf{es} Buch \\
Dat. neutre & dem geschrieben\textbf{en} Buch & einem geschrieben\textbf{en} Buch & meinem geschrieben\textbf{en} Buch \\
Gén. neutre & des geschrieben\textbf{en} Buches & eines geschrieben\textbf{en} Buches & meines geschrieben\textbf{en} Buches \\
\hline
Pluriel (tous cas) & die geschrieben\textbf{en} Bücher & keine geschrieben\textbf{en} Bücher & meine geschrieben\textbf{en} Bücher \\
\hline
\end{tabular}
\end{center}

\begin{attention}[Piège mortel : « ein lesend Junge »]
\emph{Jamais} de participe nu sans désinence !
\begin{itemize}
    \item Faux : \all{*ein lesend Junge}, \all{*das gelesen Buch}, \all{*ein berühmt Maler}.
    \item Correct : \all{ein lesend\textbf{er} Junge}, \all{das gelesen\textbf{e} Buch}, \all{ein berühmt\textbf{er} Maler}.
\end{itemize}
Même si le participe se termine déjà par une voyelle (Partizip I \all{lesende}, Partizip II faible \all{gemalte}), la désinence s'ajoute \textbf{en plus} selon le contexte : \all{der lesend\textbf{e} Junge}, \all{die gemalt\textbf{en} Bilder}.
\end{attention}

\courssub{Entraînement : exercice résolu}

\begin{exoresolu}[Transformation de relatives en participes épithètes]
\textbf{Énoncé :} Transformez les phrases suivantes en remplaçant la proposition relative (entre crochets) par un participe épithète correctement décliné. Traduisez la phrase finale.

1. Ich bewundere den Künstler, [der dieses Bild malt]. (Accusatif, masc.)
2. Wir besuchen die Ausstellung, [die im Museum gezeigt wird]. (Accusatif, fém.)
3. Er liest ein Buch, [das von Kafka geschrieben wurde]. (Accusatif, neutre)
4. Die Kinder, [die dort spielen], sind laut. (Nominatif, pluriel)
5. Sie kennt den Dichter, [der sehr berühmt ist]. (Accusatif, masc. — Attention : verbe d'état !)

\tcblower
\textbf{Solution détaillée :}

1. Relative active, simultanée $\to$ \textbf{Partizip I}. Verbe \all{malen} $\to$ \all{malend}. Accusatif masc. défini (\all{den}) $\to$ désinence \textbf{-en}. Le complément d'objet \all{dieses Bild} se place avant le participe.
   $\to$ \all{Ich bewundere \textbf{den dieses Bild malenden Künstler}.}
   \textit{Traduction : « J'admire l'artiste qui peint ce tableau. »}

2. Relative passive, simultanée (présent passif) $\to$ \textbf{Partizip II}. Verbe \all{zeigen} $\to$ \all{gezeigt}. Accusatif fém. défini (\all{die}) $\to$ désinence \textbf{-e}.
   $\to$ \all{Wir besuchen \textbf{die im Museum gezeigte Ausstellung}.}
   \textit{Traduction : « Nous visitons l'exposition présentée au musée. »}

3. Relative passive, antérieure (passif au Präteritum) $\to$ \textbf{Partizip II}. Verbe fort \all{schreiben} $\to$ \all{geschrieben}. Accusatif neutre indéfini (\all{ein}) $\to$ déclinaison mixte $\to$ désinence \textbf{-es}.
   $\to$ \all{Er liest \textbf{ein von Kafka geschriebenes Buch}.}
   \textit{Traduction : « Il lit un livre écrit par Kafka. »}

4. Relative active, simultanée $\to$ \textbf{Partizip I}. Verbe \all{spielen} $\to$ \all{spielend}. Nominatif pluriel défini (\all{die}) $\to$ désinence \textbf{-en}.
   $\to$ \all{\textbf{Die dort spielenden Kinder} sind laut.}
   \textit{Traduction : « Les enfants qui jouent là-bas sont bruyants. »}

5. \textbf{Piège !} \all{sein} (être) est un verbe d'état : il n'a pas de passif et son Partizip I \all{seiend} n'est pas employé comme épithète dans la langue courante. On ne peut donc pas transformer « qui est célèbre » en participe verbal : on utilise directement l'adjectif \all{berühmt}, décliné (accusatif masc. défini $\to$ \textbf{-en}).
   $\to$ \all{Sie kennt \textbf{den berühmten Dichter}.}
   \textit{Traduction : « Elle connaît le poète célèbre. »}
\end{exoresolu}

\courssub{Approfondissement : texte d'application}

Le texte suivant, une critique culturelle fictive, regroupe de nombreux participes épithètes. Lisez-le, repérez-les (Part. I ou Part. II ?), puis répondez aux questions.

\begin{textebox}[Kulturkritik : Die neue Ausstellung im Goethe-Institut Yaoundé]
Das Goethe-Institut Yaoundé präsentiert eine \textbf{faszinierende} Sammlung moderner Kunst. Im Mittelpunkt stehen \textbf{von jungen kamerunischen Künstlern geschaffene} Skulpturen und \textbf{gemalte} Leinwände, die traditionelle Motive mit zeitgenössischen Techniken verbinden. Besonders \textbf{beeindruckend} wirken die \textbf{aus recyceltem Material gefertigten} Figuren, die soziale Themen auf \textbf{ansprechende} Weise darstellen. Ein \textbf{begleitender} Katalog dokumentiert den Entstehungsprozess. Die \textbf{geladenen} Gäste, unter ihnen der in Yaoundé \textbf{residierende} Botschafter, zeigten sich \textbf{begeistert}. Der \textbf{eröffnende} Redner betonte die Wichtigkeit des kulturellen Austauschs. Am Ende gab es \textbf{live gespielte} Musik und \textbf{vorgetragene} Gedichte \textbf{schreibender} Poeten aus Douala.
\end{textebox}

\begin{definition}[Glossaire — Mots difficiles du texte]
\begin{itemize}
    \item \textbf{die Sammlung, -en} : la collection.
    \item \textbf{zeitgenössisch} : contemporain.
    \item \textbf{verbinden} (verband, hat verbunden) : relier, lier.
    \item \textbf{recycelt} : recyclé (Part. II du verbe faible \all{recyceln} ; dans \all{aus recyceltem Material}, désinence datif neutre sans article $\to$ \textbf{-em}).
    \item \textbf{fertigen} (fertigte, hat gefertigt) : fabriquer, confectionner.
    \item \textbf{ansprechend} : attrayant, expressif (Part. I lexicalisé de \all{ansprechen}).
    \item \textbf{die Weise, -n} : la manière ; \all{auf ansprechende Weise} : de manière expressive.
    \item \textbf{dokumentieren} : documenter (verbe en \textit{-ieren} $\to$ Part. II \all{dokumentiert}).
    \item \textbf{der Entstehungsprozess, -e} : le processus de création.
    \item \textbf{begleitend} : d'accompagnement (Part. I de \all{begleiten}).
    \item \textbf{der Botschafter, -} : l'ambassadeur.
    \item \textbf{residierend} : résidant, en poste (Part. I, terme diplomatique).
    \item \textbf{begeistert} : enthousiaste (Part. II adjectivé de \all{begeistern}).
    \item \textbf{der Redner, -} : l'orateur.
    \item \textbf{der kulturelle Austausch} : l'échange culturel.
    \item \textbf{live} (invariable) : en direct.
    \item \textbf{vortragen} (trug vor, hat vorgetragen) : réciter, présenter.
    \item \textbf{der Poet, -en} : le poète.
\end{itemize}
\end{definition}

\begin{exercice}[Questions de compréhension et analyse grammaticale]
\begin{enumerate}
    \item Relevez dans le texte \textbf{6 participes présents (Part. I)} et \textbf{6 participes passés (Part. II)} employés en épithète. Pour chacun, indiquez le nom auquel il se rapporte et sa désinence.
    \item Traduisez en français les groupes nominaux suivants extraits du texte :
    \begin{itemize}
        \item \all{von jungen kamerunischen Künstlern geschaffene Skulpturen}
        \item \all{aus recyceltem Material gefertigte Figuren}
        \item \all{vorgetragene Gedichte schreibender Poeten}
    \end{itemize}
    \item Pourquoi \all{faszinierende}, \all{beeindruckend}, \all{ansprechende}, \all{begleitender}, \all{residierende}, \all{eröffnende} et \all{schreibender} sont-ils des Part. I ? Justifiez par la valeur (actif / simultané).
    \item Pourquoi \all{geschaffene}, \all{gemalte}, \all{gefertigten}, \all{geladenen}, \all{begeistert}, \all{gespielte} et \all{vorgetragene} sont-ils des Part. II ? Justifiez par la valeur (passif / antérieur) et précisez pour chaque verbe s'il est fort ou faible : \all{schaffen} (fort : schuf, geschaffen), \all{malen} (faible : gemalt), \all{fertigen} (faible : gefertigt), \all{laden} (fort : lud, geladen), \all{begeistern} (faible à préfixe inséparable \textit{be-} : begeistert, sans \textit{ge-}), \all{spielen} (faible : gespielt), \all{vortragen} (fort à particule séparable : trug vor, vorgetragen).
    \item \textbf{Production écrite :} rédigez 5 phrases en allemand pour décrire une exposition imaginaire à Douala ou à Berlin, en utilisant \textbf{au moins 3 Part. I} et \textbf{3 Part. II} épithètes différents (avec compléments si possible : adverbe, \textit{von} + datif, groupe prépositionnel).
\end{enumerate}
\end{exercice}

\begin{corrige}[Exercice — éléments de correction]
\textbf{1.} Part. I : \all{faszinierende} (Sammlung, acc. fém. \textbf{-e}), \all{beeindruckend} (attribut après \all{wirken}, pas de désinence), \all{ansprechende} (Weise, acc. fém. \textbf{-e}), \all{begleitender} (Katalog, nom. masc. \textbf{-er}), \all{residierende} (Botschafter, nom. masc. \textbf{-e} après \all{der}), \all{eröffnende} (Redner, nom. masc. \textbf{-e} après \all{der}), \all{schreibender} (Poeten, gén. plur. fort \textbf{-er}). Part. II : \all{geschaffene} (Skulpturen, nom. plur. \textbf{-e} sans article), \all{gemalte} (Leinwände, nom. plur. \textbf{-e}), \all{gefertigten} (Figuren, nom. plur. \textbf{-en} après \all{die}), \all{geladenen} (Gäste, nom. plur. \textbf{-en} après \all{die}), \all{begeistert} (attribut, pas de désinence), \all{gespielte} (Musik, acc. fém. \textbf{-e} sans article), \all{vorgetragene} (Gedichte, acc. plur. \textbf{-e} sans article).

\textbf{2.} « des sculptures créées par de jeunes artistes camerounais » ; « des figures fabriquées à partir de matériaux recyclés » ; « des poèmes récités de poètes qui écrivent / poètes en activité ».

\textbf{5.} Modèle de production : \all{Die Ausstellung zeigt beeindruckende Gemälde junger Künstler. Die von einer Doualaer Malerin gemalten Bilder faszinieren die Besucher. Ein singender Musiker begleitet die Eröffnung. Die aus Holz geschnitzten Figuren erzählen alte Geschichten. Die begeisterten Gäste diskutieren mit den anwesenden Künstlern.}
\end{corrige}

\begin{savaistu}[L'allemand, langue des groupes épithétiques géants]
L'allemand adore empiler les compléments \textbf{avant} le participe épithète, créant des groupes nominaux compacts que le français doit déployer en phrases complètes.
\begin{center}
\all{der von Goethe im Jahr 1808 in Weimar uraufgeführte „Faust“}
\end{center}
\emph{Traduction :} « le \emph{Faust} créé (joué pour la première fois) par Goethe en 1808 à Weimar ».
En français, il faut une relative passive (\emph{qui a été créé}) plus les compléments. En allemand, tout tient dans le groupe nominal : \textbf{article + complément d'agent (von) + complément de temps + complément de lieu + Participe II + nom}. C'est le \textbf{style nominal} (\all{Nominalstil}) typique de l'allemand écrit (presse, droit, sciences). Au Bac, savoir « déplier » (version) ou « replier » (thème) ces structures est la marque des excellents élèves.
\end{savaistu}

\coursec{Grammaire 2 : le génitif}

Après avoir étudié les participes utilisés comme épithètes, qui permettent de qualifier le nom de manière compacte (\all{der schreibende Dichter}, \all{das geschriebene Gedicht}), nous abordons maintenant un cas fondamental pour lier les noms entre eux : \textbf{le génitif}. C'est le cas de la \cle{possession}, de l'\cle{appartenance} et de la \cle{complémentation nominale} (le « de » français). Dans notre texte sur les arts et la littérature, il apparaît naturellement pour relier l'œuvre à son créateur ou le lieu à ses institutions.

\courssub{Observation dans le texte}

Reprenons trois segments extraits de notre texte support \emph{« Kunst und Literatur im deutschsprachigen Raum »} :

\begin{textebox}[Extraits du texte support]
\begin{enumerate}
    \item \all{Die Werke des Dichters sind in der ganzen Welt bekannt.}
    \item \all{Das berühmte Gemälde des Malers hängt in der Nationalgalerie.}
    \item \all{Die Museen Berlins ziehen jedes Jahr Millionen Besucher an.}
\end{enumerate}
\end{textebox}

\noindent \textbf{Analyse :}
\begin{itemize}
    \item Phrase 1 : \all{des Dichters} (masculin singulier) $\rightarrow$ « du poète ». Le nom \all{der Dichter} prend un \cle{-s}.
    \item Phrase 2 : \all{des Malers} (masculin singulier) $\rightarrow$ « du peintre ». Le nom \all{der Maler} prend un \cle{-s}.
    \item Phrase 3 : \all{Berlins} (nom propre) $\rightarrow$ « de Berlin ». Le nom propre prend un \cle{-s} \textbf{sans apostrophe}.
\end{itemize}
Dans les trois cas, le mot au génitif (\all{des Dichters}, \all{des Malers}, \all{Berlins}) \textbf{suit} le nom qu'il détermine (\all{die Werke}, \all{das Gemälde}, \all{die Museen}). C'est l'ordre inverse du français (« les œuvres \emph{du} poète »).

\begin{propriete}[Fonction et principe de base]
Le \textbf{génitif} (\all{der Genitiv}, 2. Fall, « Wesfall » : \all{wessen?} — « de qui ? ») exprime principalement :
\begin{enumerate}
    \item La \textbf{possession / appartenance} : \all{das Buch des Schülers} (le livre de l'élève).
    \item La \textbf{relation entre deux noms} (complément du nom) : \all{die Hauptstadt Deutschlands} (la capitale de l'Allemagne).
    \item Certaines \textbf{prépositions} (au programme : \all{während}, \all{wegen}, \all{statt}, \all{trotz}, \all{innerhalb}, \all{außerhalb}) : \all{wegen des Regens} (à cause de la pluie), \all{während des Konzerts} (pendant le concert).
\end{enumerate}
\cle{Règle d'or :} en allemand, le possesseur (au génitif) suit la chose possédée.
\end{propriete}

\courssub{Formation du génitif : articles et noms}

La marque principale du génitif se voit sur \textbf{l'article} (et l'adjectif). Les noms masculins et neutres au singulier prennent en plus une terminaison.

\begin{center}
\begin{tabularx}{\linewidth}{|X|X|X|X|}
\hline
\rowcolor{popblueL} \textbf{Cas / Genre} & \textbf{Masculin} & \textbf{Féminin} & \textbf{Neutre} \\
\hline
\rowcolor{popblue!10} \textbf{Nominatif} & \all{der Dichter} & \all{die Dichterin} & \all{das Gedicht} \\
\hline
\rowcolor{popgreen!10} \textbf{Génitif (défini)} & \all{\textbf{des} Dichter\textbf{s}} & \all{\textbf{der} Dichterin} & \all{\textbf{des} Gedicht\textbf{es}} \\
\hline
\rowcolor{poporange!10} \textbf{Génitif (indéfini)} & \all{\textbf{eines} Dichter\textbf{s}} & \all{\textbf{einer} Dichterin} & \all{\textbf{eines} Gedicht\textbf{es}} \\
\hline
\rowcolor{poppurple!10} \textbf{Négatif (kein)} & \all{\textbf{keines} Dichter\textbf{s}} & \all{\textbf{keiner} Dichterin} & \all{\textbf{keines} Gedicht\textbf{es}} \\
\hline
\rowcolor{poppink!10} \textbf{Possessif (mein)} & \all{\textbf{meines} Dichter\textbf{s}} & \all{\textbf{meiner} Dichterin} & \all{\textbf{meines} Gedicht\textbf{es}} \\
\hline
\end{tabularx}
\end{center}

\begin{aretenir}[Terminaisons des noms au génitif singulier]
\begin{itemize}
    \item \textbf{Masculins et neutres :} ajout de \textbf{-(e)s}.
    \begin{itemize}
        \item \textbf{-s} après les polysyllabes, notamment ceux en -el, -er, -en, -ling, -ig : \all{des Vaters, des Malers, des Dichters, des Königs, des Schülers}.
        \item \textbf{-es} pour les monosyllabes et les noms en \cle{-s, -ß, -z, -x, -tz} (euphonie) : \all{des Hauses, des Buches, des Mannes, des Preises, des Sitzes, des Arztes}.
    \end{itemize}
    \item \textbf{Féminins :} \textbf{pas de terminaison} sur le nom : \all{der Dichterin, der Kunst, der Schweiz}.
    \item \textbf{Pluriel :} \textbf{pas de terminaison} au génitif (l'article \all{der} marque seul le cas) : \all{der Dichter, der Dichterinnen, der Bilder}.
\end{itemize}
\end{aretenir}

\begin{popfigure}
\begin{tikzpicture}[
    box/.style={draw=popblue, thick, rounded corners, fill=popblue!5, minimum width=3.2cm, minimum height=1.1cm, align=center},
    arrow/.style={-Stealth, thick, popdark}
]
\node[box, align=center] (possessed){\textbf{Chose possédée}\\ \all{die Werke}};
\node[box, right=3.5cm of possessed, align=center] (possessor){\textbf{Possesseur (GÉNITIF)}\\ \all{des Dichters}};
\draw[arrow] (possessed) -- (possessor) node[midway, above, font=\small, align=center] {ordre allemand :\\ N + Génitif};
\node[draw=poporange, thick, rounded corners, fill=poporange!5, align=center, below=1.6cm of possessed] (fr) {\textbf{Français :} chose + \textbf{de} + possesseur\\ \emph{les œuvres \textbf{du} poète}};
\node[draw=popgreen, thick, rounded corners, fill=popgreen!5, align=center, below=1.6cm of possessor] (de) {\textbf{Allemand :} chose + possesseur au génitif\\ \all{die Werke \textbf{des Dichters}}};
\draw[arrow, poppurple, dashed] (fr) -- (de) node[midway, left, font=\small] {Inversion !};
\end{tikzpicture}
\legende{Ordre des mots : le possesseur au génitif suit la chose possédée (inversion par rapport au français).}
\end{popfigure}

\courssub{Cas particuliers et points de vigilance}

\begin{definition}[Les noms faibles (N-Deklination) au génitif]
Les noms masculins dits « faibles » (êtres vivants, souvent en -e au nominatif : \all{der Junge, der Student, der Mensch, der Nachbar, der Herr}) ajoutent \textbf{-n / -en} à \textbf{tous les cas sauf le nominatif singulier}, donc aussi au génitif.
\begin{center}
\begin{tabular}{|l|l|l|l|}
\hline
\rowcolor{popblueL} Cas & \all{der Student} & \all{der Mensch} & \all{der Herr} \\
\hline
Nominatif & \all{der Student} & \all{der Mensch} & \all{der Herr} \\
\hline
\textbf{Génitif} & \all{des Student\textbf{en}} & \all{des Mensch\textbf{en}} & \all{des Herr\textbf{n}} \\
\hline
\end{tabular}
\end{center}
\emph{Exemple :} \all{Die Abschlussarbeit des Studenten ist hervorragend.} « Le mémoire de l'étudiant est excellent. »
\end{definition}

\begin{definition}[Les noms propres au génitif]
Les noms propres (personnes, villes, pays sans article) prennent un \textbf{-s} \textbf{sans apostrophe} et se placent en général \textbf{avant} la chose possédée.
\end{definition}

\begin{exemplebox}[Noms propres au génitif]
\all{Goethes Faust} « le Faust de Goethe » — et non \all{*Goethe's Faust}.\\
\all{Mozarts Opern} « les opéras de Mozart ».\\
\all{Berlins Museen} « les musées de Berlin ».\\
\all{Deutschlands Hauptstadt} « la capitale de l'Allemagne ».\\
\emph{Cas limite :} un nom propre terminé par -s, -ß, -z, -x prend une apostrophe \textbf{muette} : \all{Hans' Gedicht}, \all{Franz' Bilder} (on ne prononce pas de -s supplémentaire).
\end{exemplebox}

\begin{attention}[Piège majeur : l'apostrophe anglaise]
\cle{Jamais d'apostrophe} au génitif allemand ! L'influence de l'anglais (\emph{Goethe's}) conduit à l'erreur \all{*Goethe's Werk}.
\begin{itemize}
    \item \textbf{Correct :} \all{Goethes Werk}, \all{Annas Buch}, \all{das Buch meines Vaters}.
    \item \textbf{Faux :} \all{*Goethe's Werk}, \all{*das Buch meines Vater's}.
\end{itemize}
\end{attention}

\begin{attention}[L'alternative « von + Datif » : oral contre écrit]
À l'oral familier, on remplace souvent le génitif par \all{von} + \textbf{datif} : \all{das Bild von dem Maler} au lieu de \all{das Bild des Malers}.\\
\cle{Au baccalauréat (écrit et expression soignée), le génitif est obligatoire.} L'emploi de \all{von} + datif dans une rédaction ou un commentaire de texte est une faute de registre. \all{von} reste correct uniquement quand le nom est seul, sans article : \all{ein Bild von großer Schönheit} « un tableau d'une grande beauté ».
\end{attention}

\courssub{Tableau récapitulatif : génitif avec adjectif}

L'adjectif épithète au génitif prend \textbf{toujours la terminaison -en} lorsqu'un article ou un possessif porte déjà la marque du cas (\all{des, der, eines, meiner}...).

\begin{center}
\begin{tabularx}{\linewidth}{|>{\bfseries}l|X|X|X|X|}
\hline
\rowcolor{popblueL} & \textbf{Masculin} & \textbf{Féminin} & \textbf{Neutre} & \textbf{Pluriel} \\
\hline
\textbf{Article défini} & \all{des guten Dichters} & \all{der berühmten Dichterin} & \all{des schönen Gedichts} & \all{der alten Bilder} \\
\hline
\textbf{Article indéfini} & \all{eines guten Dichters} & \all{einer berühmten Dichterin} & \all{eines schönen Gedichts} & --- \\
\hline
\textbf{Négatif} & \all{keines guten Dichters} & \all{keiner berühmten Dichterin} & \all{keines schönen Gedichts} & \all{keiner alten Bilder} \\
\hline
\textbf{Possessif} & \all{meines guten Dichters} & \all{meiner berühmten Dichterin} & \all{meines schönen Gedichts} & \all{meiner alten Bilder} \\
\hline
\end{tabularx}
\end{center}

\courssub{Comparaison systématique français / allemand}

\begin{center}
\begin{tabularx}{\linewidth}{|X|X|X|}
\hline
\rowcolor{popblueL} \textbf{Français} & \textbf{Allemand (génitif)} & \textbf{Remarque structurelle} \\
\hline
le roman \textbf{de l'écrivain} & \all{der Roman \textbf{des Schriftstellers}} & Inversion : N + génitif. \\
les poèmes \textbf{de la poétesse} & \all{die Gedichte \textbf{der Dichterin}} & Féminin : pas de -s sur le nom. \\
la capitale \textbf{de l'Allemagne} & \all{die Hauptstadt \textbf{Deutschlands}} & Nom propre : -s sans apostrophe. \\
les tableaux \textbf{des peintres} & \all{die Bilder \textbf{der Maler}} & Pluriel : pas de -s sur le nom. \\
le livre \textbf{de l'étudiant} & \all{das Buch \textbf{des Studenten}} & Nom faible : -en au génitif. \\
la maison \textbf{du père} & \all{das Haus \textbf{des Vaters}} & Masculin : -s. \\
\hline
\end{tabularx}
\end{center}

\courssub{Méthode : traduire le « de » possessif (thème)}

\begin{methode}[Du français vers l'allemand : « de » = génitif]
Pour traduire correctement une relation de possession (« le X de Y ») :
\begin{enumerate}
    \item \textbf{Identifiez les deux noms :} la \cle{chose possédée} (ex. \emph{roman}) et le \cle{possesseur} (ex. \emph{écrivain}).
    \item \textbf{Déterminez le genre et le nombre du possesseur} et choisissez l'article au génitif : \all{des} (masc./neutre sg.), \all{der} (fém. sg. et pluriel), \all{eines/einer}, \all{meines/meiner}.
    \item \textbf{Appliquez la terminaison du nom :}
    \begin{itemize}
        \item masculin / neutre singulier $\rightarrow$ \textbf{-(e)s} : \all{des Schriftstellers, des Buches} ;
        \item féminin et pluriel $\rightarrow$ \textbf{rien} : \all{der Dichterin, der Maler} ;
        \item nom faible masculin $\rightarrow$ \textbf{-(e)n} : \all{des Studenten, des Nachbarn} ;
        \item nom propre $\rightarrow$ \textbf{-s sans apostrophe} : \all{Goethes, Berlins}.
    \end{itemize}
    \item \textbf{Placez le possesseur après la chose possédée :} \all{der Roman des Schriftstellers}.
    \item \textbf{Mettez l'adjectif éventuel en -en :} \all{der Roman des berühmten Schriftstellers}.
\end{enumerate}
\end{methode}

\courssub{Exemples traduits et phrases modèles}

\begin{exemplebox}[Possession simple]
\all{Die Romane des Schriftstellers sind spannend.} « Les romans de l'écrivain sont passionnants. » (masc. sg. : \all{der Schriftsteller} $\rightarrow$ \all{des Schriftstellers})\\
\all{Die Bilder der Malerin sind farbenfroh.} « Les tableaux de la peintre sont colorés. » (fém. sg. : \all{die Malerin} $\rightarrow$ \all{der Malerin}, sans -s)\\
\all{Die Gedichte der jungen Dichterin erscheinen im Herbst.} « Les poèmes de la jeune poétesse paraissent en automne. »
\end{exemplebox}

\begin{exemplebox}[Noms propres et géographie]
\all{Goethes Geburtstag ist am 28. August.} « L'anniversaire de Goethe est le 28 août. »\\
\all{Der Kölner Dom ist weltberühmt.} « La cathédrale de Cologne est mondialement connue. » (forme adjectivale fréquente pour les villes ; génitif possible : \all{die Baugeschichte des Kölner Doms}).\\
\all{Die Museen Wiens ziehen viele Touristen an.} « Les musées de Vienne attirent beaucoup de touristes. »
\end{exemplebox}

\begin{exemplebox}[Avec adjectif, article indéfini et possessif]
\all{Er liest ein Buch eines berühmten Dichters.} « Il lit un livre d'un poète célèbre. » (indéfini masc. sg. : \all{eines berühmten Dichters})\\
\all{Das Auto meines Vaters ist neu.} « La voiture de mon père est neuve. » (possessif masc. sg. : \all{meines Vaters})\\
\all{Die Ausstellung der modernen Maler beginnt morgen.} « L'exposition des peintres modernes commence demain. » (pluriel : \all{der modernen Maler})
\end{exemplebox}

\courssub{Exercice résolu : mise en application}

\begin{exoresolu}[Traduire avec le génitif (thème)]
\textbf{Énoncé :} Traduisez en allemand en utilisant le génitif.
\begin{enumerate}
    \item La sœur du professeur. (\all{der Lehrer})
    \item Les idées de la jeune fille. (\all{das Mädchen})
    \item Le nouveau livre de Thomas Mann. (\all{das neue Buch})
    \item La décision du gouvernement. (\all{die Regierung})
    \item Les œuvres des grands poètes. (\all{der große Dichter})
    \item La voiture de mon frère. (\all{mein Bruder})
\end{enumerate}
\tcblower
\textbf{Solution détaillée :}
\begin{enumerate}
    \item \all{die Schwester \textbf{des Lehrers}} — \all{der Lehrer}, masc. sg. $\rightarrow$ article \all{des} + nom en -s.
    \item \all{die Ideen \textbf{des Mädchens}} — \all{das Mädchen}, neutre sg. $\rightarrow$ \all{des Mädchens} (-s).
    \item \all{\textbf{Thomas Manns} neues Buch} ou \all{das neue Buch \textbf{Thomas Manns}} — nom propre $\rightarrow$ -s sans apostrophe ; placé avant, il déclenche l'adjectif fort \all{neues}.
    \item \all{die Entscheidung \textbf{der Regierung}} — \all{die Regierung}, fém. sg. $\rightarrow$ article \all{der}, nom inchangé.
    \item \all{die Werke \textbf{der großen Dichter}} — pluriel : article \all{der}, adjectif en -en (\all{großen}), nom sans terminaison (\all{Dichter}).
    \item \all{das Auto \textbf{meines Bruders}} — \all{mein Bruder}, masc. sg. $\rightarrow$ \all{meines} + \all{Bruders}.
\end{enumerate}
\end{exoresolu}

\courssub{Exercices d'entraînement (auto-évaluation)}

\begin{exercice}[Génitif : forme et emploi]
\textbf{A.} Mettez le nom entre parenthèses au génitif (article défini).
\begin{enumerate}
    \item Die Mutter \_\_\_\_\_\_\_\_\_\_\_\_ (\all{der Junge}).
    \item Das Ende \_\_\_\_\_\_\_\_\_\_\_\_ (\all{der Film}).
    \item Die Farbe \_\_\_\_\_\_\_\_\_\_\_\_ (\all{das Auto}).
    \item Der Name \_\_\_\_\_\_\_\_\_\_\_\_ (\all{die Frau}).
    \item Die Freunde \_\_\_\_\_\_\_\_\_\_\_\_ (\all{mein Bruder}).
    \item Die Hauptstadt \_\_\_\_\_\_\_\_\_\_\_\_ (\all{Österreich}).
\end{enumerate}

\textbf{B.} Traduisez en allemand (ordre des mots et terminaisons).
\begin{enumerate}
    \item Le chat de ma sœur.
    \item Les enfants du voisin. (\all{der Nachbar})
    \item La signature de l'auteur. (\all{der Autor})
    \item Les musées de Vienne. (\all{Wien})
    \item Le nouveau roman de l'écrivaine. (\all{die Schriftstellerin})
\end{enumerate}

\textbf{C.} Corrigez l'erreur dans chaque phrase (une erreur par phrase).
\begin{enumerate}
    \item \all{Das ist das Buch meines Vater's.}
    \item \all{Die Werke der Dichters sind bekannt.}
    \item \all{Goethe's Faust ist ein Klassiker.}
    \item \all{Das Haus von dem Lehrer ist groß.} (style écrit !)
    \item \all{Die Entscheidung des Regierung war schnell.}
\end{enumerate}
\end{exercice}

\begin{corrige}[Exercice : le génitif]
\textbf{A.}
\begin{enumerate}
    \item \all{des Jungen} (nom faible : -en)
    \item \all{des Films} (masc. sg. : -s)
    \item \all{des Autos} (neutre sg. : -s)
    \item \all{der Frau} (fém. sg. : inchangé)
    \item \all{meines Bruders} (possessif masc. sg. + -s)
    \item \all{Österreichs} (nom propre : -s sans apostrophe)
\end{enumerate}

\textbf{B.}
\begin{enumerate}
    \item \all{die Katze meiner Schwester}
    \item \all{die Kinder des Nachbarn} (\all{der Nachbar} est un nom faible $\rightarrow$ -n)
    \item \all{die Unterschrift des Autors}
    \item \all{die Museen Wiens} (toponyme $\rightarrow$ \all{Wiens})
    \item \all{der neue Roman der Schriftstellerin}
\end{enumerate}

\textbf{C.}
\begin{enumerate}
    \item \all{*Vater's} $\rightarrow$ \all{Vaters} (jamais d'apostrophe).
    \item \all{*der Dichters} $\rightarrow$ \all{der Dichter} (pluriel : pas de -s sur le nom).
    \item \all{*Goethe's} $\rightarrow$ \all{Goethes}.
    \item \all{*von dem Lehrer} $\rightarrow$ \all{des Lehrers} (génitif exigé à l'écrit).
    \item \all{*des Regierung} $\rightarrow$ \all{der Regierung} (féminin : article \all{der}, nom inchangé).
\end{enumerate}
\end{corrige}

\begin{savaistu}[Le génitif dans la culture germanophone]
Le génitif est souvent considéré comme le « cas noble » de la langue allemande. Dans la presse de qualité (\all{Die Zeit}, \all{Frankfurter Allgemeine Zeitung}, \all{Süddeutsche Zeitung}) et dans la littérature, il est omniprésent. À l'oral courant et dans plusieurs dialectes (bavarois, alémanique), on lui préfère volontiers \all{von} + datif, voire la tournure populaire \all{dem Vater sein Buch} (« au père son livre »), à proscrire à l'écrit. \textbf{Pour le baccalauréat et l'écrit académique, maîtrisez le génitif parfaitement :} c'est un marqueur de niveau élevé. Au Cameroun, dans la rédaction de rapports formels en allemand, l'usage correct de tournures comme \all{die Entscheidung des Direktors} « la décision du directeur » témoigne d'une compétence linguistique remarquable.
\end{savaistu}

\coursec{Conjugaison : le prétérit des verbes forts}

Après avoir étudié le génitif, qui marque l'appartenance et la relation entre les noms, nous abordons maintenant le temps verbal qui structure le récit écrit en allemand : le \cle{Präteritum}. Alors que le français utilise le passé simple et l'imparfait pour narrer au passé, l'allemand \textbf{écrit} (presse, littérature, contes, rapports) utilise presque exclusivement le prétérit. Maîtriser les \textbf{verbes forts} (verbes irréguliers à changement de voyelle) est donc indispensable pour lire et écrire un récit cohérent.

\courssub{Emploi du prétérit : le temps du récit écrit}

En allemand, le choix du temps du passé dépend du \cle{registre de langue} (écrit vs oral) et non de l'aspect (accompli / inaccompli) comme en français.

\begin{propriete}[Emploi du Präteritum]
Le \textbf{Präteritum} (preterite) est le temps du récit \textbf{écrit} (littérature, journalisme, contes, histoire, protocoles).
\begin{itemize}
    \item \textbf{Récit au passé :} \all{Goethe schrieb „Faust“.} « Goethe écrivit « Faust ». »
    \item \textbf{Description d'états passés :} \all{Es war ein kalter Tag.} « C'était un jour froid. »
    \item \textbf{Presse / Rapport :} \all{Der Maler kam nach Berlin.} « Le peintre vint à Berlin. »
\end{itemize}
\emph{Attention :} À l'oral (conversation, radio, TV), on utilise le \textbf{Perfekt} (\all{er hat geschrieben}). Ne traduisez pas le prétérit allemand par le passé composé français dans un texte littéraire !
\end{propriete}

\begin{attention}[Piège francophone : Prétérit vs Perfekt]
\begin{itemize}
    \item \textbf{Écrit (récit) :} \all{Der Schriftsteller las sein Gedicht. Der Applaus war groß.} (Prétérit)
    \item \textbf{Oral (raconter) :} \all{Der Schriftsteller hat sein Gedicht gelesen. Der Applaus war groß.} (Perfekt + Prétérit pour \all{sein/war})
    \item \textbf{Français (les deux) :} « L'écrivain lut son poème. L'applaudissement fut grand. » OU « L'écrivain a lu son poème... »
\end{itemize}
N'utilisez \textbf{JAMAIS} le Perfekt dans une rédaction littéraire, un résumé de texte ou un commentaire d'œuvre au lycée.
\end{attention}

\courssub{Formation des verbes forts : radical modifié + désinences}

Contrairement aux verbes faibles (réguliers : \all{machen $\to$ machte}) et aux verbes mixtes, les \textbf{verbes forts} subissent une \cle{modification de la voyelle radicale} (Ablaut) au prétérit. Le radical change, mais les désinences sont régulières (sauf 1ère et 3ème personne du singulier sans \cle{-e}).

\begin{propriete}[Désinences du Prétérit (verbes forts)]
Pas de \cle{-e} à la 1ère et 3ème personne du singulier ! Le radical modifié porte l'accent.
\begin{center}
\begin{tabular}{|l|l|l|l|}
\hline
\rowcolor{popblueL} \textbf{Personne} & \textbf{Désinence} & \textbf{Exemple : \all{schreiben}} & \textbf{Traduction} \\
\hline
ich & -- & \all{ich schrieb} & j'écrivis \\
du & \cle{-st} & \all{du schriebst} & tu écrivis \\
er/sie/es & -- & \all{er schrieb} & il/elle écrivit \\
wir & \cle{-en} & \all{wir schrieben} & nous écrivîmes \\
ihr & \cle{-t} & \all{ihr schriebt} & vous écrivîtes \\
sie/Sie & \cle{-en} & \all{sie schrieben} & ils/elles écrivirent \\
\hline
\end{tabular}
\end{center}
\emph{Précision :} Les verbes en \cle{-eln} ou \cle{-ern} (ex. \all{wandern, kosten, handeln}) sont \textbf{toujours faibles} (réguliers) : ils gardent leur \cle{-e-} de liaison au prétérit (\all{ich wanderte, ich kostete}). Il n'existe pas de verbes forts en \cle{-eln/-ern}.
\end{propriete}

\courssub{Les verbes forts du texte et du thème : tableaux et séries vocaliques}

Pour ne pas apprendre chaque verbe isolément, regroupez-les par \cle{séries de changement vocalique}. C'est la méthode la plus efficace pour la mémoire à long terme.

\begin{textebox}[Verbes forts du thème « Kunst und Literatur »]
Voici les verbes forts (et un faible piège) rencontrés dans notre texte support et le vocabulaire du thème. Apprenez les \textbf{triplets} : Infinitif -- Prétérit -- Participe II.
\end{textebox}

\begin{popfigure}
\legende{Les triplets principaux du thème (Infinitiv -- Präteritum -- Partizip II)}
\begin{tikzpicture}
\matrix (verbtable) [
  matrix of nodes,
  nodes={font=\small, inner sep=3pt, minimum height=0.65cm, anchor=center, draw=popline, fill=popblue!10},
  column 1/.style={nodes={minimum width=3.5cm}},
  column 2/.style={nodes={minimum width=3.5cm}},
  column 3/.style={nodes={minimum width=3.5cm}},
  row 1/.style={nodes={fill=popblue, text=white, font=\bfseries\small, minimum height=0.7cm}},
  row 4/.style={nodes={fill=poporange!15}}, % malen (faible)
  row 13/.style={nodes={fill=poppurple!15}}, % denken/bringen (mixtes)
  row 14/.style={nodes={fill=poppurple!15}}, % kennen/wissen (mixtes)
  column sep=-\pgflinewidth, row sep=-\pgflinewidth
] {
\all{Infinitiv} & \all{Präteritum} & \all{Partizip II} \\
\all{schreiben} & \all{schrieb} & \all{geschrieben} \\
\all{lesen} & \all{las} & \all{gelesen} \\
\all{malen} \emph{(faible!)} & \all{malte} & \all{gemalt} \\
\all{kommen} & \all{kam} & \all{gekommen} \\
\all{sehen} & \all{sah} & \all{gesehen} \\
\all{gehen} & \all{ging} & \all{gegangen} \\
\all{finden} & \all{fand} & \all{gefunden} \\
\all{bleiben} & \all{blieb} & \all{geblieben} \\
\all{singen} & \all{sang} & \all{gesungen} \\
\all{trinken} & \all{trank} & \all{getrunken} \\
\all{beginnen} & \all{begann} & \all{begonnen} \\
\all{sterben} & \all{starb} & \all{gestorben} \\
\all{denken / bringen} & \all{dachte / brachte} & \all{gedacht / gebracht} \\
\all{kennen / wissen} & \all{kannte / wusste} & \all{gekannt / gewusst} \\
};
\end{tikzpicture}
\end{popfigure}

\begin{attention}[Le piège « malen » vs « mahlen »]
\all{malen} (peindre) est un verbe \textbf{FAIBLE} (régulier) !
\begin{itemize}
    \item Prétérit : \all{ich malte} (radical inchangé + \cle{-te})
    \item Participe : \all{gemalt}
    \item \textbf{Erreur fréquente :} *\all{er mahlte} (confusion orthographique avec \all{mahlen} = moudre).
\end{itemize}
\emph{Note :} \all{mahlen} (moudre) est aussi un verbe faible en allemand moderne (\all{mahlte, gemahlen}), mais l'orthographe \cle{ah} vs \cle{a} distingue les deux verbes. De même : \all{behandeln, gestalten, kopieren, organisieren} sont faibles.
\end{attention}

\begin{propriete}[Les 4 grandes séries de changement vocalique (Ablautreihen)]
Classez les verbes forts par série pour les mémoriser. La voyelle du radical change selon un schéma prévisible.

\begin{center}
\begin{tabularx}{\linewidth}{|X|X|X|X|}
\hline
\rowcolor{popblueL} \textbf{Série} & \textbf{Modèle} & \textbf{Verbes du thème / fréquents} & \textbf{Participe II (souvent -en)} \\
\hline
\cle{ei $\to$ ie} & \all{schreiben $\to$ schrieb} & \all{schreiben, bleiben, treiben, beißen, scheinen} & \all{geschrieben, geblieben} \\
\hline
\cle{e $\to$ a / ä} & \all{geben $\to$ gab} & \all{lesen $\to$ las, sehen $\to$ sah, geben, nehmen, treffen, essen $\to$ aß, helfen $\to$ half} & \all{gelesen, gesehen, gegeben} \\
\hline
\cle{i $\to$ a} & \all{trinken $\to$ trank} & \all{singen $\to$ sang, trinken, beginnen $\to$ begann, finden $\to$ fand, binden, gewinnen, schwimmen} & \cle{gesungen, getrunken, begonnen, gefunden} \\
\hline
\cle{ie $\to$ o} & \all{bieten $\to$ bot} & \all{ziehen $\to$ zog, fliegen $\to$ flog, bieten, frieren $\to$ fror, verlieren $\to$ verlor} & \cle{gezogen, geflogen, geboten} \\
\hline
\end{tabularx}
\end{center}
\emph{Note :} La série \cle{e $\to$ a} comporte souvent un \cle{Umlaut} au pluriel du présent (\all{er liest, wir lesen}) mais pas au prétérit (\all{er las, wir lasen}). La série \cle{i $\to$ a} donne souvent un participe en \cle{-u-} (\all{getrunken, gesungen}).
\end{propriete}

\begin{popfigure}
\legende{Carte mentale des séries de changement vocalique (Ablautreihen) pour verbes forts}
\begin{center}\tcbox[colback=popink,colframe=popink,coltext=white,arc=9pt,boxsep=4pt,left=6pt,right=6pt]{\bfseries\small Verbes forts Präteritum}\end{center}
\vspace{-2pt}
\begin{tcbraster}[raster columns=2,raster equal height=rows,raster column skip=6pt,raster row skip=6pt,enhanced,blankest]
\begin{tcolorbox}[enhanced,colback=white,colframe=popgreen,boxrule=0.9pt,arc=3pt,title={ei $\to$ ie},fonttitle=\bfseries\scriptsize,coltitle=white,colbacktitle=popgreen,left=4pt,right=4pt,top=2pt,bottom=2pt,boxsep=1pt,toptitle=1.5pt,bottomtitle=1.5pt,halign=left]
\begin{itemize}[nosep,leftmargin=*,itemsep=1pt,font=\scriptsize]
\item \all{schreiben $\to$ schrieb}
\item \all{bleiben $\to$ blieb}
\item \all{beißen $\to$ biss}
\end{itemize}
\end{tcolorbox}
\begin{tcolorbox}[enhanced,colback=white,colframe=poporange,boxrule=0.9pt,arc=3pt,title={e $\to$ a/ä},fonttitle=\bfseries\scriptsize,coltitle=white,colbacktitle=poporange,left=4pt,right=4pt,top=2pt,bottom=2pt,boxsep=1pt,toptitle=1.5pt,bottomtitle=1.5pt,halign=left]
\begin{itemize}[nosep,leftmargin=*,itemsep=1pt,font=\scriptsize]
\item \all{lesen $\to$ las}
\item \all{sehen $\to$ sah}
\item \all{geben $\to$ gab}
\item \all{essen $\to$ aß}
\end{itemize}
\end{tcolorbox}
\begin{tcolorbox}[enhanced,colback=white,colframe=poppurple,boxrule=0.9pt,arc=3pt,title={i $\to$ a},fonttitle=\bfseries\scriptsize,coltitle=white,colbacktitle=poppurple,left=4pt,right=4pt,top=2pt,bottom=2pt,boxsep=1pt,toptitle=1.5pt,bottomtitle=1.5pt,halign=left]
\begin{itemize}[nosep,leftmargin=*,itemsep=1pt,font=\scriptsize]
\item \all{singen $\to$ sang}
\item \all{trinken $\to$ trank}
\item \all{beginnen $\to$ begann}
\item \all{finden $\to$ fand}
\end{itemize}
\end{tcolorbox}
\begin{tcolorbox}[enhanced,colback=white,colframe=poppink,boxrule=0.9pt,arc=3pt,title={ie $\to$ o},fonttitle=\bfseries\scriptsize,coltitle=white,colbacktitle=poppink,left=4pt,right=4pt,top=2pt,bottom=2pt,boxsep=1pt,toptitle=1.5pt,bottomtitle=1.5pt,halign=left]
\begin{itemize}[nosep,leftmargin=*,itemsep=1pt,font=\scriptsize]
\item \all{ziehen $\to$ zog}
\item \all{fliegen $\to$ flog}
\item \all{bieten $\to$ bot}
\end{itemize}
\end{tcolorbox}
\end{tcbraster}

\end{popfigure}

\courssub{Les verbes mixtes : changement de voyelle + désinence faible}

Les \cle{verbes mixtes} (gemischte Verben) combinent les deux systèmes : ils changent de voyelle radicale \textbf{ET} prennent les désinences faibles (\cle{-te}, \cle{-test}, \cle{-ten}...). Ils sont peu nombreux mais très fréquents.

\begin{propriete}[Conjugaison des verbes mixtes au Prétérit]
Radical modifié + désinences en \cle{-t-} (comme les faibles).
\begin{center}
\begin{tabular}{|l|l|l|l|l|}
\hline
\rowcolor{poppurpleL} \textbf{Infinitif} & \textbf{Présent (er)} & \textbf{Präteritum (ich/er)} & \textbf{Participe II} & \textbf{Sens} \\
\hline
\all{denken} & \all{denkt} & \all{dachte} & \all{gedacht} & penser \\
\all{bringen} & \all{bringt} & \all{brachte} & \all{gebracht} & apporter \\
\all{kennen} & \all{kennt} & \all{kannte} & \all{gekannt} & connaître (qn/qch) \\
\all{wissen} & \all{weiß} & \all{wusste} & \all{gewusst} & savoir (fait) \\
\all{brennen} & \all{brennt} & \all{brannte} & \all{gebrannt} & brûler \\
\all{nennen} & \all{nennt} & \all{nannte} & \all{genannt} & nommer \\
\all{rennen} & \all{rennt} & \all{rannte} & \all{gerannt} & courir \\
\hline
\end{tabular}
\end{center}
\emph{Astuce :} La plupart ont \cle{a} au prétérit (\all{dachte, brachte, kannte, wusste, brannte, nannte, rannte}).
\end{propriete}

\courssub{Conjugaison complète modèle : \all{schreiben} (ei $\to$ ie)}

\begin{exemplebox}[Modèle complet : \all{schreiben} -- écrire]
\begin{center}
\begin{tabular}{|l|l|l|}
\hline
\rowcolor{popblueL} \textbf{Personne} & \textbf{Forme} & \textbf{Traduction} \\
\hline
ich & \all{schrieb} & j'écrivis \\
du & \all{schriebst} & tu écrivis \\
er / sie / es & \all{schrieb} & il/elle/on écrivit \\
wir & \all{schrieben} & nous écrivîmes \\
ihr & \all{schriebt} & vous écrivîtes \\
sie / Sie & \all{schrieben} & ils/elles/Vous écrivîtes \\
\hline
\end{tabular}
\end{center}
\textbf{Exemples contextuels (thème Arts/Littérature) :}
\begin{itemize}
    \item \all{Der Dichter \textbf{schrieb} seinen Roman in Berlin.} « Le poète écrivit son roman à Berlin. »
    \item \all{Wir \textbf{schrieben} eine Analyse über das Gedicht.} « Nous écrivîmes une analyse sur le poème. »
    \item \all{Warum \textbf{schriebst} du nicht mehr?} « Pourquoi n'écrivais-tu plus ? » (narratif)
\end{itemize}
\end{exemplebox}

\courssub{Le triplet infinitif – prétérit – participe passé : clé de la mémorisation}

\begin{methode}[Comment apprendre les verbes forts efficacement]
\begin{enumerate}
    \item \textbf{Apprenez par triplets :} \all{Infinitiv -- Präteritum -- Partizip II}. Ex. : \all{lesen -- las -- gelesen}. Ne séparez jamais les trois formes.
    \item \textbf{Groupez par série vocale :} Apprenez tous les \cle{ei $\to$ ie} ensemble, puis tous les \cle{e $\to$ a}, etc. (Voir carte mentale ci-dessus).
    \item \textbf{Associez au thème :} Créez une phrase-type par verbe du thème « Arts ». Ex. : \all{Der Maler malte (faible!) ein Bild. Der Kritiker sah es. Er fand es gut.}
    \item \textbf{Entraînez-vous à l'oral :} Dites les triplets à voix haute : « lesen – las – gelesen ». Le rythme aide la mémoire auditive.
    \item \textbf{Attention aux verbes mixtes :} Liste à part (\all{denken, bringen, kennen, wissen, brennen, nennen, rennen}).
\end{enumerate}
\end{methode}

\begin{aretenir}[Récapitulatif : Les 3 formes principales (Hauptformen)]
Pour \textbf{chaque} verbe fort ou mixte, vous devez connaître par cœur :
\begin{center}
\begin{tabular}{|c|c|c|}
\hline
\rowcolor{popgoldL} \textbf{Infinitiv} & \textbf{Präteritum (ich/er)} & \textbf{Partizip II} \\
\hline
\all{schreiben} & \all{schrieb} & \all{geschrieben} \\
\all{lesen} & \all{las} & \all{gelesen} \\
\all{kommen} & \all{kam} & \all{gekommen} \\
\all{sehen} & \all{sah} & \all{gesehen} \\
\all{gehen} & \all{ging} & \all{gegangen} \\
\all{finden} & \all{fand} & \all{gefunden} \\
\all{bleiben} & \all{blieb} & \all{geblieben} \\
\all{singen} & \all{sang} & \all{gesungen} \\
\all{trinken} & \all{trank} & \all{getrunken} \\
\all{beginnen} & \all{begann} & \all{begonnen} \\
\all{sterben} & \all{starb} & \all{gestorben} \\
\all{entstehen} & \all{entstand} & \all{entstanden} \\
\all{fließen} & \all{floss} & \all{geflossen} \\
\all{schlafen} & \all{schlief} & \all{geschlafen} \\
\all{verstehen} & \all{verstand} & \all{verstanden} \\
\all{nehmen} & \all{nahm} & \all{genommen} \\
\all{denken} & \all{dachte} & \all{gedacht} \\
\all{bringen} & \all{brachte} & \all{gebracht} \\
\all{kennen} & \all{kannte} & \all{gekannt} \\
\all{wissen} & \all{wusste} & \all{gewusst} \\
\hline
\end{tabular}
\end{center}
\cle{malen} (faible) : \all{malen -- malte -- gemalt}.
\end{aretenir}

\courssub{Entraînement : Texte authentique et application}

Lisez le texte suivant, extrait d'une critique littéraire fictive parue dans la \all{Frankfurter Allgemeinen Zeitung}. Repérez tous les verbes au prétérit (forts, mixtes, faibles) et identifiez leur infinitif.

\begin{textebox}[Critique : „Die Stimme der Stille“ — Romanneuerscheinung]
\textbf{Der junge Schriftsteller Jonas K. \textbf{schrieb} seinen zweiten Roman in einer kleinen Hütte im Schwarzwald. Er \textbf{blieb} dort drei Monate lang, fern vom Lärm der Großstadt. Jeden Morgen \textbf{ging} er früh in den Wald und \textbf{sah} das Licht zwischen den Bäumen tanzen. Diese Stille \textbf{inspirierte} ihn zutiefst.}
\textbf{Er \textbf{las} viel: Goethe, Hesse, aber auch afrikanische Autoren wie Mongo Beti. Ihre Worte \textbf{klangen} in seinem Kopf nach. Eines Abends \textbf{fand} er ein altes Tagebuch seiner Großmutter. Er \textbf{öffnete} es und \textbf{las} von ihrer Jugend in Kamerun. Die Geschichten \textbf{handelten} von Kakaoplantagen, von \textbf{gesungenen} Liedern am Abend und von einem Fluss, der nie \textbf{schlief}.}
\textbf{Jonas \textbf{verstand} plötzlich: Literatur \textbf{bedeutete} nicht nur Schönheit, sondern Erinnerung. Er \textbf{nahm} den Stift und \textbf{schrieb} die erste Zeile: „Die Vergangenheit ist ein Land, das wir nie verlassen.“ Der Roman \textbf{entstand} in wenigen Wochen. Die Kritik \textbf{lobte} das Werk später als „Meisterwerk der leisen Töne“.}
\end{textebox}

\begin{definition}[Glossaire du texte (mots difficiles)]
\begin{itemize}
    \item \all{die Hütte, die Hütten} : la cabane, le chalet
    \item \all{der Schwarzwald} : la Forêt-Noire (Allemagne)
    \item \all{zutiefst} : profondément
    \item \all{nachklingen} (faible, \all{klang nach}) : résonner, résonner encore
    \item \all{das Tagebuch, die Tagebücher} : le journal intime
    \item \all{die Kakaoplantage, -n} : la plantation de cacao
    \item \all{der Stift, die Stifte} : le stylo, le crayon
    \item \all{entstehen} (fort, \all{entstand, entstanden}) : naître, voir le jour, être créé
    \item \all{das Meisterwerk, -e} : le chef-d'œuvre
    \item \all{der Ton, die Töne} : le ton, la note (musicale/figurée)
\end{itemize}
\end{definition}

\begin{exoresolu}[Analyse verbale du texte]
\textbf{Consigne :} Relevez 10 verbes au prétérit dans le texte, donnez leur infinitif et leur série (forte, mixte, faible).

\textbf{Solution détaillée :}
\begin{center}
\begin{tabularx}{\linewidth}{|X|X|X|X|}
\hline
\rowcolor{popblueL} \textbf{Forme (Texte)} & \textbf{Infinitif} & \textbf{Type} & \textbf{Série / Remarque} \\
\hline
\all{schrieb} & \all{schreiben} & Fort & ei $\to$ ie \\
\all{blieb} & \all{bleiben} & Fort & ei $\to$ ie \\
\all{ging} & \all{gehen} & Fort & Irrégulier (g$\to$ng) \\
\all{sah} & \all{sehen} & Fort & e $\to$ a \\
\all{inspirierte} & \all{inspirieren} & Faible & -te (verbe en -ieren) \\
\all{las} (x2) & \all{lesen} & Fort & e $\to$ a \\
\all{klangen} & \all{klingen} & Fort & i $\to$ a (Part. \all{geklungen}) \\
\all{fand} & \all{finden} & Fort & i $\to$ a \\
\all{öffnete} & \all{öffnen} & Faible & -te \\
\all{handelten} & \all{handeln} & Faible & -te \\
\all{gesungenen} & \all{singen} & Part. II (adj.) & i $\to$ u (Part. \all{gesungen}) \\
\all{schlief} & \all{schlafen} & Fort & ä $\to$ ie (Présent \all{schläft}) \\
\all{verstand} & \all{verstehen} & Fort & e $\to$ a (Part. \all{verstanden}) \\
\all{bedeutete} & \all{bedeuten} & Faible & -te \\
\all{nahm} & \all{nehmen} & Fort & e $\to$ a (Part. \all{genommen}) \\
\all{entstand} & \all{entstehen} & Fort & i $\to$ a \\
\all{lobte} & \all{loben} & Faible & -te \\
\hline
\end{tabularx}
\end{center}
\emph{Observation :} Le texte narratif utilise quasi exclusivement le Prétérit. Les verbes forts dominent pour les actions principales (\all{schrieb, ging, sah, fand, verstand, nahm, entstand}), les verbes faibles pour les actions secondaires ou verbes récents (\all{inspirierte, öffnete, handelten, bedeutete, lobte}).
\end{exoresolu}

\courssub{Exercices d'application}

\begin{exercice}[Conjugaison au Prétérit : Mettez les verbes entre parenthèses au Prétérit]
Conjuguez les verbes au prétérit. Attention : verbes forts, mixtes et faibles sont mélangés. Pensez aux séries vocaliques !

\begin{enumerate}
    \item Der Maler Max Liebermann \_\_\_\_\_\_\_\_\_\_ (malen) viele Impressionen am Wannsee.
    \item Die Schriftstellerin \_\_\_\_\_\_\_\_\_\_ (schreiben) ihren ersten Roman mit 25 Jahren.
    \item Wir \_\_\_\_\_\_\_\_\_\_ (lesen) das Gedicht „Erlkönig“ in der Schule.
    \item Der Kritiker \_\_\_\_\_\_\_\_\_\_ (finden) das Buch sehr originell.
    \item Ich \_\_\_\_\_\_\_\_\_\_ (denken) oft an meine Kindheit in Yaoundé.
    \item Die Sänger \_\_\_\_\_\_\_\_\_\_ (singen) ein altes Volkslied.
    \item Er \_\_\_\_\_\_\_\_\_\_ (bringen) mir ein Buch von Goethe mit.
    \item Das Festival \_\_\_\_\_\_\_\_\_\_ (beginnen) pünktlich um 20 Uhr.
    \item Wer \_\_\_\_\_\_\_\_\_\_ (wissen) die Antwort auf diese Frage?
    \item Der Fluss \_\_\_\_\_\_\_\_\_\_ (fließen) ruhig durch das Tal. (fort: fließen -- floss -- geflossen)
\end{enumerate}
\end{exercice}

\begin{corrige}[Exercice : Conjugaison au Prétérit]
\begin{enumerate}
    \item \all{malte} (faible : malen -- malte -- gemalt) — \emph{Piège : pas *mahlte !}
    \item \all{schrieb} (fort, ei$\to$ie : schreiben -- schrieb -- geschrieben)
    \item \all{lasen} (fort, e$\to$a : lesen -- las -- gelesen)
    \item \all{fand} (fort, i$\to$a : finden -- fand -- gefunden)
    \item \all{dachte} (mixte : denken -- dachte -- gedacht)
    \item \all{sangen} (fort, i$\to$a : singen -- sang -- gesungen)
    \item \all{brachte} (mixte : bringen -- brachte -- gebracht)
    \item \all{begann} (fort, i$\to$a : beginnen -- begann -- begonnen)
    \item \all{wusste} (mixte : wissen -- wusste -- gewusst) — \emph{Présent : er weiß}
    \item \all{floss} (fort, i$\to$o/ß : fließen -- floss -- geflossen) — \emph{Attention : ß après voyelle longue !}
\end{enumerate}
\end{corrige}

\begin{exercice}[Production écrite guidée : Mini-récit]
Écrivez un court paragraphe (6-8 phrases) en allemand au Prétérit pour raconter la visite d'un musée ou d'une exposition.
\textbf{Contraintes :}
\begin{itemize}
    \item Utilisez au moins \textbf{5 verbes forts} différents (parmi : \all{gehen, sehen, finden, bleiben, kommen, lesen, schreiben, singen, trinken, beginnen, sterben, fließen}).
    \item Utilisez au moins \textbf{2 verbes mixtes} (\all{denken, bringen, kennen, wissen}).
    \item Utilisez \textbf{3 verbes faibles} du thème (\all{malen, besuchen, bewundern, erklären, organisieren}).
    \item Contexte : \all{Letzten Sonntag} (Dimanche dernier) ou \all{Vorgestern} (Avant-hier).
\end{itemize}
\textbf{Amorce :} \all{Letzten Sonntag besuchte ich das Nationalmuseum in Yaoundé...}
\end{exercice}

\begin{savaistu}[Le Prétérit dans la culture germanophone : Le « Märchen »]
Les contes de fées (\all{Märchen}) des Frères Grimm sont le modèle parfait du Prétérit narratif. L'ouverture type : \all{„Es war einmal...“} (Il était une fois...). En allemand, le conte \textbf{ne se raconte jamais au Perfekt} à l'écrit. Si vous lisez \all{„Rotkäppchen ging in den Wald“}, c'est du Prétérit pur. Au cinéma ou à l'oral, on dira : \all{„Rotkäppchen ist in den Wald gegangen“}. Cette distinction stricte Écrit/Oral est une marque de compétence C1/C2.
\end{savaistu}

\coursec{Méthode du commentaire, commentaire modèle et production guidée}

Après avoir maîtrisé le prétérit des verbes forts, les participes en épithète et le génitif, vous disposez maintenant de tous les outils grammaticaux pour analyser et produire un texte structuré sur les arts et la littérature. Cette dernière section vous guide pas à pas dans la rédaction d'un commentaire de texte — exercice incontournable du Baccalauréat — puis vous propose des activités d'expression écrite et orale pour réinvestir le lexique et la grammaire vus dans cette leçon.

\courssub{La méthode du commentaire de texte en 4 étapes}

\begin{methode}[Réussir le commentaire de texte au Bac]
\textbf{Étape 1 — Présenter le texte} : Auteur (si connu), source, date, type de texte (article, extrait de roman, critique, interview), thème général. \\
\textbf{Étape 2 — Dégager l'idée générale} : En une phrase, résumer le sujet principal et la thèse de l'auteur. \\
\textbf{Étape 3 — Analyser les idées principales et le lexique} : Identifier 2 à 3 parties logiques, citer des exemples précis (mots, expressions, structures grammaticales : participes, génitifs, prétérit), expliquer leur fonction. \\
\textbf{Étape 4 — Donner son avis argumenté} : Exprimer une opinion personnelle, la justifier par des arguments et des exemples (du texte ou de sa culture personnelle), conclure.
\end{methode}

\begin{exemplebox}[Expressions utiles pour le commentaire]
\begin{itemize}
    \item \all{Der Text handelt von der kulturellen Vielfalt im deutschsprachigen Raum.} « Le texte traite de la diversité culturelle dans l'espace germanophone. »
    \item \all{Der Autor beschreibt die Entwicklung der modernen Kunst.} « L'auteur décrit le développement de l'art moderne. »
    \item \all{Im ersten Teil geht es um die Rolle der Literatur.} « Dans la première partie, il s'agit du rôle de la littérature. »
    \item \all{Ein wichtiges Beispiel ist das Gemälde „Der Wanderer über dem Nebelmeer“.} « Un exemple important est le tableau "Le Voyageur au-dessus de la mer de brouillard". »
    \item \all{Meiner Meinung nach ist Kunst ein Spiegel der Gesellschaft.} « À mon avis, l'art est un miroir de la société. »
    \item \all{Ich finde das Gedicht beeindruckend, weil es tiefe Gefühle ausdrückt.} « Je trouve le poème impressionnant, car il exprime des sentiments profonds. »
    \item \all{Zusammenfassend lässt sich sagen, dass...} « En résumé, on peut dire que... »
\end{itemize}
\end{exemplebox}

\begin{popfigure}
\begin{tikzpicture}[
    node distance=1.2cm and 2.5cm,
    box/.style={rectangle, draw=popblue, thick, fill=popblueL, rounded corners=6pt, minimum width=3.2cm, minimum height=1.6cm, align=center, font=\small},
    arrow/.style={->, >=Stealth, thick, popdark}
]
\node[box, align=center] (e1){\textbf{Étape 1}\\\textbf{Présenter}\\\textit{Auteur, source, thème}};
\node[box, right=of e1, align=center] (e2){\textbf{Étape 2}\\\textbf{Idée générale}\\\textit{Sujet + thèse en 1 phrase}};
\node[box, right=of e2, align=center] (e3){\textbf{Étape 3}\\\textbf{Analyser}\\\textit{Parties, citations, grammaire}};
\node[box, right=of e3, align=center] (e4){\textbf{Étape 4}\\\textbf{Avis personnel}\\\textit{Opinion + arguments}};
\draw[arrow] (e1.east) -- (e2.west);
\draw[arrow] (e2.east) -- (e3.west);
\draw[arrow] (e3.east) -- (e4.west);
\end{tikzpicture}
\legende{Organigramme des 4 étapes du commentaire de texte}
\end{popfigure}

\courssub{Texte support de référence}

Le texte suivant servira de base au commentaire modèle, au résumé guidé et aux exercices d'application. Il réutilise le vocabulaire de la section 2, les participes épithètes (section 3), le génitif (section 4) et le prétérit (section 5).

\begin{textebox}[Kunst und Literatur im deutschsprachigen Raum — Texte original pour la leçon AL22]
Die kulturelle Landschaft Deutschlands, Österreichs und der Schweiz ist reich und vielfältig. Berühmte Maler wie Albrecht Dürer, Caspar David Friedrich oder Paula Modersohn-Becker prägten die europäische Kunstgeschichte nachhaltig. Friedrichs Gemälde „Der Wanderer über dem Nebelmeer“ zeigt einen einsamen Menschen, der die gewaltige Natur betrachtet. Das Bild fasziniert durch seine mystische Atmosphäre und die feine Malerei.

In der Literatur ragen Johann Wolfgang von Goethe, Friedrich Schiller und Thomas Mann heraus. Goethes „Faust“ und Schillers „Die Räuber“ gehören zum Kanon der Weltliteratur. Auch die Lyrik spielt eine große Rolle: Rainer Maria Rilkes Gedichte, voller tiefer Symbolik, berühren Leser bis heute. In der Schweiz schrieb Friedrich Dürrenmatt dramatische Werke, die moralische Fragen stellen. Österreichs Literaturnobelpreisträgerin Elfriede Jelinek kritisiert in ihren Romanen gesellschaftliche Missstände.

Heute blüht die zeitgenössische Szene: junge Schriftstellerinnen wie Saša Stanišić oder Julia Franck erzählen von Migration und Identität. In Kamerun ebenfalls entstehen bemerkenswerte Werke: die Romane von Patrice Nganang oder die Gedichte von Bate Besong zeigen die afrikanische Perspektive. Kunst und Literatur verbinden Menschen über Grenzen hinweg — sie sind der Spiegel der Seele einer Nation.
\end{textebox}

\begin{center}
\begin{tabularx}{\linewidth}{|X|X|}\hline
\rowcolor{popblueL} \textbf{Deutsch} & \textbf{Français} \\ \hline
die kulturelle Landschaft & le paysage culturel \\ \hline
vielfältig & varié, divers \\ \hline
prägen (Prät. prägte) & marquer, façonner \\ \hline
nachhaltig & durablement \\ \hline
der Wanderer & le voyageur, le promeneur \\ \hline
das Nebelmeer & la mer de brouillard \\ \hline
betrachten & contempler, observer \\ \hline
mystisch & mystique \\ \hline
die Atmosphäre & l'atmosphère \\ \hline
die Malerei & la peinture (art/technique) \\ \hline
ragen (Prät. ragte) & dépasser, ressortir \\ \hline
hervorragen & se distinguer, émerger \\ \hline
der Kanon & le canon (œuvres de référence) \\ \hline
die Weltliteratur & la littérature mondiale \\ \hline
die Lyrik & la poésie, la littérature lyrique \\ \hline
die Symbolik & le symbolisme \\ \hline
berühren (Prät. berührte) & toucher (émotionnellement) \\ \hline
das dramatische Werk & l'œuvre dramatique \\ \hline
die moralische Frage & la question morale \\ \hline
der Literaturnobelpreis & le prix Nobel de littérature \\ \hline
kritisieren & critiquer \\ \hline
der gesellschaftliche Missstand & le dysfonctionnement social \\ \hline
die zeitgenössische Szene & la scène contemporaine \\ \hline
die Migration & la migration \\ \hline
die Identität & l'identité \\ \hline
bemerkenswert & remarquable \\ \hline
die Perspektive & la perspective \\ \hline
verbinden & relier, connecter \\ \hline
der Spiegel & le miroir \\ \hline
die Seele & l'âme \\ \hline
\end{tabularx}
\end{center}

\courssub{Commentaire modèle (niveau Bac)}

\begin{exemplebox}[Commentaire modèle — environ 9 lignes]
\all{Der Text handelt von der reichen Kunst- und Literaturgeschichte im deutschsprachigen Raum (Deutschland, Österreich, Schweiz) und nennt wichtige Vertreter wie Dürer, Friedrich, Goethe, Schiller, Mann, Rilke, Dürrenmatt und Jelinek. Im ersten Teil beschreibt der Autor die Malerei: Friedrichs berühmtes Gemälde „Der Wanderer über dem Nebelmeer“ fasziniert durch seine mystische Atmosphäre. Im zweiten Teil geht es um die Literatur: Goethes „Faust“ und Schillers „Die Räuber“ gehören zum Weltliteraturkanon, während Rilkes Gedichte durch tiefe Symbolik berühren. Der dritte Teil eröffnet den Blick auf die Gegenwart: zeitgenössische Autoren wie Stanišić thematisieren Migration und Identität, und auch kamerunische Schriftsteller wie Nganang werden erwähnt. Meiner Meinung nach zeigt der Text überzeugend, dass Kunst und Literatur Brücken zwischen Kulturen bauen. Besonders beeindruckend finde ich die Verbindung von klassischer und moderner Literatur, weil sie die Kontinuität des kulturellen Erbes verdeutlicht.}
\end{exemplebox}

\begin{aretenir}[Structure du commentaire modèle]
\begin{itemize}
    \item \textbf{Présentation} : \all{Der Text handelt von...} + thème + aire culturelle + noms d'artistes.
    \item \textbf{Idée générale} : Implicitement donnée par l'énumération des époques (classique → moderne → contemporain) et l'ouverture sur le Cameroun.
    \item \textbf{Analyse en 3 parties} : 1) Peinture (Friedrich), 2) Littérature classique (Goethe, Schiller, Rilke, Dürrenmatt, Jelinek), 3) Littérature contemporaine + ouverture africaine. Citations précises (\all{Friedrichs Gemälde}, \all{Goethes „Faust“}, \all{zeitgenössische Autoren}).
    \item \textbf{Avis personnel} : \all{Meiner Meinung nach...}, \all{Besonders beeindruckend finde ich...}, argumenté (ponts entre cultures, continuité du patrimoine).
    \item \textbf{Grammaire réinvestie} : Génitifs (\all{Friedrichs Gemälde}, \all{Goethes „Faust“}, \all{Rilkes Gedichte}), participes épithètes (\all{berühmtes Gemälde}, \all{zeitgenössische Autoren}, \all{kamerunische Schriftsteller}), présent de l'analyse textuelle (\all{handelt, nennt, beschreibt, fasziniert, gehören, berühren, eröffnen, thematisieren, werden erwähnt, zeigt, finde, verdeutlicht}). Le prétérit des verbes forts est travaillé en section 5 et réinvesti dans les exercices de production.
\end{itemize}
\end{aretenir}

\courssub{Résumé guidé du texte}

\begin{methode}[Rédiger un résumé en 3--4 phrases]
\begin{enumerate}
    \item Lisez le texte et soulignez les \textbf{mots-clés} (noms propres, thèmes, verbes principaux).
    \item Regroupez les idées par paragraphes (Peinture — Littérature classique — Littérature contemporaine).
    \item Rédigez une phrase par paragraphe en reliant les mots-clés par des connecteurs logiques (\all{und}, \all{außerdem}, \all{jedoch}, \all{schließlich}).
    \item Vérifiez : accords (participes, génitifs), prétérit si narration passée, orthographe des noms propres.
\end{enumerate}
\end{methode}

\begin{exercice}[Résumé guidé — Complétez les phrases avec les mots-clés fournis]
\textbf{Mots-clés :} \all{kulturelle Landschaft}, \all{Maler}, \all{Gemälde}, \all{Literatur}, \all{Goethe}, \all{Schiller}, \all{Rilke}, \all{zeitgenössische Autoren}, \all{Migration}, \all{Identität}, \all{Kamerun}, \all{verbinden}, \all{Spiegel}.

\begin{enumerate}
    \item Der Text beschreibt die reiche \_\_\_\_\_\_\_\_\_\_\_\_ im deutschsprachigen Raum und nennt berühmte \_\_\_\_\_\_\_\_\_\_\_\_ wie Dürer und Friedrich.
    \item Friedrichs \_\_\_\_\_\_\_\_\_\_\_\_ „Der Wanderer über dem Nebelmeer“ zeigt einen Menschen, der die Natur betrachtet.
    \item In der \_\_\_\_\_\_\_\_\_\_\_\_ ragen \_\_\_\_\_\_\_\_\_\_\_\_, \_\_\_\_\_\_\_\_\_\_\_\_ und \_\_\_\_\_\_\_\_\_\_\_\_ hervor; ihre Werke gehören zum Weltkanon.
    \item \_\_\_\_\_\_\_\_\_\_\_\_ wie Stanišić thematisieren \_\_\_\_\_\_\_\_\_\_\_\_ und \_\_\_\_\_\_\_\_\_\_\_\_; auch Schriftsteller aus \_\_\_\_\_\_\_\_\_\_\_\_ werden erwähnt.
    \item Kunst und Literatur \_\_\_\_\_\_\_\_\_\_\_\_ Menschen über Grenzen hinweg — sie sind der \_\_\_\_\_\_\_\_\_\_\_\_ der Seele einer Nation.
\end{enumerate}
\end{exercice}

\begin{corrige}[Exercice Résumé guidé]
\begin{enumerate}
    \item Der Text beschreibt die reiche \textbf{kulturelle Landschaft} im deutschsprachigen Raum und nennt berühmte \textbf{Maler} wie Dürer und Friedrich.
    \item Friedrichs \textbf{Gemälde} „Der Wanderer über dem Nebelmeer“ zeigt einen Menschen, der die Natur betrachtet.
    \item In der \textbf{Literatur} ragen \textbf{Goethe}, \textbf{Schiller} und \textbf{Rilke} hervor; ihre Werke gehören zum Weltkanon.
    \item \textbf{Zeitgenössische Autoren} wie Stanišić thematisieren \textbf{Migration} und \textbf{Identität}; auch Schriftsteller aus \textbf{Kamerun} werden erwähnt.
    \item Kunst und Literatur \textbf{verbinden} Menschen über Grenzen hinweg — sie sind der \textbf{Spiegel} der Seele einer Nation.
\end{enumerate}
\textit{Remarque :} Le résumé complet fait 5 phrases ici ; pour l'examen, vous pouvez fusionner les phrases 1 et 2, et les phrases 4 et 5, afin d'obtenir 3--4 phrases.
\end{corrige}

\courssub{Production écrite guidée : « Mein Lieblingskünstler / Mein Lieblingsbuch »}

\begin{methode}[Rédiger une présentation personnelle (8--10 lignes)]
\begin{enumerate}
    \item \textbf{Choisir} : Un artiste (peintre, écrivain, musicien, sculpteur) ou un livre/œuvre qui vous marque. Camerounais ou germanophone.
    \item \textbf{Remplir la grille d'aide} ci-dessous (mots-clés en allemand).
    \item \textbf{Rédiger} en phrases complètes : Présentation → Description de l'œuvre → Raisons de votre choix → Conclusion.
    \item \textbf{Relire} avec la grille d'auto-évaluation (voir fin de section).
\end{enumerate}
\end{methode}

\begin{exemplebox}[Grille d'aide — Vocabulaire et structures]
\begin{center}
\begin{tabularx}{\linewidth}{|X|X|X|}\hline
\rowcolor{popblueL} \textbf{Question} & \textbf{Mots-clés / Structures} & \textbf{Exemple} \\ \hline
Wer? (Qui ?) & \all{Mein Lieblingskünstler ist...}, \all{Mein Lieblingsautor heißt...}, \all{Ich bewundere...} & \all{Mein Lieblingskünstler ist Pascal Kenfack.} \\ \hline
Was? (Quoi ?) & \all{Er/Sie malt...}, \all{er/sie schreibt...}, \all{sein/ihr bekanntestes Werk ist...}, \all{der Roman/das Gedicht/das Gemälde heißt...} & \all{Er malt farbenfrohe Bilder. Sein bekanntestes Gemälde heißt „Maske der Ahnen“.} \\ \hline
Warum? (Pourquoi ?) & \all{weil... beeindruckend/tiefgründig/originell ist}, \all{das Werk berührt mich, weil...}, \all{es zeigt...}, \all{es verbindet...} & \all{Weil seine Kunst traditionelle Symbole mit modernen Techniken verbindet.} \\ \hline
Persönlicher Bezug & \all{Ich habe das Werk in... gesehen/gelesen.}, \all{Es erinnert mich an...}, \all{Ich fühle...} & \all{Ich habe das Gemälde in Yaoundé gesehen. Es erinnert mich an meine Kindheit.} \\ \hline
Fazit & \all{Deshalb ist er/sie mein Lieblingskünstler.}, \all{Ich empfehle das Werk jedem, der... interessiert.} & \all{Deshalb ist er mein Lieblingskünstler. Ich empfehle seine Werke allen Kunstliebhabern.} \\ \hline
\end{tabularx}
\end{center}
\end{exemplebox}

\begin{exoresolu}[Production écrite modèle — Mein Lieblingskünstler]
\textbf{Consigne :} Rédigez 8--10 lignes sur votre artiste préféré en utilisant la grille d'aide. Utilisez au moins \textbf{un participe présent en épithète}, \textbf{un participe passé en épithète}, \textbf{un génitif} et \textbf{un verbe fort au prétérit}.

\vspace{0.3cm}
\tcblower
\textbf{Solution détaillée :}

\all{Mein Lieblingskünstler ist der kamerunische Maler Pascal Kenfack. Er wurde 1966 in Bafoussam geboren und studierte in Düsseldorf. Kenfacks Gemälde sind farbenfroh und expressiv; sie verbinden traditionelle afrikanische Motive mit modernen Ausdrucksformen. Sein bekanntestes Werk, das Gemälde „Die Maske der Ahnen“, fasziniert mich durch seine mystische Ausstrahlung. Die dargestellten Figuren scheinen lebendig zu sein. Ich sah das Bild vor zwei Jahren in einer Ausstellung in Yaoundé und war tief beeindruckt. Meiner Meinung nach schafft Kenfack eine Brücke zwischen der kamerunischen Tradition und der zeitgenössischen Kunst. Deshalb bewundere ich ihn sehr und empfehle seine Werke jedem Kunstliebhaber.}

\textbf{Analyse grammaticale :}
\begin{itemize}
    \item Participe présent en épithète (Partizip I) : \all{faszinierend} (ex. \all{ein faszinierendes Werk}) — non utilisé dans ce texte mais requis pour l'exercice. Le texte utilise \all{lebendig} (adjectif) et \all{zeitgenössischen} (adjectif).
    \item Participe passé en épithète (Partizip II) : \all{dargestellten} (de \all{darstellen}, accordé pluriel nominatif/accusatif), \all{beeindruckt} (de \all{beeindrucken}, ici prédicatif : \all{war tief beeindruckt}).
    \item Génitif : \all{Kenfacks Gemälde} (nom propre + \all{s}), \all{die Maske der Ahnen} (article \all{der} = génitif pluriel de \all{die Ahnen}).
    \item Prétérit verbes forts/irréguliers : \all{wurde ... geboren} (\all{werden}, fort), \all{sah} (\all{sehen}, fort), \all{war} (\all{sein}, irrégulier). Le verbe \all{schafft} est au présent ; son prétérit serait \all{schuf}.
\end{itemize}
\end{exoresolu}

\begin{exercice}[Production écrite — À vous de jouer !]
Rédigez votre propre texte « \all{Mein Lieblingskünstler / Mein Lieblingsbuch} » (8--10 lignes) en respectant la grille d'aide et en intégrant :
\begin{itemize}
    \item au moins \textbf{un participe présent en épithète} (\all{lesend, malend, beeindruckend...})
    \item au moins \textbf{un participe passé en épithète} (\all{geschrieben, gemalt, berührt...})
    \item au moins \textbf{un génitif} (\all{Goethes Werk, des Künstlers Stil, der Titel des Buches})
    \item au moins \textbf{un verbe fort au prétérit} (\all{las, sah, ging, fand, schrieb, schuf...})
\end{itemize}
Utilisez le vocabulaire des sections 2 à 5. Écrivez sur une feuille séparée, puis auto-évaluez-vous avec la grille page suivante.
\end{exercice}

\courssub{Production orale guidée : Présentation en binôme}

\begin{experience}[Présentation orale — Un artiste camerounais ou germanophone]
\textbf{Durée :} 15 min préparation + 2--3 min par binôme devant la classe. \\
\textbf{Consigne :} En binôme, choisissez \textbf{un artiste camerounais} (peintre, écrivain, musicien, chorégraphe) \textbf{OU} \textbf{un artiste germanophone} (vu dans le texte ou un autre). Préparez une présentation structurée de 2 minutes. L'un présente la biographie, l'autre l'œuvre et votre avis. Alternez la parole.

\textbf{Structure recommandée :}
\begin{enumerate}
    \item \all{Guten Tag! Wir stellen Ihnen ... vor.} (Présentation du binôme et de l'artiste)
    \item \textbf{Biographie} (naissance, formation, parcours) : \all{Er/Sie wurde ... geboren. Er/Sie studierte in ...}
    \item \textbf{Œuvre majeure} : \all{Sein/Ihr bekanntestes Werk ist ... Es handelt von ...}
    \item \textbf{Analyse} (style, thèmes, techniques) : \all{Die Kunst ist ... / Der Roman thematisiert ...}
    \item \textbf{Avis personnel} : \all{Wir finden das Werk beeindruckend, weil ... Meiner Meinung nach ...}
    \item \all{Vielen Dank für Ihre Aufmerksamkeit. Haben Sie Fragen?}
\end{enumerate}

\textbf{Exemples d'artistes camerounais :} \all{Pascal Kenfack} (peintre), \all{Patrice Nganang} (écrivain), \all{Bate Besong} (poète/dramaturge), \all{Manu Dibango} (musicien), \all{Werewere Liking} (écrivaine/choreographe), \all{Francis Bebey} (musicien/écrivain).

\textbf{Critères d'évaluation (sur 10) :} Prononciation/Intonation (3) — Correction grammaticale (génitif, participes, prétérit) (3) — Richesse lexicale (2) — Structure/cohérence (2).
\end{experience}

\courssub{Auto-évaluation : Grille de relecture avant de rendre votre copie}

\begin{aretenir}[Check-list finale — Vérifiez chaque point avant de remettre votre production]
\begin{center}
\begin{tabular}{|l|c|c|}\hline
\rowcolor{popblueL} \textbf{Critère} & \textbf{Contrôle} & \textbf{Exemple de correction} \\ \hline
\textbf{Articles \& Pluriels} & Tous les noms ont leur article (\all{der/die/das}) et le pluriel correct ? & \all{das Gemälde} → \all{die Gemälde} ; \all{der Künstler} → \all{die Künstler} \\ \hline
\textbf{Génitif} & \textbf{Noms propres} : -\all{s} (\all{Goethes Werk}) ? \textbf{Noms communs masc./neutre} : article \all{des} + -\all{(e)s} (\all{des Künstlers Name}) ? \textbf{Féminin/pluriel} : article \all{der} (\all{der Titel der Bücher}) ? & \all{Friedrichs Gemälde} ✓ — \all{des Buches Titel} ✓ — \all{der Titel der Bücher} ✓ \\ \hline
\textbf{Participes épithètes} & Accord en genre, nombre, cas avec le nom ? \all{der malende Künstler}, \all{das gemalte Bild}, \all{die geschriebene Geschichte} & \all{ein faszinierendes Werk} (neutre, accusatif) — \all{die dargestellten Figuren} (pluriel, nominatif/accusatif) \\ \hline
\textbf{Prétérit verbes forts} & Voyelle correcte ? Pas de -\all{te} ajouté ? \all{schrieb} (pas \all{schreibte}), \all{fand} (pas \all{fandete}), \all{sah} (pas \all{sehte}) & \all{Ich las das Buch.} — \all{Er ging nach Hause.} — \all{Wir schufen etwas Neues.} \\ \hline
\textbf{Ordre des mots} & Verbe en 2e position (phrase principale) ? Verbe à la fin (subordonnée) ? & \all{Meiner Meinung nach \textbf{ist} Kunst wichtig.} — \all{..., \textbf{weil} Kunst \textbf{ist} wichtig.} \\ \hline
\textbf{Majuscules} & \textbf{Tous les noms} en majuscule ? (y compris \all{Kunst, Literatur, Maler, Roman, Gedicht}) & \all{Die Kunst und die Literatur} ✓ \\ \hline
\textbf{Orthographe Umlaute/ß} & \all{ä, ö, ü, ß} corrects ? \all{Größe, Bücher, Künstler, Straße} & \all{für, über, Größe, Bücher} \\ \hline
\end{tabular}
\end{center}
\end{aretenir}

\begin{attention}[Pièges fréquents des francophones — Dernière vérification !]
\begin{itemize}
    \item \textbf{Oublier le -\all{s} du génitif} : \all{*Goethe Werk} → \all{Goethes Werk} ; \all{*der Titel das Buches} → \all{der Titel des Buches}.
    \item \textbf{Accord du participe épithète} : \all{*ein faszinierend Werk} → \all{ein faszinierendes Werk} (neutre, nominatif/accusatif) ; \all{*die dargestellte Figuren} → \all{die dargestellten Figuren} (pluriel).
    \item \textbf{Confondre Partizip I et II} : \all{*das gelesene Buch} (lu — action passée) vs \all{das lesende Kind} (qui lit — action simultanée).
    \item \textbf{Prétérit faible au lieu de fort} : \all{*schreibte} → \all{schrieb} ; \all{*fandete} → \all{fand} ; \all{*gabte} → \all{gab}.
    \item \textbf{Majuscules manquantes} : \all{*die kunst} → \all{die Kunst} ; \all{*ein roman} → \all{ein Roman}.
    \item \textbf{Faux ami} : \all{Kunst} = art (pas « kunst » en français) ; \all{Roman} = roman (pas « romain ») ; \all{eventuell} = éventuel (pas « actuel » → \all{aktuell} = actuel/current).
\end{itemize}
\end{attention}

\begin{savaistu}[Le Cameroun, terre d'artistes — Atout pour l'oral du Bac]
Le Cameroun possède une scène culturelle vibrante : littérature (\all{Patrice Nganang}, \all{Werewere Liking}, \all{Bate Besong}, \all{Ferdinand Oyono}), peinture (\all{Pascal Kenfack}, \all{Joseph Francis Sumégné}), musique (\all{Manu Dibango}, \all{Richard Bona}, \all{Charlotte Dipanda}), théâtre (\all{Jean Ikellé Matiba}). Savoir présenter l'un de ces artistes \textbf{en allemand} lors de l'oral du Baccalauréat (épreuve de compréhension/expression orale) est un atout majeur : cela montre votre capacité à transférer votre culture dans la langue cible, à utiliser le vocabulaire spécialisé (\all{der Maler, der Schriftsteller, das Werk, die Ausstellung, der Roman, das Gedicht}) et à maîtriser les structures grammaticales (génitif : \all{Kenfacks Gemälde}, \all{Nganangs Roman} ; participes : \all{dargestellte Traditionen}, \all{geschriebene Geschichte} ; prétérit : \all{schuf}, \all{schrieb}, \all{begeisterte}). Préparez dès maintenant une fiche de 5 lignes sur un artiste camerounais de votre choix — elle vous servira pour l'oral !
\end{savaistu}

\begin{exoresolu}[Entraînement oral — Fiche de présentation type]
\textbf{Consigne :} Complétez cette fiche pour l'artiste camerounais de votre choix, puis entraînez-vous à la présenter à l'oral (1 min 30) sans lire.

\vspace{0.3cm}
\tcblower
\textbf{Modèle rempli (exemple : Patrice Nganang) :}
\begin{itemize}
    \item \textbf{Artiste :} \all{Patrice Nganang} (écrivain, né en 1970 à Yaoundé)
    \item \textbf{Études :} \all{Er studierte in Yaoundé und in Deutschland (Frankfurt/Main).}
    \item \textbf{Œuvre majeure :} \all{Sein Roman „Der Schatten der Sultane“ (2017) erhielt den Grand Prix de la Littérature Africaine.}
    \item \textbf{Thèmes :} \all{Das Werk thematisiert Kolonialgeschichte, Identität und Migration.}
    \item \textbf{Style :} \all{Die Sprache ist poetisch und politisch.}
    \item \textbf{Avis :} \all{Meiner Meinung nach ist Nganang eine wichtige Stimme der afrikanischen Literatur, weil er Geschichte und Gegenwart verbindet.}
    \item \textbf{Grammaire intégrée :} \all{Nganangs Roman} (génitif), \all{die thematisierte Geschichte} (Part. II épithète), \all{verbindet} (présent), \all{erhielt} (prétérit fort de \all{erhalten}).
\end{itemize}
\end{exoresolu}

\begin{exercice}[Auto-évaluation finale de la leçon AL22]
Cochez les cases correspondant à votre niveau de maîtrise. Pour chaque item non maîtrisé, reportez-vous à la section indiquée.

\begin{center}
\begin{tabular}{|l|c|c|c|l|}\hline
\rowcolor{popblueL} \textbf{Compétence} & \textbf{😊 Maîtrisé} & \textbf{😐 En cours} & \textbf{☹ Difficile} & \textbf{Section à réviser} \\ \hline
Comprendre un texte sur les arts/littérature & & & & 1 (Texte support) \\ \hline
Utiliser le vocabulaire thématique (articles, pluriels) & & & & 2 (Lexique) \\ \hline
Former et accorder le Partizip I en épithète & & & & 3 (Grammaire 1) \\ \hline
Former et accorder le Partizip II en épithète & & & & 3 (Grammaire 1) \\ \hline
Construire le génitif (noms propres, communs, articles) & & & & 4 (Grammaire 2) \\ \hline
Conjuguer les verbes forts au prétérit & & & & 5 (Conjugaison) \\ \hline
Rédiger un commentaire structuré (4 étapes) & & & & 6 (Cette section) \\ \hline
Produire un résumé cohérent & & & & 6 (Résumé guidé) \\ \hline
Rédiger une présentation personnelle (écrit) & & & & 6 (Production écrite) \\ \hline
Présenter un artiste à l'oral (binôme) & & & & 6 (Production orale) \\ \hline
\end{tabular}
\end{center}

\textbf{Objectif :} Tous les items en 😊 avant le contrôle continu ou le Bac. Bon travail et \all{viel Erfolg} pour la suite de votre apprentissage de l'allemand !
\end{exercice}

\newpage

\coursec{Activité d'intégration}

\courssub{Objectifs de l'activité}
\noindent Cette activité d'intégration vise à mobiliser l'ensemble des savoirs et savoir-faire travaillés au cours de la leçon AL22 dans une tâche complexe, authentique et motivante. À l'issue de cette séance, l'élève doit être capable de :
\begin{itemize}
    \item \textbf{Mobiliser le lexique thématique} : employer avec exactitude les noms (avec article et pluriel), verbes et adjectifs du champ des arts (\all{die Kunst, der Maler, das Gemälde, die Ausstellung, der Bildhauer, die Skulptur}) et de la littérature (\all{der Schriftsteller, der Roman, das Gedicht, die Lesung, der Verlag, die Erstveröffentlichung}).
    \item \textbf{Appliquer la grammaire} : construire des syntagmes nominaux au \textbf{génitif} (\all{die Werke des Malers, die Themen der Romane}) et utiliser les \textbf{participes présent et passé comme épithètes} (\all{ein beeindruckendes Gemälde, der wandernde Mann, die gespielte Musik}) pour enrichir la description.
    \item \textbf{Conjuguer au prétérit} : narrer une biographie ou un événement passé en employant correctement le \textbf{Präteritum des verbes forts} (\all{schrieb, las, ging, fand, begann}) et des verbes faibles (\all{malte, lebte, arbeitete}).
    \item \textbf{Produire un texte cohérent} : rédiger un texte de présentation structuré (6 à 8 phrases) respectant la progression thématique (biographie $\rightarrow$ œuvre majeure $\rightarrow$ analyse formelle $\rightarrow$ portée actuelle).
    \item \textbf{Communiquer à l'oral} : présenter l'exposition devant la classe en s'appuyant sur l'affiche, avec une prononciation soignée (voyelles longues/courtes, \all{ß}, \all{ch}, accent tonique) et une intonation expressive.
\end{itemize}

\courssub{Situation}
\noindent Le \textbf{Lycée Général Leclerc de Yaoundé} organise, en partenariat avec le \textbf{Goethe-Institut Kamerun}, la \emph{Woche der deutsch-kamerunischen Kultur} (Semaine de la culture germano-camerounaise). Le point d'orgue est une grande exposition temporaire dans le hall principal : \all{„Brücken bauen – Kunst und Literatur im Dialog“} (« Construire des ponts — arts et littérature en dialogue »).

\begin{textebox}[Appel à projets affiché dans le hall et diffusé sur la plateforme de l'établissement]
\textbf{Projet : „Unsere Kulturausstellung“}

Chers élèves de Terminale A,

Dans le cadre de la semaine culturelle, votre classe de LV2 Allemand est invitée à concevoir \textbf{un stand d'exposition} consacré à un duo d'artistes : \textbf{un artiste germanophone} (Allemagne, Autriche, Suisse) et \textbf{un artiste camerounais}. L'objectif est de montrer des résonances, des contrastes ou des influences entre deux univers créatifs.

\textbf{Votre mission (par groupes de 4 élèves) :}
\begin{enumerate}
    \item Choisir un duo pertinent (ex. : Caspar David Friedrich / Pascale Marthine Tayou \textbf{ou} Johann Wolfgang von Goethe / Mongo Beti \textbf{ou} Käthe Kollwitz / Francis Bebey \textbf{ou} Franz Kafka / Werewere Liking).
    \item Réaliser \textbf{une affiche A3} (titre, visuels, dates, noms, citation courte).
    \item Rédiger \textbf{le texte du petit catalogue} (6 à 8 phrases en allemand) qui sera imprimé et distribué aux visiteurs.
    \item Présenter votre stand à l'oral (2 minutes par groupe) lors de l'inauguration.
\end{enumerate}

\textbf{Contraintes linguistiques obligatoires dans le texte du catalogue :}
\begin{itemize}
    \item Au moins \textbf{3 verbes au Präteritum} (verbes forts de préférence) pour la biographie.
    \item Au moins \textbf{2 participes I ou II en fonction d'épithète} (attribut du nom) pour décrire les œuvres.
    \item Au moins \textbf{2 constructions au Génitif} (possession, appartenance, thème).
    \item Orthographe soignée : majuscules aux noms, \all{ß} / \all{ss}, tréma, guillemets allemands „…“.
\end{itemize}

Le jury (professeur d'allemand, documentaliste, représentant du Goethe-Institut) récompensera les trois meilleures productions (originalité du duo, qualité de l'allemand, esthétique de l'affiche, aisance orale). Les affiches resteront exposées deux semaines.
\vspace{0.2cm}

\noindent \textit{Bonne création ! / Viel Erfolg und viel Kreativität!}
\end{textebox}

\courssub{Consignes}
\noindent Le travail se déroule en \textbf{quatre phases} sur la séance de 2 heures (et éventuellement un début de séance suivante pour les présentations orales).

\textbf{Phase 1 : Constitution des groupes et choix du duo (15 min)}
\begin{enumerate}
    \item Former les groupes de 4 (rôles suggérés : \textit{rédacteur principal}, \textit{documentaliste}, \textit{graphiste}, \textit{orateur principal}).
    \item Choisir le duo d'artistes. Valider le choix auprès du professeur (éviter les doublons). Utiliser les tablettes ou la salle informatique pour vérifier 3 à 4 dates clés et 2 œuvres majeures par artiste.
\end{enumerate}

\textbf{Phase 2 : Rédaction du texte de catalogue (45 min)}
\begin{enumerate}
    \item \textbf{Étape 1 (brouillon) :} chaque groupe rédige 6 à 8 phrases en allemand sur papier brouillon.
        \begin{itemize}
            \item Phrases 1--2 : biographie de l'artiste A (Präteritum).
            \item Phrases 3--4 : biographie de l'artiste B (Präteritum).
            \item Phrase 5 : description de l'œuvre A (participe épithète + génitif).
            \item Phrase 6 : description de l'œuvre B (participe épithète + génitif).
            \item Phrase 7 : lien, comparaison, portée actuelle (présent ou prétérit).
            \item Phrase 8 (optionnelle) : invitation au public ou citation.
        \end{itemize}
    \item \textbf{Étape 2 (auto-correction guidée) :} vérifier à l'aide de la fiche \textbf{Ressources à mobiliser} (ci-dessous) :
        \begin{itemize}
            \item Les formes du Präteritum des verbes forts (voyelle radicale modifiée, pas de \textit{-t} à la 3\textsuperscript{e} personne du singulier, \textit{-en} à \textit{wir/sie/Sie}).
            \item L'accord des participes épithètes (genre, nombre, cas du nom noyau $\rightarrow$ \all{das beeindruckend\textbf{e} Gemälde}, \all{der wandernd\textbf{e} Mann}).
            \item Les terminaisons du génitif (\all{des Malers}, \all{der Schriftstellerin}, \all{der Künstler}, \all{des Volkes}).
            \item La place du verbe (position 2 en phrase principale, position finale en subordonnée si \all{weil} ou \all{dass} sont utilisés).
        \end{itemize}
    \item \textbf{Étape 3 (version finale) :} recopier le texte corrigé proprement sur la fiche catalogue fournie (écriture lisible, interligne 1,5).
\end{enumerate}

\textbf{Phase 3 : Création de l'affiche (30 min)}
\begin{enumerate}
    \item Sur feuille A3 (fournie) ou logiciel (Canva, PowerPoint), composer l'affiche :
        \begin{itemize}
            \item Titre accrocheur en allemand (\all{„Zwei Welten – Ein Blick“}).
            \item Photos ou reproductions des deux artistes et d'une œuvre de chacun (imprimées ou dessinées).
            \item Noms, dates de vie (ex. \all{1749–1832} pour Goethe), nationalité.
            \item Une citation courte (originale ou plausible) en allemand entre guillemets „…“.
            \item Le texte du catalogue (version finale) imprimé en petit format ou manuscrit soigné, collé au dos ou sur le côté.
        \end{itemize}
    \item Préparer l'oral : l'orateur principal répète la lecture expressive du texte ; les autres préparent chacun une phrase de commentaire personnel (\all{„Mich beeindruckt …“, „Ich finde die Verbindung interessant, weil …“}).
\end{enumerate}

\textbf{Phase 4 : Inauguration et vote (30 min — peut déborder sur la séance suivante)}
\begin{enumerate}
    \item Installation des affiches sur les murs ou les tables du hall (ou de la salle de classe).
    \item Passage de chaque groupe (2 min maximum) : lecture du catalogue + commentaire personnel.
    \item Vote à bulletin secret par tous les élèves (critères : \textit{qualité linguistique}, \textit{originalité du lien}, \textit{esthétique}, \textit{expression orale}).
    \item Annonce des résultats et remise de prix symboliques (goodies du Goethe-Institut, livres, chocolats).
\end{enumerate}

\courssub{Ressources à mobiliser}
\noindent Cette fiche mémo est distribuée à chaque groupe pour l'auto-correction (Phase 2, Étape 2).

\begin{definition}[Lexique clé : arts et littérature (rappel)]
\textbf{Personnes :} \all{der Maler, die Maler} (le peintre) | \all{der Bildhauer, die Bildhauer} (le sculpteur) | \all{der Schriftsteller, die Schriftsteller} (l'écrivain) | \all{die Schriftstellerin, die Schriftstellerinnen} (la femme écrivain) | \all{der Dichter, die Dichter} (le poète) | \all{der Künstler, die Künstler} / \all{die Künstlerin, die Künstlerinnen} (l'artiste).
\textbf{Œuvres et lieux :} \all{das Gemälde, die Gemälde} (le tableau) | \all{das Bild, die Bilder} (l'image, le tableau) | \all{die Skulptur, die Skulpturen} (la sculpture) | \all{der Roman, die Romane} (le roman) | \all{das Gedicht, die Gedichte} (le poème) | \all{die Novelle, die Novellen} (la nouvelle) | \all{die Ausstellung, die Ausstellungen} (l'exposition) | \all{das Museum, die Museen} (le musée) | \all{die Galerie, die Galerien} (la galerie) | \all{die Vernissage, die Vernissagen} (le vernissage).
\textbf{Verbes utiles (Präteritum) :} \all{schreiben $\rightarrow$ schrieb, geschrieben} | \all{lesen $\rightarrow$ las, gelesen} | \all{malen $\rightarrow$ malte, gemalt} (faible) | \all{zeichnen $\rightarrow$ zeichnete, gezeichnet} (faible) | \all{schaffen $\rightarrow$ schuf, geschaffen} (fort, « créer ») | \all{denken $\rightarrow$ dachte, gedacht} (mixte) | \all{finden $\rightarrow$ fand, gefunden} | \all{geben $\rightarrow$ gab, gegeben} | \all{gehen $\rightarrow$ ging, gegangen} | \all{stehen $\rightarrow$ stand, gestanden} | \all{sehen $\rightarrow$ sah, gesehen} | \all{sprechen $\rightarrow$ sprach, gesprochen} | \all{werden $\rightarrow$ wurde, geworden}.
\textbf{Adjectifs et participes utiles :} \all{berühmt} (célèbre), \all{bedeutend} (important), \all{kritisch} (critique), \all{romantisch} (romantique), \all{expressionistisch} (expressionniste), \all{zeitgenössisch} (contemporain), \all{beeindruckend} (impressionnant), \all{bewegend} (émouvant), \all{umstritten} (controversé), \all{bekannt} (connu), \all{unbekannt} (inconnu).
\end{definition}

\begin{propriete}[Le génitif : terminaisons et usage]
\textbf{1. Article défini + nom} (masculin et neutre : \textbf{-s} ou \textbf{-es} sur le nom ; féminin et pluriel : pas de marque sur le nom, article \all{der}) :
\begin{center}
\begin{tabular}{|l|l|l|l|l|}
\hline
\rowcolor{popblueL} Cas & Maskulin & Neutrum & Feminin & Plural \\ \hline
Nominativ & \all{der Maler} & \all{das Werk} & \all{die Schriftstellerin} & \all{die Künstler} \\ \hline
Genitiv & \all{des Maler\textbf{s}} & \all{des Werk\textbf{es}} & \all{der Schriftstellerin} & \all{der Künstler} \\ \hline
\end{tabular}
\end{center}
\textit{Règle du choix -s / -es :} les noms monosyllabiques et les noms terminés par \textit{-s, -ß, -x, -z} prennent en général \textbf{-es} (\all{des Buches, des Volkes, des Schmerzes}) ; les noms polysyllabiques prennent \textbf{-s} (\all{des Malers, des Romans, des Dichters}). Les deux formes sont souvent possibles (\all{des Werks / des Werkes}).
\textbf{2. Article indéfini, possessif, négation (\all{kein}) :} l'article prend la marque du génitif (\all{eines, meines, keines} au masculin et au neutre ; \all{einer, meiner, keiner} au féminin et au pluriel) et l'adjectif prend la terminaison \textbf{-en} : \all{eines berühmten Malers, keines bekannten Romans, die Werke meiner Lieblingsschriftstellerin}.
\textbf{3. Emplois :} appartenance (\all{der Pinsel des Malers} — le pinceau du peintre), thème (\all{die Themen der Romane} — les thèmes des romans), après certaines prépositions (\all{wegen des Regens, während der Ausstellung, statt des Konzerts, außerhalb des Museums}).
\end{propriete}

\begin{propriete}[Les participes comme épithètes : formation et déclinaison]
\textbf{Partizip I (participe présent) :} \all{Infinitiv + -d} $\rightarrow$ \all{wandernd, schreibend, lesend, malend, beeindruckend}. \textit{Sens :} action simultanée et active. \all{der wandernde Mann} = « l'homme qui se promène ».
\textbf{Partizip II (participe passé) :} verbes forts \all{ge-...-en} (\all{geschrieben, gelesen, geschaffen}), verbes faibles \all{ge-...-t} (\all{gemalt, gearbeitet}). \textit{Sens :} passif (action subie), résultat, antériorité. \all{das gemalte Bild} = « le tableau qui a été peint ».
\textbf{Déclinaison (après article défini, nominatif/accusatif) :}
\begin{center}
\begin{tabular}{|l|l|l|l|}
\hline
\rowcolor{popblueL} & Maskulin & Neutrum & Feminin \\ \hline
Partizip I & \all{der wandernd\textbf{e} Mann} & \all{das lesend\textbf{e} Kind} & \all{die singend\textbf{e} Frau} \\ \hline
Partizip II & \all{der geschrieben\textbf{e} Brief} & \all{das gemalt\textbf{e} Bild} & \all{die geschrieben\textbf{e} Geschichte} \\ \hline
\end{tabular}
\end{center}
\cle{Attention} : après l'article défini, le participe épithète prend la terminaison faible \textbf{-e} au nominatif singulier (et à l'accusatif neutre et féminin), \textbf{-en} partout ailleurs. Après l'article indéfini \all{ein}, le masculin et le neutre au nominatif prennent la marque forte : \all{ein wandernd\textbf{er} Mann, ein lesend\textbf{es} Kind, ein beeindruckend\textbf{es} Gemälde}.
\end{propriete}

\begin{propriete}[Präteritum des verbes forts : 3\textsuperscript{e} personne du singulier et du pluriel]
\begin{center}
\begin{tabular}{|l|l|l|l|}
\hline
\rowcolor{popblueL} Infinitif & Präteritum \textit{ich/er/sie/es} & Präteritum \textit{wir/sie/Sie} & Partizip II \\ \hline
\all{schreiben} & \all{schrieb} & \all{schrieben} & \all{geschrieben} \\ \hline
\all{lesen} & \all{las} & \all{lasen} & \all{gelesen} \\ \hline
\all{schaffen} & \all{schuf} & \all{schufen} & \all{geschaffen} \\ \hline
\all{finden} & \all{fand} & \all{fanden} & \all{gefunden} \\ \hline
\all{geben} & \all{gab} & \all{gaben} & \all{gegeben} \\ \hline
\all{gehen} & \all{ging} & \all{gingen} & \all{gegangen} \\ \hline
\all{stehen} & \all{stand} & \all{standen} & \all{gestanden} \\ \hline
\all{sprechen} & \all{sprach} & \all{sprachen} & \all{gesprochen} \\ \hline
\all{treffen} & \all{traf} & \all{trafen} & \all{getroffen} \\ \hline
\all{werden} & \all{wurde} & \all{wurden} & \all{geworden} \\ \hline
\end{tabular}
\end{center}
\cle{Rappel} : pas de \textit{-t} à la 3\textsuperscript{e} personne du singulier pour les verbes forts (\all{er schrieb}, jamais \all{*er schriebt}) ; la voyelle radicale est modifiée (Ablaut). La 2\textsuperscript{e} personne du singulier prend \textit{-st} : \all{du schriebst, du lasest (du lasst)}.
\end{propriete}

\courssub{Proposition de production et corrigé commenté}
\noindent Voici une production modèle pour le duo \textbf{Caspar David Friedrich (1774–1840), peintre romantique allemand} et \textbf{Mongo Beti (1932–2001), écrivain camerounais}. Ce texte sert de référence pour l'enseignant et de modèle de correction en classe.

\begin{exoresolu}[Texte de catalogue modèle — duo Friedrich / Beti]
\textbf{Énoncé :} rédiger le texte de présentation (7 à 8 phrases) pour l'affiche du duo Friedrich / Beti, en respectant les contraintes (3 Präteritum, 2 participes épithètes, 2 génitifs).

\tcblower

\textbf{Produktion (Modelltext) :}

\all{Caspar David Friedrich lebte im 19. Jahrhundert in Deutschland.} \\
\all{Er malte viele berühmte Landschaften mit mystischer Stimmung.} \\
\all{Sein Hauptwerk „Der Wanderer über dem Nebelmeer“ zeigt einen \textbf{wandernden} Mann auf einem Felsen.} \\
\all{Die \textbf{gemalten} Farben \textbf{des Bildes} sind dunkel und tief.} \\
\all{Mongo Beti \textbf{schrieb} in der Kolonialzeit kritische Romane über Kamerun.} \\
\all{Sein Roman „Le pauvre Christ de Bomba“ \textbf{kritisiert} die Missionare und die Kolonialverwaltung.} \\
\all{Die Themen \textbf{des Schriftstellers} sind Identität, Freiheit und Widerstand.} \\
\all{Beide Künstler verbinden europäische und afrikanische Perspektiven im Dialog der Kulturen.}

\textbf{Traduction :}
Caspar David Friedrich vécut au XIX\textsuperscript{e} siècle en Allemagne. Il peignit de nombreux paysages célèbres à l'atmosphère mystique. Son chef-d'œuvre « Le Voyageur au-dessus de la mer de brouillard » montre un homme \textbf{se promenant} (participe I) sur un rocher. Les couleurs \textbf{peintes} (participe II) \textbf{du tableau} (génitif) sont sombres et profondes. Mongo Beti \textbf{écrivit} à l'époque coloniale des romans critiques sur le Cameroun. Son roman « Le Pauvre Christ de Bomba » \textbf{critique} (présent : portée actuelle) les missionnaires et l'administration coloniale. Les thèmes \textbf{de l'écrivain} (génitif) sont l'identité, la liberté et la résistance. Les deux artistes relient des perspectives européennes et africaines dans le dialogue des cultures.

\textbf{Commentaire linguistique détaillé (pour la correction collective) :}
\begin{enumerate}
    \item \textbf{\all{lebte, malte}} : Präteritum des verbes faibles (\all{leben, malen}), terminaison \textit{-te}. Correct.
    \item \textbf{\all{zeigt}} : présent (description intemporelle d'une œuvre). Le prétérit \all{zeigte} serait possible, mais le présent est préféré pour l'analyse d'œuvre.
    \item \textbf{\all{wandernden}} : \textbf{Partizip I épithète}, attribut de \all{einen Mann} (masculin accusatif singulier, article défini $\rightarrow$ terminaison \textit{-en}). Sens actif : l'homme \textit{est en train de} se promener. Accord correct.
    \item \textbf{\all{gemalten}} : \textbf{Partizip II épithète}, attribut de \all{die Farben} (pluriel, article défini $\rightarrow$ \textit{-en}). Sens passif et résultatif : les couleurs \textit{ont été peintes}. \textbf{\all{des Bildes}} : \textbf{génitif} neutre singulier (\all{das Bild} $\rightarrow$ \all{des Bild\textbf{es}}) ; marque \textit{-es} car le nom est monosyllabique. Double marque du cas (article \all{des} + \textit{-es} sur le nom). Excellent.
    \item \textbf{\all{schrieb}} : \textbf{Präteritum du verbe fort} \all{schreiben} (Ablaut \textit{ei} $\rightarrow$ \textit{ie}). 3\textsuperscript{e} personne du singulier sans \textit{-t}. Correct.
    \item \textbf{\all{kritisiert}} : présent (fait littéraire permanent). \all{kritisierte} (prétérit) serait accepté pour rester dans le récit, mais le présent souligne l'actualité de la critique.
    \item \textbf{\all{des Schriftstellers}} : \textbf{génitif} masculin singulier (\all{der Schriftsteller} $\rightarrow$ \all{des Schriftsteller\textbf{s}}). Correct.
    \item \textbf{\all{verbinden}} : présent (énoncé de portée générale). \textbf{\all{europäische, afrikanische}} : adjectifs au pluriel avec terminaison \textit{-n} après \all{beide} sous-entendu (\all{beide Künstler}).
\end{enumerate}

\textbf{Points de vigilance pour les groupes (erreurs fréquentes attendues) :}
\begin{itemize}
    \item \all{der wandernde Mann} (correct), mais \all{ein wandernd\textbf{er} Mann} : avec l'article indéfini, déclinaison forte en \textit{-er}.
    \item \all{*des Maler} : oubli du \textit{-s} au génitif masculin/neutre $\rightarrow$ \all{des Maler\textbf{s}}.
    \item \all{*er schreibt} (présent au lieu du prétérit) ou \all{*er schriebt} (le \textit{-t} est interdit à la 3\textsuperscript{e} personne du singulier des verbes forts).
    \item \all{*das gemalt Bild} : oubli de la déclinaison du participe II $\rightarrow$ \all{das gemalt\textbf{e} Bild}.
    \item Confusion entre \all{malen} (faible : \all{malte}) et \all{schaffen} (fort : \all{schuf}) ; attention aussi à \all{schaffen} faible (\all{schaffte}) au sens de « réussir, gérer ».
\end{itemize}
\end{exoresolu}

\courssub{Bilan de l'activité}
\noindent Cette activité de synthèse permet de valider les acquis de la leçon AL22 dans une situation de communication réelle et valorisante. Le bilan se fait à trois niveaux : auto-évaluation de l'élève, retour collectif de l'enseignant, ouverture culturelle.

\begin{aretenir}[Bilan linguistique — grille d'auto-évaluation élève (à distribuer)]
L'élève coche : \textbf{A} (acquis) / \textbf{EC} (en cours) / \textbf{NA} (non acquis).
\begin{center}
\begin{tabularx}{\linewidth}{|X|c|c|c|}
\hline
\rowcolor{popblueL} \textbf{Compétence / savoir} & \textbf{A} & \textbf{EC} & \textbf{NA} \\ \hline
Je connais le vocabulaire des arts et de la littérature (noms + articles + pluriels). & & & \\ \hline
J'écris correctement le \textbf{génitif} (article + nom + terminaison \textit{-s}/\textit{-es}). & & & \\ \hline
J'utilise le \textbf{Partizip I} (\textit{-nd}) comme adjectif (ex. \all{wandernd}). & & & \\ \hline
J'utilise le \textbf{Partizip II} (\textit{ge-...-t/en}) comme adjectif (ex. \all{gemalt}). & & & \\ \hline
Je décline correctement le participe épithète (accord en genre, nombre, cas). & & & \\ \hline
Je conjugue les \textbf{verbes forts au Präteritum} (3\textsuperscript{e} personne sg./pl.). & & & \\ \hline
Je raconte une biographie au passé (Präteritum) de façon fluide. & & & \\ \hline
Je produis un texte cohérent de 6 à 8 phrases (connecteurs, ordre des mots). & & & \\ \hline
Je présente mon travail à l'oral (prononciation, regard, débit). & & & \\ \hline
\end{tabularx}
\end{center}
\end{aretenir}

\begin{methode}[Bilan pédagogique — retour enseignant (oral, 5 min en fin de séance)]
\begin{enumerate}
    \item \textbf{Valoriser :} souligner la créativité des duos choisis (ex. \all{Käthe Kollwitz / Francis Bebey} : gravure sociale et musique engagée ; \all{Franz Kafka / Werewere Liking} : absurdité bureaucratique et théâtre rituel). Féliciter les usages spontanés réussis du génitif (\all{die Stimme des Volkes}) ou du participe I (\all{die singende Menge}).
    \item \textbf{Corriger collectivement 2 à 3 erreurs types} relevées pendant la circulation (ex. confusion entre \all{des Buches} [singulier] et \all{der Bücher} [pluriel], ou \all{*er schreibt} au lieu de \all{er schrieb} dans la biographie).
    \item \textbf{Ancrer :} rappeler que le \textbf{génitif} reste vivant dans la langue écrite formelle (presse, catalogues, administration) même s'il recule à l'oral au profit de \all{von + Datif}. Les \textbf{participes épithètes} sont la marque d'un style écrit élaboré (\all{Schriftsprache}) : \all{die betreffende Person, die beschlossene Sache}.
    \item \textbf{Ouvrir :} cette compétence « décrire une œuvre, présenter un artiste » est réinvestissable en Terminale pour l'\textbf{épreuve orale du Baccalauréat} (exposé sur un thème culturel) et dans les études supérieures (fiches de lecture, critiques d'exposition).
\end{enumerate}
\end{methode}

\begin{savaistu}[Culture : le Goethe-Institut au Cameroun]
Le \textbf{Goethe-Institut Kamerun} (siège à Yaoundé, antenne à Douala) est l'institution culturelle officielle de la République fédérale d'Allemagne. Il promeut la langue allemande (cours, examens \all{Goethe-Zertifikat}, \all{TestDaF}), la coopération universitaire (avec le DAAD, \all{Deutscher Akademischer Austauschdienst}) et les échanges artistiques. Il organise régulièrement des expositions d'artistes allemands contemporains (photographie, art vidéo) et soutient des projets d'artistes camerounais (résidences à Berlin, Weimar ou Düsseldorf). Des semaines culturelles germano-camerounaises y sont effectivement organisées. Participer à de tels événements (vernissages, lectures publiques — \all{Lesungen} —, concerts) est le meilleur moyen de pratiquer l'allemand hors de la classe et de découvrir la \all{Kulturszene} (scène culturelle) vibrante des deux pays.
\end{savaistu}

\newpage

\coursec{Méthodes et exercices résolus}

\courssub{Méthodes et exercices résolus}

\begin{methode}[Commentaire de texte : compréhension globale et détaillée (Épreuve du Bac)]
\textbf{Objectif :} Répondre en allemand à des questions fermées et ouvertes sur un texte littéraire ou journalistique.
\begin{enumerate}
    \item \textbf{Lecture globale (2 min) :} Identifier le \cle{type de texte} (article, extrait de roman, critique), le \cle{thème} (arts, littérature), la \cle{source} et l'\cle{intention} (informer, critiquer, divertir).
    \item \textbf{Lecture détaillée (10 min) :} Surligner les \cle{mots-clés} du sujet, les \cle{connecteurs logiques} (\all{denn, aber, deshalb, zwar ... aber}), les \cle{participes en épithète} (indices de densité informative) et les \cle{verbes au prétérit} (marque du récit).
    \item \textbf{Analyse des questions :} Repérer l'opérateur : \all{Was?} (contenu), \all{Warum? / Weshalb?} (cause), \all{Wie?} (manière), \all{Inwiefern?} (évaluation/nuance).
    \item \textbf{Rédaction des réponses :}
    \begin{itemize}
        \item Répondre \textbf{en allemand}, en phrases complètes.
        \item Citer le texte : \all{Im Text heißt es: „...“} ou \all{Laut Text ...}.
        \item Reformuler avec son propre vocabulaire (pas de copier-coller brut de trois lignes).
        \item Respecter l'\cle{ordre des mots} (verbe conjugué en 2\up{e} position en principale, verbe à la fin en subordonnée).
    \end{itemize}
    \item \textbf{Relecture (3 min) :} Vérifier les cas (génitif !), les majuscules aux noms, les terminaisons adjectivales, la concordance des temps.
\end{enumerate}
\cle{Réflexe Bac} : Pour une question \all{Inwiefern ...}, structurez : \all{Einerseits ... Andererseits ... / Zum einen ... Zum anderen ...}.
\end{methode}

\begin{exoresolu}[Compréhension du texte : „Kunst und Literatur im deutschsprachigen Raum“]
\textbf{Énoncé :} Lisez le texte ci-dessous et répondez aux questions en allemand (phrases complètes).

\all{Die zeitgenössische Kunstszene in Deutschland, Österreich und der Schweiz ist vielfältig und international vernetzt. Junge Malerinnen und Maler stellen ihre Werke nicht nur in klassischen Galerien, sondern auch im öffentlichen Raum aus. Street Art und Graffiti, lange als Vandalismus verpönt, gelten heute als anerkannte Kunstformen. Ein bekanntes Beispiel ist die East Side Gallery in Berlin, ein bemaltes Stück der ehemaligen Mauer, das für Freiheit und Einheit steht.}

\all{Auch die Literatur erlebt einen Aufschwung. Junge Schriftstellerinnen und Schriftsteller thematisieren Identität, Migration und Zugehörigkeit in ihren Romanen und Essays. Der Deutsche Buchpreis, jährlich im Oktober verliehen, lenkt die Aufmerksamkeit auf vielversprechende Debüts. Lesungen in Buchhandlungen und Literaturfestivals ziehen tausende Besucher an.}

\all{Dennoch bleibt die Finanzierung prekär. Viele Künstler leben von Projektförderungen und Nebenjobs. Die Corona-Pandemie hat die Lage verschärft: Ausstellungen wurden abgesagt, Ateliers standen leer. Der Staat reagierte mit Hilfsprogrammen, doch die strukturelle Unsicherheit bleibt.}

\textbf{Glossaire :} \all{verpönt} = banni, mal vu ; \all{die Mauer} = le Mur ; \all{der Aufschwung} = l'essor ; \all{das Debüt} = le premier livre ; \all{prekär} = précaire ; \all{die Projektförderung} = la subvention de projet.

\textbf{Fragen :}
\begin{enumerate}
    \item Was versteht man unter der „East Side Gallery“ und was symbolisiert sie?
    \item Nennen Sie zwei Themen, die junge deutschsprachige Schriftsteller in ihren Werken behandeln.
    \item Warum ist die finanzielle Lage vieler Künstler „prekär“? Nennen Sie zwei Gründe aus dem Text.
    \item \textbf{Expression personnelle :} „Kunst ist ein Luxus, den sich eine Gesellschaft in Krisenzeiten nicht leisten kann.“ Nehmen Sie Stellung (ca. 60--80 Wörter).
\end{enumerate}

\tcblower
\textbf{Solution détaillée :}

\textbf{1. Analyse de la question 1 :} \all{Was ... symbolisiert sie?} $\rightarrow$ Information explicite + sens symbolique.
\textit{Repérage dans le texte :} \all{„Ein bekanntes Beispiel ist die East Side Gallery in Berlin, ein bemaltes Stück der ehemaligen Mauer, das für Freiheit und Einheit steht.“}
\textbf{Réponse :} \all{Die East Side Gallery ist ein bemaltes Stück der ehemaligen Berliner Mauer. Sie symbolisiert Freiheit und Einheit.}

\textbf{2. Analyse de la question 2 :} \all{Nennen Sie zwei Themen ...} $\rightarrow$ Recherche ciblée.
\textit{Repérage :} \all{„... thematisieren Identität, Migration und Zugehörigkeit ...“}
\textbf{Réponse :} \all{Junge Schriftsteller thematisieren Identität und Migration (oder: Zugehörigkeit) in ihren Werken.}

\textbf{3. Analyse de la question 3 :} \all{Warum ... prekär? Zwei Gründe.} $\rightarrow$ Causalité.
\textit{Repérage :} \all{„Viele Künstler leben von Projektförderungen und Nebenjobs.“} / \all{„Die Corona-Pandemie hat die Lage verschärft: Ausstellungen wurden abgesagt, Ateliers standen leer.“}
\textbf{Réponse :} \all{Die Lage ist prekär, weil viele Künstler von Projektförderungen und Nebenjobs abhängen. Zudem hat die Corona-Pandemie zu abgesagten Ausstellungen und leeren Ateliers geführt.}

\textbf{4. Expression personnelle (Méthode : Thèse -- Arguments -- Exemple -- Conclusion) :}
\textit{Thèse :} \all{Ich stimme nicht zu. Gerade in Krisenzeiten ist Kunst wichtig.}
\textit{Argument 1 :} \all{Sie bietet Trost, Identität und kritische Reflexion.}
\textit{Argument 2 + exemple :} \all{Während der Pandemie halfen digitale Konzerte und Lesungen gegen die Isolation.}
\textit{Conclusion :} \all{Kunst ist kein Luxus, sondern Nahrung für die Seele.}

\textbf{Production écrite (modèle) :}
\all{Ich teile diese Meinung nicht. Kunst ist gerade in Krisenzeiten unverzichtbar, kein Luxus. Sie spendet Trost, stärkt die Identität und ermöglicht kritische Reflexion über die Gesellschaft. Während der Corona-Pandemie haben digitale Konzerte, Lesungen und virtuelle Museumsbesuche vielen Menschen geholfen, Isolation und Angst zu überwinden. Kunst schafft Gemeinschaft und hält die Demokratie lebendig. Daher muss der Staat Kultur auch in Krisen fördern.}
\end{exoresolu}

\begin{methode}[Exploitation lexicale : familles de mots, genres, pluriels et champs sémantiques]
\textbf{Objectif :} Maîtriser le vocabulaire thématique (arts, littérature) pour l'expression écrite et orale et pour la version.
\begin{enumerate}
    \item \textbf{Fiche par famille :} Notez le \cle{nom de base} (article + pluriel), le \cle{verbe} associé, l'\cle{adjectif}, le \cle{nom d'agent} (\all{-er / -in}) et le \cle{nom d'action} (\all{-ung, -erei, -keit}).
    \item \textbf{Genre et pluriel :} Apprenez \textbf{systématiquement} l'article et le pluriel (\all{das Gemälde, die Gemälde} ; \all{der Roman, die Romane} ; \all{das Gedicht, die Gedichte}). Attention aux \cle{Umlaute} au pluriel : \all{der Maler -- die Maler} (pas d'Umlaut), mais \all{die Kunst -- die Künste}, \all{der Roman -- die Romane}.
    \item \textbf{Champs lexicaux :} Regroupez : \textit{Lieux} (\all{die Galerie, das Museum, das Theater, die Bühne, das Atelier}), \textit{Acteurs} (\all{der Künstler, der Kritiker, der Verleger, das Publikum}), \textit{Genres} (\all{der Roman, das Gedicht, das Drama, die Novelle, die Kurzgeschichte}), \textit{Verbes d'action} (\all{veröffentlichen, ausstellen, rezensieren, inszenieren, komponieren}).
    \item \textbf{Faux amis et pièges :} \all{das Kunstwerk} (l'œuvre d'art) $\neq$ \all{das Kunststück} (le tour de force, l'acrobatie) ; \all{der Roman} (le roman) $\neq$ \all{die Romanze} (la romance, pièce lyrique) ; \all{die Handlung} (l'intrigue, l'action d'un récit) $\neq$ \all{der Handel} (le commerce).
    \item \textbf{Entraînement :} Transformez des noms en verbes (\all{die Ausstellung -- ausstellen}), des verbes en noms d'agent (\all{schreiben -- der Schriftsteller}), des adjectifs en noms (\all{kreativ -- die Kreativität}).
\end{enumerate}
\cle{Astuce} : Utilisez la couleur pour coder le genre (bleu = masculin, rouge = féminin, vert = neutre) sur vos fiches de révision.
\end{methode}

\begin{exoresolu}[Lexique : arts et littérature -- dérivation, genres, pluriels, contexte]
\textbf{Énoncé :}
\begin{enumerate}
    \item Complétez le tableau de dérivation (nom de base $\rightarrow$ verbe / adjectif / nom d'agent / nom d'action).
    \item Mettez les noms suivants au pluriel et donnez l'article : \all{der Maler, das Gemälde, der Schriftsteller, das Gedicht, die Lesung, der Verleger, das Theater, die Novelle}.
    \item Complétez les phrases avec le mot juste : \all{veröffentlichen, inszenieren, rezensieren, ausstellen, uraufführen}.
    \item Version : Traduisez en allemand : « Le jeune écrivain a publié son premier roman chez un grand éditeur. La critique l'a très bien accueilli. »
\end{enumerate}

\tcblower
\textbf{Solution détaillée :}

\textbf{1. Tableau de dérivation (modèle) :}

\begin{center}
\begin{tabularx}{\linewidth}{|X|X|X|X|X|}
\hline
\rowcolor{popblueL} \textbf{Nom de base} & \textbf{Verbe} & \textbf{Adjectif} & \textbf{Nom d'agent} & \textbf{Nom d'action} \\
\hline
die Kunst & (kein Verb) & künstlerisch & der Künstler / die Künstlerin & die Kunstfertigkeit \\
\hline
das Gemälde & malen & malerisch & der Maler / die Malerin & die Malerei \\
\hline
der Roman & erzählen & romanhaft & der Romanschriftsteller & das Erzählen \\
\hline
das Gedicht & dichten & poetisch & der Dichter / die Dichterin & die Dichtung \\
\hline
die Lesung & vorlesen & -- & der Vorleser / die Vorleserin & das Vorlesen \\
\hline
\end{tabularx}
\end{center}

\textbf{2. Pluriels (avec articles) :}
\begin{itemize}
    \item \all{der Maler -- die Maler} (pluriel inchangé, pas d'Umlaut)
    \item \all{das Gemälde -- die Gemälde} (pluriel inchangé)
    \item \all{der Schriftsteller -- die Schriftsteller} (pluriel inchangé)
    \item \all{das Gedicht -- die Gedichte} (pluriel en \textbf{-e}, sans Umlaut)
    \item \all{die Lesung -- die Lesungen} (féminin en \all{-ung} $\rightarrow$ \all{-ungen})
    \item \all{der Verleger -- die Verleger} (masculin en \all{-er} $\rightarrow$ inchangé)
    \item \all{das Theater -- die Theater} (pluriel inchangé)
    \item \all{die Novelle -- die Novellen} (féminin en \all{-e} $\rightarrow$ \all{-n})
\end{itemize}

\textbf{3. Complétion (contexte : la création artistique) :}
\begin{enumerate}
    \item Der Verleger will den Roman im Herbst \textbf{veröffentlichen}. (rendre public)
    \item Der Regisseur wird das Drama nächstes Jahr am Staatstheater \textbf{inszenieren}. (mettre en scène)
    \item Die Kritikerin wird das Buch in der Zeitung \textbf{rezensieren}. (faire la critique)
    \item Die Galerie plant, die neuen Werke im Frühjahr \textbf{auszustellen}. (exposer ; infinitif avec \all{zu} : particule séparable)
    \item Die Oper wird nächste Woche \textbf{uraufgeführt}. (passif : création mondiale)
\end{enumerate}

\textbf{4. Version (thème) :}
\textit{Analyse :} « Le jeune écrivain » (sujet) $\rightarrow$ \all{der junge Schriftsteller}. « a publié » $\rightarrow$ prétérit de récit : \all{veröffentlichte} (ou Perfekt : \all{hat ... veröffentlicht}). « son premier roman » (COD) $\rightarrow$ \all{seinen ersten Roman} (accusatif masculin). « chez un grand éditeur » $\rightarrow$ \all{bei einem großen Verlag} (datif après \all{bei}). « La critique » $\rightarrow$ \all{die Kritik}. « l'a très bien accueilli » $\rightarrow$ \all{nahm ihn sehr gut auf} (verbe à particule \all{aufnehmen} : \all{nahm ... auf}).

\textbf{Proposition (prétérit -- style récit) :}
\all{Der junge Schriftsteller veröffentlichte seinen ersten Roman bei einem großen Verlag. Die Kritik nahm ihn sehr gut auf.}

\textbf{Proposition (Perfekt -- style journalistique) :}
\all{Der junge Schriftsteller hat seinen ersten Roman bei einem großen Verlag veröffentlicht. Die Kritik hat ihn sehr gut aufgenommen.}
\end{exoresolu}

\begin{methode}[Grammaire : les participes présent et passé en épithète (Partizipialattribute)]
\textbf{Objectif :} Transformer une subordonnée relative en épithète participiale (style soutenu, densité informative) et l'accorder correctement.
\begin{enumerate}
    \item \textbf{Identifier la relative :} \all{der Mann, der liest} $\rightarrow$ antécédent \all{der Mann}, verbe \all{liest}.
    \item \textbf{Participe présent (Partizip I) : action active, simultanée :} infinitif + \cle{-d} : \all{lesend}. \all{der lesende Mann} = « l'homme qui lit ».
    \item \textbf{Participe passé (Partizip II) : action subie (passif) ou achevée :} verbes forts : \all{gelesen} ; verbes faibles : \all{gemacht}. \all{der gelesene Brief} = « la lettre qui a été lue ».
    \item \textbf{Accord (déclinaison) :} Le participe se comporte comme un \cle{adjectif épithète}. Sa terminaison dépend de l'article, du cas, du genre et du nombre :
    \begin{itemize}
        \item Article défini (\all{der, die, das}) $\rightarrow$ déclinaison \cle{faible} (\all{-e, -en}).
        \item Article indéfini (\all{ein, kein}) ou possessif $\rightarrow$ déclinaison \cle{mixte} (\all{-er, -e, -es, -en}).
        \item Sans article $\rightarrow$ déclinaison \cle{forte} (\all{-er, -e, -es, -en}).
    \end{itemize}
    \item \textbf{Position :} \cle{Toujours entre l'article et le nom} : \all{das von Kritikern gelobte Buch}.
    \item \textbf{Compléments du participe :} Ils se placent \textbf{avant} le participe : \all{das in Berlin gemalte Bild} (jamais *\all{das gemalte in Berlin Bild}).
\end{enumerate}
\cle{Repère Bac} : \all{Partizip I} = \all{-d} (actif, simultané) ; \all{Partizip II} = \all{ge-...-t/-en} (passif, achevé).
\end{methode}

\begin{exoresolu}[Grammaire : participes en épithète -- transformation et accord]
\textbf{Énoncé :}
\begin{enumerate}
    \item Transformez la relative en épithète participiale (Partizip I ou II) avec la terminaison correcte.
    \begin{enumerate}
        \item Ich sehe den Maler. Er malt das Bild. $\rightarrow$ Ich sehe den \_\_\_\_\_\_\_\_ Maler.
        \item Das ist das Gedicht. Es wurde von Goethe geschrieben. $\rightarrow$ Das ist das von Goethe \_\_\_\_\_\_\_\_ Gedicht.
        \item Wir besuchen die Ausstellung. Sie wurde gestern eröffnet. $\rightarrow$ Wir besuchen die gestern \_\_\_\_\_\_\_\_ Ausstellung.
        \item Er liest einen Roman. Der Roman fesselt ihn. $\rightarrow$ Er liest einen \_\_\_\_\_\_\_\_ Roman.
    \end{enumerate}
    \item Traduisez en allemand (attention à la place des compléments) :
    \begin{enumerate}
        \item Un livre écrit par un auteur célèbre.
        \item Une exposition visitée hier par beaucoup de gens.
        \item L'artiste qui peint le portrait (en ce moment).
    \end{enumerate}
\end{enumerate}

\tcblower
\textbf{Solution détaillée :}

\textbf{1. Transformation et accord :}

\begin{itemize}
    \item[a)] \all{Ich sehe den Maler. Er malt das Bild.}
    \textit{Analyse :} le sujet de la relative \all{der Maler} est agent actif $\rightarrow$ \textbf{Partizip I} (\all{malend}).
    \textit{Cas :} \all{den Maler} = accusatif masculin, article défini $\rightarrow$ déclinaison faible $\rightarrow$ \textbf{-en}.
    \textbf{Réponse :} \all{Ich sehe den \textbf{malenden} Maler.}

    \item[b)] \all{Das ist das Gedicht. Es wurde von Goethe geschrieben.}
    \textit{Analyse :} \all{das Gedicht} subit l'action (passif) $\rightarrow$ \textbf{Partizip II} (\all{geschrieben}).
    \textit{Cas :} nominatif neutre, article défini \all{das} $\rightarrow$ faible $\rightarrow$ \textbf{-e}.
    \textbf{Réponse :} \all{Das ist das von Goethe \textbf{geschriebene} Gedicht.}

    \item[c)] \all{Wir besuchen die Ausstellung. Sie wurde gestern eröffnet.}
    \textit{Analyse :} passif $\rightarrow$ \textbf{Partizip II} (\all{eröffnet}, verbe faible sans \all{ge-} car préfixe \all{er-}).
    \textit{Cas :} accusatif féminin, article défini \all{die} $\rightarrow$ faible $\rightarrow$ \textbf{-e}.
    \textbf{Réponse :} \all{Wir besuchen die gestern \textbf{eröffnete} Ausstellung.}

    \item[d)] \all{Er liest einen Roman. Der Roman fesselt ihn.}
    \textit{Analyse :} le roman « captive » : sens actif $\rightarrow$ \textbf{Partizip I} (\all{fesselnd} = captivant).
    \textit{Cas :} \all{einen Roman} = accusatif masculin, article indéfini $\rightarrow$ mixte $\rightarrow$ \textbf{-en}.
    \textbf{Réponse :} \all{Er liest einen \textbf{fesselnden} Roman.}
\end{itemize}

\textbf{2. Traduction (thème) :}

\begin{itemize}
    \item[a)] « Un livre écrit par un auteur célèbre. »
    \textit{Structure :} article + complément d'agent (\all{von einem berühmten Autor}) + Partizip II (\all{geschrieben}) + nom (\all{Buch}).
    \textbf{Réponse :} \all{\textbf{ein von einem berühmten Autor geschriebenes Buch}} (mixte, neutre : \all{-es}).

    \item[b)] « Une exposition visitée hier par beaucoup de gens. »
    \textit{Ordre des compléments :} temps (\all{gestern}) -- agent (\all{von vielen Leuten}) -- participe.
    \textbf{Réponse :} \all{\textbf{eine gestern von vielen Leuten besuchte Ausstellung}} (mixte, féminin : \all{-e}).

    \item[c)] « L'artiste qui peint le portrait (en ce moment). »
    \textit{Action active et simultanée} $\rightarrow$ Partizip I ; complément d'objet \all{das Porträt} placé avant le participe.
    \textbf{Réponse :} \all{\textbf{der das Porträt malende Künstler}} (faible, masculin nominatif : \all{-e}).
\end{itemize}

\cle{Point clé :} En allemand, le groupe participial \textbf{englobe tous les compléments} et se place \textbf{avant} le nom : \all{Artikel -- [compléments] -- Partizip -- Nomen}. C'est l'inverse du français (« le tableau peint à Berlin » $\rightarrow$ \all{das in Berlin gemalte Bild}).
\end{exoresolu}

\begin{methode}[Grammaire : le génitif (der Genitiv) -- forme, emplois, substituts]
\textbf{Objectif :} Maîtriser la déclinaison du génitif (noms, articles, adjectifs), les prépositions qui le régissent et l'alternative \all{von} + datif.
\begin{enumerate}
    \item \textbf{Formation des noms (marqueurs du génitif singulier) :}
    \begin{itemize}
        \item \textbf{Masculin / neutre :} \cle{-(e)s} : \all{des Vaters, des Buches, des Landes, des Lehrers}. \all{-es} surtout pour les monosyllabes et les noms en \all{-s, -ß, -x, -z} (\all{des Preises, des Fleißes}) ; \all{-s} pour les polysyllabes (\all{des Computers, des Malers}).
        \item \textbf{Féminin / pluriel :} \cle{pas de marqueur} sur le nom : \all{der Mutter, der Bücher}. Seuls l'article et l'adjectif marquent le cas.
        \item \textbf{Noms faibles masculins} (\all{der Student, der Mensch, der Junge, der Herr}) : \cle{-n/-en} à tous les cas sauf le nominatif singulier : \all{des Studenten, des Menschen}. Cas particulier : \all{der Name $\rightarrow$ des Namens}.
        \item \textbf{Noms propres :} \cle{-s} sans apostrophe, placés \textbf{avant} le nom : \all{Goethes Werke, Marias Buch}.
    \end{itemize}
    \item \textbf{Déclinaison des articles et adjectifs au génitif :}
    \begin{center}
    \begin{tabular}{|l|l|l|l|}
    \hline
    \rowcolor{popblueL} & \textbf{Masculin} & \textbf{Féminin} & \textbf{Neutre} \\
    \hline
    \textbf{Article défini} & \all{des guten Mannes} & \all{der guten Frau} & \all{des guten Kindes} \\
    \hline
    \textbf{Article indéfini} & \all{eines guten Mannes} & \all{einer guten Frau} & \all{eines guten Kindes} \\
    \hline
    \textbf{Négation} & \all{keines guten Mannes} & \all{keiner guten Frau} & \all{keines guten Kindes} \\
    \hline
    \textbf{Pluriel} & \multicolumn{3}{l|}{\all{der guten Männer / Frauen / Kinder} (article \all{der}, adjectif en \all{-en})} \\
    \hline
    \end{tabular}
    \end{center}
    \item \textbf{Emplois principaux :}
    \begin{itemize}
        \item \textbf{Complément du nom} (possession, appartenance, relation) : \all{das Buch \textbf{des Lehrers}}. Position : \textbf{après} le nom déterminé (sauf noms propres : \all{Marias Buch}).
        \item \textbf{Prépositions + génitif (niveau Tle) :} \all{während, wegen, statt (anstatt), trotz, innerhalb, außerhalb}. \textit{Note :} à l'oral familier on entend le datif (\all{wegen dem Wetter}), mais \textbf{à l'écrit et au Bac : génitif exigé}.
    \end{itemize}
    \item \textbf{Substitution par \all{von} + datif} (langue courante, ou obligatoire quand le nom est seul sans article) : \all{das Buch von meinem Vater} ; \all{die Werke von Shakespeare}.
\end{enumerate}
\cle{Piège francophone} : Ne pas confondre \all{des} (génitif singulier masculin/neutre de l'article défini) avec le français « des » (pluriel). \textit{Test :} remplacer par \all{eines} : si c'est possible, c'est le génitif singulier.
\end{methode}

\begin{exoresolu}[Grammaire : le génitif -- déclinaison, prépositions, réécriture]
\textbf{Énoncé :}
\begin{enumerate}
    \item Déclinez au génitif singulier et pluriel : \all{der junge Maler}, \all{das berühmte Gemälde}, \all{die alte Schriftstellerin}, \all{der kritische Verleger}.
    \item Complétez avec la préposition correcte (+ génitif) : \all{während, wegen, trotz, statt, außerhalb}.
    \begin{enumerate}
        \item \_\_\_\_\_\_\_\_ des Regens fiel das Konzert aus.
        \item \_\_\_\_\_\_\_\_ der Pause las er ein Gedicht.
        \item \_\_\_\_\_\_\_\_ seiner Krankheit kam er zur Vernissage.
        \item \_\_\_\_\_\_\_\_ eines Romans schrieb er eine Novelle.
        \item Das Atelier liegt \_\_\_\_\_\_\_\_ der Stadt.
    \end{enumerate}
    \item Réécrivez avec \all{von} + datif quand c'est naturel, ou gardez le génitif :
    \begin{enumerate}
        \item die Entscheidung des Komitees
        \item das Atelier des Künstlers
        \item die Werke Shakespeares
        \item die Eröffnung der Ausstellung
    \end{enumerate}
    \item Thème : Traduisez : « Malgré la fatigue des artistes, l'exposition a eu lieu. » (Utiliser \all{trotz} + génitif.)
\end{enumerate}

\tcblower
\textbf{Solution détaillée :}

\textbf{1. Déclinaison au génitif (singulier / pluriel) :}

\begin{center}
\begin{tabularx}{\linewidth}{|X|X|X|}
\hline
\rowcolor{popblueL} \textbf{Nominatif} & \textbf{Génitif singulier} & \textbf{Génitif pluriel} \\
\hline
der junge Maler & \textbf{des jungen Malers} (masc. : \all{-s} sur le nom, adjectif faible \all{-en}) & \textbf{der jungen Maler} (article \all{der}, adjectif \all{-en}) \\
\hline
das berühmte Gemälde & \textbf{des berühmten Gemäldes} (neutre : \all{-s}, adjectif \all{-en}) & \textbf{der berühmten Gemälde} \\
\hline
die alte Schriftstellerin & \textbf{der alten Schriftstellerin} (fém. : pas de \all{-s} sur le nom) & \textbf{der alten Schriftstellerinnen} (nom en \all{-nen}) \\
\hline
der kritische Verleger & \textbf{des kritischen Verlegers} (masc. : \all{-s}) & \textbf{der kritischen Verleger} \\
\hline
\end{tabularx}
\end{center}

\textbf{2. Prépositions + génitif :}
\begin{enumerate}
    \item \textbf{Wegen} \all{des Regens} fiel das Konzert aus. (cause ; \all{der Regen $\rightarrow$ des Regens})
    \item \textbf{Während} \all{der Pause} las er ein Gedicht. (simultanéité ; féminin \all{der Pause})
    \item \textbf{Trotz} \all{seiner Krankheit} kam er zur Vernissage. (concession ; féminin \all{seiner Krankheit})
    \item \textbf{Statt} \all{eines Romans} schrieb er eine Novelle. (substitution ; \all{eines} + \all{-s} sur le nom masculin)
    \item Das Atelier liegt \textbf{außerhalb} \all{der Stadt}. (localisation ; féminin \all{der Stadt})
\end{enumerate}

\textbf{3. Substitution \all{von} + datif / génitif conservé :}
\begin{enumerate}
    \item \all{die Entscheidung des Komitees} $\rightarrow$ \all{die Entscheidung \textbf{von dem Komitee}} (\all{vom Komitee}) -- plus fluide à l'oral.
    \item \all{das Atelier des Künstlers} $\rightarrow$ \all{das Atelier \textbf{von dem Künstler}} (\all{vom Künstler}) -- possible à l'oral ; à l'écrit soutenu, garder le génitif.
    \item \all{die Werke Shakespeares} $\rightarrow$ \textbf{garder le génitif} : \all{\textbf{Shakespeares Werke}} (nom propre placé devant) ou \all{die Werke \textbf{von Shakespeare}}.
    \item \all{die Eröffnung der Ausstellung} $\rightarrow$ \all{die Eröffnung \textbf{von der Ausstellung}} -- possible, mais le génitif reste préférable à l'écrit.
\end{enumerate}

\textbf{4. Thème : « Malgré la fatigue des artistes, l'exposition a eu lieu. »}
\textit{Analyse :} « malgré » = \all{trotz} + génitif. « la fatigue » = \all{die Müdigkeit} (féminin $\rightarrow$ \all{der Müdigkeit}). « des artistes » = génitif pluriel \all{der Künstler}. « a eu lieu » = \all{fand statt} (prétérit, verbe à particule \all{stattfinden}).

\textbf{Phrase complète :}
\all{\textbf{Trotz der Müdigkeit der Künstler fand die Ausstellung statt.}} (inversion sujet-verbe après le complément placé en tête)
ou : \all{Die Ausstellung fand trotz der Müdigkeit der Künstler statt.}
\end{exoresolu}

\begin{methode}[Conjugaison : le prétérit des verbes forts (das Präteritum der starken Verben)]
\textbf{Objectif :} Conjuguer correctement les verbes forts au prétérit (récit, texte journalistique et littéraire) et identifier les classes d'alternance vocalique.
\begin{enumerate}
    \item \textbf{Terminaisons (pas de terminaison à la 1\up{re} et 3\up{e} personne du singulier !) :}
    \begin{center}
    \begin{tabular}{|c|c|c|}
    \hline
    \rowcolor{popblueL} \textbf{Personne} & \textbf{Terminaison} & \textbf{Exemple \all{schreiben}} \\
    \hline
    ich & \textbf{--} (zéro) & \all{schrieb} \\
    du & \textbf{-st} & \all{schriebst} \\
    er/sie/es & \textbf{--} (zéro) & \all{schrieb} \\
    wir & \textbf{-en} & \all{schrieben} \\
    ihr & \textbf{-t} & \all{schriebt} \\
    sie/Sie & \textbf{-en} & \all{schrieben} \\
    \hline
    \end{tabular}
    \end{center}
    \item \textbf{Principales classes d'alternance (Ablautreihen) -- infinitif / prétérit / participe II :}
    \begin{itemize}
        \item \textbf{ei -- ie -- ie} : \all{bleiben -- blieb -- geblieben} ; \all{schreiben -- schrieb -- geschrieben}.
        \item \textbf{ie -- o -- o} : \all{fliegen -- flog -- geflogen} ; \all{verlieren -- verlor -- verloren} ; \all{bieten -- bot -- geboten}.
        \item \textbf{e -- a -- o} : \all{sprechen -- sprach -- gesprochen} ; \all{nehmen -- nahm -- genommen} ; \all{treffen -- traf -- getroffen} ; \all{helfen -- half -- geholfen} ; \all{werfen -- warf -- geworfen}.
        \item \textbf{i -- a -- u} : \all{finden -- fand -- gefunden} ; \all{singen -- sang -- gesungen} ; \all{trinken -- trank -- getrunken} ; \all{beginnen -- begann -- begonnen} ; \all{gelingen -- gelang -- gelungen}.
        \item \textbf{e -- a -- e} : \all{essen -- aß -- gegessen} ; \all{geben -- gab -- gegeben} ; \all{lesen -- las -- gelesen} ; \all{sehen -- sah -- gesehen} ; \all{geschehen -- geschah -- geschehen}.
        \item \textbf{a -- u -- a} : \all{fahren -- fuhr -- gefahren} ; \all{tragen -- trug -- getragen} ; \all{schlagen -- schlug -- geschlagen} ; \all{wachsen -- wuchs -- gewachsen}.
        \item \textbf{Autres verbes fréquents :} \all{gehen -- ging -- gegangen} ; \all{stehen -- stand -- gestanden} ; \all{kommen -- kam -- gekommen} ; \all{heißen -- hieß -- geheißen} ; \all{sein -- war -- gewesen} ; \all{haben -- hatte -- gehabt} ; \all{werden -- wurde -- geworden}.
    \end{itemize}
    \item \textbf{Verbes mixtes} (alternance + terminaisons faibles) : \all{bringen -- brachte -- gebracht} ; \all{denken -- dachte -- gedacht} ; \all{kennen -- kannte -- gekannt} ; \all{nennen -- nannte -- genannt} ; \all{wissen -- wusste -- gewusst}.
    \item \textbf{Verbes modaux au prétérit (sans Umlaut !) :} \all{können $\rightarrow$ konnte} ; \all{müssen $\rightarrow$ musste} ; \all{wollen $\rightarrow$ wollte} ; \all{sollen $\rightarrow$ sollte} ; \all{dürfen $\rightarrow$ durfte} ; \all{mögen $\rightarrow$ mochte}.
    \item \textbf{Méthode d'apprentissage :} Apprenez par \textbf{blocs de trois} (infinitif -- prétérit -- participe II) : \all{schreiben -- schrieb -- geschrieben}.
\end{enumerate}
\cle{Alerte Bac} : Ne confondez pas \all{Präteritum} (récit écrit) et \all{Perfekt} (oral). Dans une rédaction ou un commentaire au Bac, le \textbf{prétérit est obligatoire} pour \all{sein, haben}, les modaux et les verbes forts du récit.
\end{methode}

\begin{exoresolu}[Conjugaison : prétérit des verbes forts -- application textuelle]
\textbf{Énoncé :}
\begin{enumerate}
    \item Conjuguez au prétérit (toutes les personnes) : \all{finden}, \all{sprechen}, \all{lesen}, \all{werfen}.
    \item Complétez le texte au prétérit (verbes entre parenthèses) :

    \all{Gestern \_\_\_\_\_\_\_\_ (gehen) ich in die neue Galerie. Dort \_\_\_\_\_\_\_\_ (sehen) ich ein Bild, das ich sehr interessant \_\_\_\_\_\_\_\_ (finden). Der Künstler \_\_\_\_\_\_\_\_ (stehen) daneben. Er \_\_\_\_\_\_\_\_ (erzählen) mir die Geschichte des Bildes. Früher \_\_\_\_\_\_\_\_ (schreiben) er Romane, aber dann \_\_\_\_\_\_\_\_ (beginnen) er zu malen. Das Bild \_\_\_\_\_\_\_\_ (gelingen) ihm gut. Ich \_\_\_\_\_\_\_\_ (kaufen) es nicht, aber ich \_\_\_\_\_\_\_\_ (denken) noch lange daran.}
    \item Thème : Mettez au prétérit : « L'écrivain a lu son livre. Il a trouvé le public attentif. Il a donné une interview. »
\end{enumerate}

\tcblower
\textbf{Solution détaillée :}

\textbf{1. Conjugaison complète :}

\begin{center}
\begin{tabular}{|l|c|c|c|c|c|c|}
\hline
\rowcolor{popblueL} \textbf{Verbe} & \textbf{ich} & \textbf{du} & \textbf{er/sie/es} & \textbf{wir} & \textbf{ihr} & \textbf{sie/Sie} \\
\hline
\textbf{finden} & fand & fandest & fand & fanden & fandet & fanden \\
\hline
\textbf{sprechen} & sprach & sprachst & sprach & sprachen & spracht & sprachen \\
\hline
\textbf{lesen} & las & lasest & las & lasen & laset & lasen \\
\hline
\textbf{werfen} & warf & warfst & warf & warfen & warft & warfen \\
\hline
\end{tabular}
\end{center}

\textbf{2. Texte à trous (analyse verbe par verbe) :}
\begin{itemize}
    \item \all{gehen} (fort) $\rightarrow$ \textbf{ging} (ich).
    \item \all{sehen} (fort, \all{e -- a -- e}) $\rightarrow$ \textbf{sah} (ich).
    \item \all{finden} (fort, \all{i -- a -- u}) $\rightarrow$ \textbf{fand} (ich) ; attention à la subordonnée relative : \all{das ich sehr interessant \textbf{fand}} (verbe à la fin).
    \item \all{stehen} (fort : \all{stehen -- stand -- gestanden}) $\rightarrow$ \textbf{stand} (er).
    \item \all{erzählen} (\textbf{faible} !) $\rightarrow$ \textbf{erzählte} (er). \textit{Piège : ne pas inventer de forme forte.}
    \item \all{schreiben} (fort, \all{ei -- ie -- ie}) $\rightarrow$ \textbf{schrieb} (er).
    \item \all{beginnen} (fort, \all{i -- a -- u}) $\rightarrow$ \textbf{begann} (er).
    \item \all{gelingen} (fort : \all{gelingen -- gelang -- gelungen}) $\rightarrow$ \textbf{gelang} ; construction impersonnelle avec datif : \all{Das Bild \textbf{gelang} ihm gut} (« il réussit bien le tableau »).
    \item \all{kaufen} (\textbf{faible}) $\rightarrow$ \textbf{kaufte} (ich).
    \item \all{denken} (\textbf{mixte} : \all{denken -- dachte -- gedacht}) $\rightarrow$ \textbf{dachte} (ich).
\end{itemize}

\textbf{Texte complet :}
\all{Gestern \textbf{ging} ich in die neue Galerie. Dort \textbf{sah} ich ein Bild, das ich sehr interessant \textbf{fand}. Der Künstler \textbf{stand} daneben. Er \textbf{erzählte} mir die Geschichte des Bildes. Früher \textbf{schrieb} er Romane, aber dann \textbf{begann} er zu malen. Das Bild \textbf{gelang} ihm gut. Ich \textbf{kaufte} es nicht, aber ich \textbf{dachte} noch lange daran.}

\textbf{3. Thème (prétérit) :}
\begin{itemize}
    \item « L'écrivain a lu son livre. » $\rightarrow$ \all{Der Schriftsteller \textbf{las} sein Buch.} (\all{lesen -- las})
    \item « Il a trouvé le public attentif. » $\rightarrow$ \all{Er \textbf{fand} das Publikum aufmerksam.} (\all{finden -- fand})
    \item « Il a donné une interview. » $\rightarrow$ \all{Er \textbf{gab} ein Interview.} (\all{geben -- gab} ; locution \all{ein Interview geben})
\end{itemize}
\end{exoresolu}

\begin{methode}[Expression écrite guidée : résumé et prise de position (Zusammenfassung und Stellungnahme)]
\textbf{Objectif :} Produire un texte cohérent (100--150 mots) : résumer un texte source, puis donner son avis argumenté.
\begin{enumerate}
    \item \textbf{Étape 1 : le résumé (Zusammenfassung) -- environ un tiers du volume.}
    \begin{itemize}
        \item Lisez le texte source, notez trois ou quatre idées forces (mots-clés).
        \item Rédigez au \textbf{présent} (norme pour le résumé), à la troisième personne.
        \item Utilisez \textbf{vos propres mots} (paraphrase), pas de longues citations.
        \item Phrases d'introduction : \all{Der Text handelt von ... / Der Autor stellt fest, dass ... / Abschließend wird gesagt, dass ...}
        \item Connecteurs : \all{zunächst, außerdem, jedoch, deshalb, schließlich}.
    \end{itemize}
    \item \textbf{Étape 2 : la prise de position (Stellungnahme) -- environ deux tiers du volume.}
    \begin{itemize}
        \item Formulez une \cle{thèse} claire : \all{Meiner Meinung nach ... / Ich bin der Ansicht, dass ... / Ich stimme (nicht) zu.}
        \item \textbf{Deux arguments} distincts + \textbf{un exemple concret} (le contexte camerounais ou germanophone est le bienvenu).
        \item Structure : \all{Erstens ... Zweitens ... Ein Beispiel dafür ist ... / Zum einen ... Zum anderen ...}
        \item Contre-argument (optionnel mais valorisé) : \all{Ein Gegenargument könnte sein ... Doch ...}
        \item Conclusion : \all{Zusammenfassend lässt sich sagen ... / Daher fordere ich ...}
    \end{itemize}
    \item \textbf{Étape 3 : vérification linguistique (checklist Bac) :}
    \begin{itemize}
        \item \cle{Cas} (accusatif / datif / génitif après prépositions et verbes).
        \item \cle{Ordre des mots} (verbe en 2\up{e} position ; verbe à la fin après \all{dass, weil, ob} ; participe en fin de phrase).
        \item \cle{Majuscules} à \textbf{tous les noms}.
        \item \cle{Accords} (sujet-verbe ; article-adjectif-nom).
        \item \cle{Orthographe} (Umlaute, \all{ß} contre \all{ss}, consonnes doubles).
    \end{itemize}
\end{enumerate}
\cle{Format Bac} : « Fassen Sie den Text zusammen und nehmen Sie Stellung. » $\rightarrow$ Deux paragraphes visuellement distincts.
\end{methode}

\begin{exoresolu}[Expression écrite : résumé et Stellungnahme -- sujet type Bac]
\textbf{Énoncé :}
\textbf{Aufgabe :} Fassen Sie den Text „Kunst und Literatur im deutschsprachigen Raum“ (voir exercice 1) zusammen (ca. 60 Wörter) und nehmen Sie Stellung zu der These: „Kulturelle Bildung ist der Schlüssel zur gesellschaftlichen Teilhabe.“ (ca. 80--100 Wörter).

\tcblower
\textbf{Solution détaillée :}

\textbf{Étape 1 : analyse du texte source (idées forces) :}
\begin{enumerate}
    \item Scène artistique variée et internationale ; Street Art reconnue (East Side Gallery).
    \item Littérature dynamique : jeunes auteurs, thèmes identité/migration, prix et festivals.
    \item Financement précaire, aggravé par la pandémie ; aides d'urgence mais insécurité structurelle.
\end{enumerate}

\textbf{Étape 2 : rédaction du résumé (présent, troisième personne, ton neutre) :}
\all{Der Text thematisiert die Vielfalt der zeitgenössischen Kunst- und Literaturszene im deutschsprachigen Raum. Street Art hat sich als anerkannte Kunstform etabliert, zum Beispiel die East Side Gallery in Berlin. Junge Schriftsteller behandeln Themen wie Migration und Identität und werden durch Preise und Festivals gefördert. Trotz dieses Aufschwungs bleibt die finanzielle Lage vieler Künstler prekär, verschärft durch die Corona-Pandemie. Staatliche Hilfsprogramme lindern die Not, lösen aber das strukturelle Problem nicht.}
(\textit{environ 65 mots -- conforme.})

\textbf{Étape 3 : rédaction de la prise de position (Ich-Form, argumentatif) :}
\all{Meiner Meinung nach ist kulturelle Bildung tatsächlich der Schlüssel zur gesellschaftlichen Teilhabe. Erstens ermöglicht sie den Zugang zu Sprache und Kultur, was für Integration und beruflichen Erfolg unerlässlich ist. Wer Literatur versteht und Kunst rezipieren kann, entwickelt kritisches Denken und Empathie. Zweitens stärkt kulturelle Teilhabe das demokratische Bewusstsein: Man lernt, verschiedene Perspektiven auszuhalten und Konflikte mit Worten auszutragen. Ein Beispiel sind Leseclubs und Schultheater, die benachteiligten Jugendlichen in Deutschland -- und auch in Kamerun -- eine Stimme geben. Ein Gegenargument wäre die Kostenfrage, doch Bildung ist eine Investition, keine Ausgabe. Zusammenfassend fordere ich, dass kulturelle Bildung nicht als Luxus, sondern als Grundrecht im Lehrplan verankert wird.}
(\textit{environ 105 mots -- structure claire : thèse, argument 1, argument 2, exemple, contre-argument, conclusion.})

\textbf{Étape 4 : checklist linguistique (auto-correction) :}
\begin{itemize}
    \item \all{die Teilhabe} (féminin) -- correct ; \all{unerlässlich} -- correct.
    \item \all{aushalten} (verbe à particule : \all{hält ... aus}) -- particule en fin de phrase, correct.
    \item \all{verankert wird} (passif présent : \all{werden} + Partizip II) -- correct.
    \item Majuscules vérifiées sur tous les noms : \all{Bildung, Schlüssel, Teilhabe, Sprache, Kultur, Integration, Erfolg, Literatur, Kunst, Denken, Empathie, Bewusstsein, Perspektiven, Konflikte, Beispiel, Leseclubs, Schultheater, Jugendlichen, Stimme, Gegenargument, Kostenfrage, Investition, Ausgabe, Grundrecht, Lehrplan}.
    \item Subordonnées (\all{dass, was, wer}) : verbe conjugué à la fin -- vérifié.
\end{itemize}
\end{exoresolu}

\begin{attention}[Les 5 erreurs qui coûtent des points au Bac (AL22)]
\begin{enumerate}
    \item \textbf{Majuscules oubliées sur les noms :} *\all{die kunst}, *\all{der roman}, *\all{die literatur} $\rightarrow$ \textbf{TOUS les noms prennent la majuscule} : \all{die Kunst, der Roman, die Literatur}. \textit{Astuce :} relisez votre copie \textbf{une fois uniquement} pour les majuscules.
    \item \textbf{Génitif mal formé (ou remplacé par le datif) :} *\all{das Buch von der Lehrer} $\rightarrow$ \all{das Buch \textbf{des Lehrers}} ; *\all{wegen dem Wetter} (oral familier) $\rightarrow$ \all{wegen \textbf{des Wetters}} (écrit, Bac) ; oubli du \textbf{-s/-es} sur les noms masculins et neutres : *\all{des Vater} $\rightarrow$ \all{des Vater\textbf{s}}. \textit{Règle :} \all{während, wegen, trotz, statt, außerhalb} = \textbf{génitif obligatoire} à l'écrit.
    \item \textbf{Participes en épithète : accord et place.} *\all{das gemalte in Berlin Bild} (ordre français) $\rightarrow$ \all{das \textbf{in Berlin gemalte} Bild} ; oubli de la terminaison : *\all{ein geschrieben Buch} $\rightarrow$ \all{ein \textbf{geschriebenes} Buch} (mixte, neutre) ; confusion Partizip I / II : *\all{der gefesselte Roman} $\rightarrow$ \all{der \textbf{fesselnde} Roman} (le roman captive : sens actif).
    \item \textbf{Prétérit des verbes forts :} *\all{er fandet} $\rightarrow$ \all{er \textbf{fand}} ; *\all{er schreibte} $\rightarrow$ \all{er \textbf{schrieb}} ; *\all{ich ginge} (Konjunktiv II !) $\rightarrow$ \all{ich \textbf{ging}}. \textbf{Apprenez les trois formes par cœur} (infinitif -- prétérit -- participe II). Modaux sans Umlaut : \all{musste, konnte, wollte} (\all{müsste, könnte} = Konjunktiv II, autre sens !).
    \item \textbf{Faux amis et calques syntaxiques :}
    \begin{itemize}
        \item \all{aktuell} = « actuel, d'actualité » ; « actuellement » = \all{derzeit, zurzeit}.
        \item \all{kontrollieren} = « vérifier » ; « contrôler/diriger » = \all{steuern, leiten}.
        \item \all{realisieren} = « se rendre compte » ; « réaliser (un projet) » = \all{verwirklichen, umsetzen}.
        \item « Je suis d'accord » $\rightarrow$ \all{Ich \textbf{stimme zu}} (verbe à particule \all{zustimmen} + datif : \all{Ich stimme \textbf{dir} zu}).
        \item « Il me manque » $\rightarrow$ \all{\textbf{Er fehlt mir}} (datif !), jamais *\all{ich fehle ihn}.
    \end{itemize}
\end{enumerate}
\cle{Conseil final} : Relisez votre production \textbf{à l'envers} (de la dernière phrase vers la première) pour ne voir que la forme, pas le sens.
\end{attention}

\newpage

\coursec{Exercices}

\courssub{Je vérifie mes connaissances}

\begin{exercice}[QCM : Vocabulaire et grammaire]
\textbf{Choisissez la bonne réponse. Une seule proposition est correcte.}
\begin{enumerate}
    \item Le mot \all{der Maler} signifie :
    \begin{enumerate}
        \item le musée
        \item le peintre
        \item la peinture
        \item l'exposition
    \end{enumerate}
    \item Quel est le pluriel correct de \all{das Gedicht} ?
    \begin{enumerate}
        \item die Gedichten
        \item die Gedichte
        \item die Gedichts
        \item die Gedichte
    \end{enumerate}
    \item La forme du participe présent de \all{schreiben} utilisée comme épithète devant un nom au nominatif masculin est :
    \begin{enumerate}
        \item schreibend
        \item schreibende
        \item schreibender
        \item schreibendes
    \end{enumerate}
    \item Au génitif singulier, le nom \all{der Dichter} (le poète) devient :
    \begin{enumerate}
        \item des Dichters
        \item des Dichters
        \item des Dichter
        \item den Dichters
    \end{enumerate}
    \item Le prétérit de \all{lesen} (il/elle lisait) est :
    \begin{enumerate}
        \item las
        \item lasen
        \item lies
        \item gelesen
    \end{enumerate}
\end{enumerate}
\end{exercice}

\begin{exercice}[Vrai ou Faux ? Justifiez en allemand]
\textbf{Décidez si les affirmations suivantes sont vraies (Richtig) ou fausses (Falsch) et justifiez votre réponse par une phrase complète en allemand.}
\begin{enumerate}
    \item \all{Das Partizip II von „gehen“ ist „gegangen“.}
    \item \all{Der Genitiv Plural von „die Bilder“ lautet „der Bilder“.}
    \item \all{Ein „Roman“ ist ein kurzes Gedicht.}
    \item \item \all{Das Partizip I wird gebildet mit der Endung „-t“.}
    \item \all{Im Präteritum ändern starke Verben oft ihren Stammvokal.}
\end{enumerate}
\end{exercice}

\begin{exercice}[Vocabulaire : Familles de mots, articles et pluriels]
\textbf{Complétez le tableau suivant. Donnez l'article défini, le pluriel et, si possible, un verbe ou un adjectif de la même famille.}
\begin{center}
\begin{tabularx}{\linewidth}{|X|X|X|X|}
\hline
\rowcolor{popblueL} \textbf{Nom (singulier)} & \textbf{Article} & \textbf{Pluriel} & \textbf{Famille de mots (Verbe / Adjectif)} \\ \hline
\all{Kunst} & & & \all{künstlerisch} \\ \hline
& \all{der} & \all{Schriftsteller} & \all{schreiben} \\ \hline
\all{Gemälde} & & & \all{malen} \\ \hline
\all{Buch} & & & \\ \hline
\all{Ausstellung} & & & \all{ausstellen} \\ \hline
\end{tabularx}
\end{center}
\end{exercice}

\begin{exercice}[Conjugaison : Prétérit des verbes forts]
\textbf{Mettez les verbes entre parenthèses au prétérit (Präteritum). Attention aux verbes forts !}
\begin{enumerate}
    \item Gestern \all{gehen} wir ins Museum. (wir)
    \item Der Schriftsteller \all{schreiben} einen neuen Roman. (er)
    \item Ich \all{lesen} das Buch in einer Nacht. (ich)
    \item Die Maler \all{kommen} aus Österreich. (sie)
    \item Er \all{finden} das Gemälde sehr schön. (er)
    \item Wir \all{sehen} die Ausstellung in Berlin. (wir)
    \item Das Kind \all{schlafen} während der Lesung. (es)
    \item Sie \all{geben} dem Künstler ein Geschenk. (sie)
\end{enumerate}
\end{exercice}

\courssub{Je m'entraîne}

\begin{exercice}[Traduction : Participes en épithète]
\textbf{Traduisez les phrases suivantes en allemand en utilisant le participe présent ou passé comme épithète (attributif). Pensez aux déclinaisons !}
\begin{enumerate}
    \item Le peintre qui travaille dans l'atelier est célèbre.
    \item Le roman écrit par Goethe est connu dans le monde entier.
    \item Nous visitons une exposition intéressante. (participe présent)
    \item Les tableaux peints par cet artiste sont chers.
    \item L'écrivain vivant à Berlin lit ce soir. (participe présent)
\end{enumerate}
\end{exercice}

\begin{exercice}[Transformation : Relatives $\rightarrow$ Participes / Génitif]
\textbf{Transformez les phrases en remplaçant la proposition relative par un participe en épithète (Partizip I ou II) ou en utilisant le génitif.}
\begin{enumerate}
    \item Der Maler, \all{der das Bild malt}, ist jung. $\rightarrow$ Der \ldots Maler ist jung.
    \item Das ist das Buch, \all{das der Schriftsteller geschrieben hat}. $\rightarrow$ Das ist das \ldots Buch.
    \item Die Werke \all{von dem Dichter} sind berühmt. $\rightarrow$ \ldots Werke sind berühmt. (Génitif)
    \item Die Frau, \all{die das Gedicht liest}, ist meine Lehrerin. $\rightarrow$ Die \ldots Frau ist meine Lehrerin.
    \item Das Atelier \all{von dem Künstler} ist groß. $\rightarrow$ \ldots Atelier ist groß. (Génitif)
\end{enumerate}
\end{exercice}

\begin{exercice}[Texte à trous : Le génitif]
\textbf{Mettez les noms entre parenthèses au génitif (article + nom + terminaison si nécessaire).}
\begin{textebox}[Lückentext: Der Genitiv]
Die Werke \_\_\_\_\_\_\_\_\_\_\_\_ (der Dichter) sind in der Bibliothek.
Wir bewundern die Farben \_\_\_\_\_\_\_\_\_\_\_\_ (das Gemälde).
Der Name \_\_\_\_\_\_\_\_\_\_\_\_ (der Maler) steht auf der Liste.
Die Lesung \_\_\_\_\_\_\_\_\_\_\_\_ (der Schriftstellerin) beginnt um 20 Uhr.
Das Atelier \_\_\_\_\_\_\_\_\_\_\_\_ (der Künstlerin) ist hell.
\end{textebox}
\end{exercice}

\begin{exercice}[Réécriture au prétérit : Biographie]
\textbf{Réécrivez le paragraphe suivant au prétérit (Präteritum). Tous les verbes sont forts ou irréguliers.}
\begin{textebox}[Präsens $\rightarrow$ Präteritum]
Heute \all{kommt} der Schriftsteller nach Köln. Er \all{geht} in sein Hotel und \all{schreibt} an seinem neuen Roman. Am Abend \all{liest} er in einer Buchhandlung. Viele Leute \all{kommen} und \all{hören} zu. Er \all{trinkt} danach ein Glas Wein und \all{denkt} an seine Kindheit.
\end{textebox}
\end{exercice}

\courssub{Je me prépare au Bac}

\begin{exercice}[Compréhension écrite : Texte original — „Ein Abend im Literaturhaus“]
\textbf{Lisez le texte suivant et répondez aux questions en allemand.}
\begin{textebox}[Ein Abend im Literaturhaus]
Gestern Abend öffnete das Literaturhaus in München seine Türen für eine besondere Lesung. Der bekannte Schriftsteller Thomas Berger, dessen neuer Roman „Schatten der Zeit“ gerade erschienen war, las aus seinem Werk. Der Saal war bis auf den letzten Platz gefüllt; die Zuhörer hingen gebannt an seinen Lippen. Bergers Stimme, ruhig und tief, ließ die Figuren lebendig werden. Besonders der Abschnitt über die Kindheit des Protagonisten in einer kleinen Stadt am Rhein berührte viele Anwesende. Nach der Lesung signierte der Autor geduldig die mitgebrachten Bücher und beantwortete Fragen zum Schreibprozess. Ein Besucher meinte: „Diese Geschichten gehen unter die Haut.“ Der Abend klang bei einem Glas Wein im Foyer aus, wo sich Literaturfreunde noch lange über das Gehörte unterhielten.
\end{textebox}

\textbf{A. Vrai ou Faux ? Justifiez par une phrase du texte (copiez le segment exact).}
\begin{enumerate}
    \item Thomas Berger a lu un vieux roman.
    \item La salle était à moitié vide.
    \item L'auteur a signé des livres après la lecture.
    \item Le protagoniste a grandi à Munich.
\end{enumerate}

\textbf{B. Questions ouvertes (répondez en allemand, phrases complètes).}
\begin{enumerate}
    \item Wie heißt der neue Roman von Thomas Berger?
    \item Welche Passage hat die Zuhörer besonders berührt?
    \item Was hat der Autor nach der Lesung gemacht?
    \item Wie war die Atmosphäre im Foyer nach der Veranstaltung?
\end{enumerate}

\textbf{C. Travail linguistique (sur le texte).}
\begin{enumerate}
    \item Trouvez dans le texte deux exemples de \textbf{Partizip II en épithète} et donnez leur infinitif.
    \item Trouvez un exemple de \textbf{génitif} et analysez-le (nom + article + terminaison).
    \item Conjuguez le verbe \all{erscheinen} (apparaître/paraître) au prétérit aux 3e personne du singulier et du pluriel.
\end{enumerate}
\end{exercice}

\begin{exercice}[Thème et Version]
\textbf{A. Thème : Traduisez en allemand. (Utilisez le vocabulaire de la leçon : participes épithètes, génitif, prétérit forts.)}
\begin{enumerate}
    \item Le poète qui vit à Vienne écrit des poèmes célèbres.
    \item Les tableaux du peintre (génitif) sont dans le musée.
    \item Hier, l'écrivain est allé à la bibliothèque et a lu un vieux manuscrit.
    \item Nous avons vu une exposition intéressante sur l'art moderne.
    \item Le roman écrit par elle est un best-seller.
\end{enumerate}

\textbf{B. Version : Traduisez en français le paragraphe suivant.}
\begin{textebox}[Versionsaufgabe]
Das Atelier des berühmten Malers liegt in der Altstadt. Die dort gemalten Bilder zeigen Landschaften am Rhein. Ein Besucher, der die Ausstellung betrachtet, fühlt die Ruhe der Natur. Der Künstler, der seit zwanzig Jahren hier arbeitet, hat viele Preise gewonnen. Seine Werke sind in der ganzen Welt bekannt.
\end{textebox}
\end{exercice}

\begin{exercice}[Production écrite : Mon artiste préféré]
\textbf{Rédigez un texte cohérent d'environ 10 lignes (80--100 mots) en allemand sur le sujet suivant :}
\all{„Schreiben Sie einen kurzen Text über Ihren Lieblingsschriftsteller oder Lieblingsmaler. Beschreiben Sie sein bekanntestes Werk und erklären Sie, warum Sie ihn mögen.“}

\textbf{Consignes :}
\begin{itemize}
    \item Utilisez au moins 3 verbes forts au prétérit.
    \item Intégrez au moins 2 participes en épithète (Partizip I ou II).
    \item Utilisez au moins 2 génitifs (possession ou apposition).
    \item Structurez votre texte : Introduction -- Description de l'œuvre -- Avis personnel -- Conclusion.
    \item Soignez l'orthographe (majuscules aux noms, Umlaute, ß) et la ponctuation.
\end{itemize}
\end{exercice}

\newpage

\coursec{Corrigés des exercices}

\begin{corrige}[Exercice 1 — QCM : Vocabulaire et grammaire]
\begin{enumerate}
    \item \textbf{2. le peintre} — \all{der Maler, die Maler} désigne le peintre (artiste). Le musée est \all{das Museum, die Museen}, le tableau \all{das Gemälde, die Gemälde}, la peinture (l'activité) \all{die Malerei}, l'exposition \all{die Ausstellung, die Ausstellungen}.
    \item \textbf{2. die Gedichte} — Les noms neutres en \all{-icht} font leur pluriel en \all{-e} : \all{das Gedicht} $\rightarrow$ \all{die Gedichte}. (Les propositions 2 et 4 sont identiques, c'est la bonne réponse.)
    \item \textbf{3. schreibender} — Participe présent : \all{schreibend}. Devant un nom masculin au nominatif après l'article indéfini \all{ein} (cas fort) : terminaison \all{-er} $\rightarrow$ \all{ein schreibender Maler}. Après l'article défini, on aurait \all{der schreibende Maler} (terminaison faible \all{-e}).
    \item \textbf{1. des Dichters} — \all{der Dichter} (n-déclinaison) : génitif singulier \all{des Dichters} (article \all{des} + nom \all{Dichter} + \all{-s}).
    \item \textbf{1. las} — Verbe fort \all{lesen} (e $\rightarrow$ a) : \all{ich las, du lasest/last, er/sie/es las, wir lasen, ihr last, sie lasen}. \all{gelesen} est le Partizip II.
\end{enumerate}
\end{corrige}

\begin{corrige}[Exercice 2 — Vrai ou Faux ? Justifiez en allemand]
\begin{enumerate}
    \item \textbf{Richtig.} \all{Das Partizip II von „gehen“ ist „gegangen“.} (Verbe fort : \all{gehen – ging – gegangen}.)
    \item \textbf{Richtig.} \all{Der Genitiv Plural von „die Bilder“ lautet „der Bilder“.} (Pluriel \all{die Bilder} ; au génitif pluriel, article \all{der}, et pas de \all{-n} ajouté car le pluriel ne se termine pas en \all{-n} ni en \all{-s}.)
    \item \textbf{Falsch.} \all{Ein Roman ist ein langer Prosatext, kein kurzes Gedicht.} (Un roman est un long texte en prose ; \all{das Gedicht} = le poème.)
    \item \textbf{Falsch.} \all{Das Partizip I wird mit der Endung „-d“ gebildet (z. B. lesend, schreibend).} (La terminaison est \all{-d}, pas \all{-t} ; la terminaison \all{-t} est celle du Partizip II des verbes faibles : \all{gemacht, gelernt}.)
    \item \textbf{Richtig.} \all{Im Präteritum ändern starke Verben oft ihren Stammvokal (Ablaut), z. B. lesen – las, schreiben – schrieb.}
\end{enumerate}
\end{corrige}

\begin{corrige}[Exercice 3 — Vocabulaire : Familles de mots, articles et pluriels]
\begin{center}
\begin{tabularx}{\linewidth}{|X|X|X|X|}
\hline
\rowcolor{popblueL} \textbf{Nom (singulier)} & \textbf{Article} & \textbf{Pluriel} & \textbf{Famille de mots (Verbe / Adjectif)} \\ \hline
\all{Kunst} & \all{die} & \all{die Künste} & \all{künstlerisch (Adj.), der Künstler, die Künstlerin} \\ \hline
\all{Schriftsteller} & \all{der} & \all{die Schriftsteller} & \all{schreiben (V.), schriftstellerisch (Adj.)} \\ \hline
\all{Gemälde} & \all{das} & \all{die Gemälde} & \all{malen (V.), malerisch (Adj.), der Maler} \\ \hline
\all{Buch} & \all{das} & \all{die Bücher} & \all{lesen (V.), der Buchstabe, die Buchhandlung} \\ \hline
\all{Ausstellung} & \all{die} & \all{die Ausstellungen} & \all{ausstellen (V.), der Aussteller} \\ \hline
\end{tabularx}
\end{center}
\emph{Remarques :} \all{die Kunst} $\rightarrow$ Umlaut au pluriel (\all{die Künste}). \all{der Schriftsteller} (masculin en \all{-er}) : pluriel inchangé, génitif \all{des Schriftstellers}. \all{das Gemälde} (neutre en \all{-e}) : pluriel inchangé. \all{das Buch} : Umlaut + \all{-er} (\all{die Bücher}). \all{die Ausstellung} (féminin en \all{-ung}) : pluriel en \all{-en}.
\end{corrige}

\begin{corrige}[Exercice 4 — Conjugaison : Prétérit des verbes forts]
\begin{enumerate}
    \item \all{gingen} (\all{gehen} : ging, gingst, ging, gingen, gingt, gingen)
    \item \all{schrieb} (\all{schreiben} : schrieb, schriebst, schrieb, schrieben, schriebt, schrieben)
    \item \all{las} (\all{lesen} : las, lasest/last, las, lasen, last, lasen)
    \item \all{kamen} (\all{kommen} : kam, kamst, kam, kamen, kamt, kamen)
    \item \all{fand} (\all{finden} : fand, fandest, fand, fanden, fandet, fanden)
    \item \all{sahen} (\all{sehen} : sah, sahst, sah, sahen, saht, sahen)
    \item \all{schlief} (\all{schlafen} : schlief, schliefst, schlief, schliefen, schlieft, schliefen)
    \item \all{gaben} (\all{geben} : gab, gabst, gab, gaben, gabt, gaben)
\end{enumerate}
\emph{Rappel :} au prétérit des verbes forts, la 1\iere{} et la 3\ieme{} personne du singulier sont identiques et sans terminaison (\all{ich ging, er ging}) ; la 2\ieme{} personne du pluriel prend \all{-t} directement sur le radical du prétérit (\all{ihr schriebt, ihr schlieft}).
\end{corrige}

\begin{corrige}[Exercice 5 — Traduction : Participes en épithète]
\begin{enumerate}
    \item \all{Der in der Werkstatt arbeitende Maler ist berühmt.} (Partizip I \all{arbeitend} + complément \all{in der Werkstatt} placé avant ; déclinaison nominatif masculin faible \all{-e} après \all{der}.)
    \item \all{Der von Goethe geschriebene Roman ist weltbekannt.} (Partizip II \all{geschrieben} + agent \all{von Goethe} ; déclinaison nominatif masculin faible \all{-e} après l'article défini \all{der}.)
    \item \all{Wir besuchen eine die Besucher interessierende Ausstellung.} (Partizip I \all{interessierend} + objet \all{die Besucher} ; déclinaison accusatif féminin faible \all{-e} après \all{eine}.)
    \item \all{Die von diesem Künstler gemalten Bilder sind teuer.} (Partizip II \all{gemalt} + agent \all{von diesem Künstler} ; déclinaison nominatif pluriel faible \all{-en} après \all{die}.)
    \item \all{Der in Berlin lebende Schriftsteller liest heute Abend.} (Partizip I \all{lebend} + complément \all{in Berlin} ; déclinaison nominatif masculin faible \all{-e} après \all{der}.)
\end{enumerate}
\emph{Point de méthode :} le participe épithète se décline comme un adjectif ; les compléments (lieu, agent, objet) se placent \textbf{avant} le participe, jamais après.
\end{corrige}

\begin{corrige}[Exercice 6 — Transformation : Relatives $\rightarrow$ Participes / Génitif]
\begin{enumerate}
    \item \all{Der das Bild malende Maler ist jung.} (Relatif sujet $\rightarrow$ Partizip I \all{malend} + objet \all{das Bild} préposé ; nominatif masculin \all{-e} après \all{der}.)
    \item \all{Das ist das von dem Schriftsteller geschriebene Buch.} (Relatif objet au passif $\rightarrow$ Partizip II \all{geschrieben} + agent \all{von dem Schriftsteller} ; nominatif neutre \all{-e} après \all{das}.)
    \item \all{Die Werke des Dichters sind berühmt.} (Génitif singulier masculin : \all{des Dichters}. La forme \all{des Dichters Werke} est possible mais soutenue ; on préfère \all{die Werke des Dichters}.)
    \item \all{Die das Gedicht lesende Frau ist meine Lehrerin.} (Relatif sujet $\rightarrow$ Partizip I \all{lesend} + objet \all{das Gedicht} ; nominatif féminin \all{-e} après \all{die}.)
    \item \all{Das Atelier des Künstlers ist groß.} (Génitif singulier masculin : \all{des Künstlers}.)
\end{enumerate}
\end{corrige}

\begin{corrige}[Exercice 7 — Texte à trous : Le génitif]
\begin{enumerate}
    \item \all{des Dichters} (masculin : article \all{des} + \all{Dichter} + \all{-s})
    \item \all{des Gemäldes} (neutre : article \all{des} + \all{Gemälde} + \all{-s})
    \item \all{des Malers} (masculin : article \all{des} + \all{Maler} + \all{-s})
    \item \all{der Schriftstellerin} (féminin : article \all{der} + \all{Schriftstellerin} sans \all{-s})
    \item \all{der Künstlerin} (féminin : article \all{der} + \all{Künstlerin} sans \all{-s})
\end{enumerate}
\textbf{Texte complet :} \all{Die Werke des Dichters sind in der Bibliothek. Wir bewundern die Farben des Gemäldes. Der Name des Malers steht auf der Liste. Die Lesung der Schriftstellerin beginnt um 20 Uhr. Das Atelier der Künstlerin ist hell.}
\end{corrige}

\begin{corrige}[Exercice 8 — Réécriture au prétérit : Biographie]
\textbf{Texte au prétérit :}

\all{Gestern kam der Schriftsteller in Köln an. Er ging in sein Hotel und schrieb an seinem neuen Roman. Am Abend las er in einer Buchhandlung. Viele Leute kamen und hörten zu. Er trank danach ein Glas Wein und dachte an seine Kindheit.}

\emph{Remarques :}
\begin{itemize}
    \item Verbes forts : \all{kam ... an} (\all{ankommen}, particule \all{an} en fin de phrase), \all{ging} (\all{gehen}), \all{schrieb} (\all{schreiben}), \all{las} (\all{lesen}), \all{kamen} (\all{kommen}), \all{trank} (\all{trinken}).
    \item Verbes faibles ou irréguliers : \all{hörten} (\all{hören}, faible), \all{dachte} (\all{denken}, irrégulier mixte : \all{denken – dachte – gedacht}).
    \item Attention à \all{ankommen} : verbe à particule séparable, la particule \all{an} va en fin de phrase : \all{Der Schriftsteller kam in Köln an.}
\end{itemize}
\end{corrige}

\begin{corrige}[Exercice 9 — Compréhension écrite : „Ein Abend im Literaturhaus“]
\textbf{A. Vrai ou Faux ?}
\begin{enumerate}
    \item \textbf{Falsch.} Justification : \all{„… dessen neuer Roman ‚Schatten der Zeit‘ gerade erschienen war“} (le roman vient de paraître, ce n'est pas un vieux roman).
    \item \textbf{Falsch.} Justification : \all{„Der Saal war bis auf den letzten Platz gefüllt“} (la salle était comble, pas à moitié vide).
    \item \textbf{Richtig.} Justification : \all{„Nach der Lesung signierte der Autor geduldig die mitgebrachten Bücher“}.
    \item \textbf{Falsch.} Justification : \all{„… Kindheit des Protagonisten in einer kleinen Stadt am Rhein“} (pas à Munich, mais dans une petite ville sur le Rhin).
\end{enumerate}

\textbf{B. Questions ouvertes}
\begin{enumerate}
    \item \all{Der neue Roman von Thomas Berger heißt „Schatten der Zeit“.}
    \item \all{Der Abschnitt über die Kindheit des Protagonisten in einer kleinen Stadt am Rhein hat die Zuhörer besonders berührt.}
    \item \all{Nach der Lesung signierte der Autor geduldig die mitgebrachten Bücher und beantwortete Fragen zum Schreibprozess.}
    \item \all{Die Atmosphäre im Foyer war gesellig und entspannt; die Literaturfreunde unterhielten sich bei einem Glas Wein noch lange über das Gehörte.}
\end{enumerate}

\textbf{C. Travail linguistique}
\begin{enumerate}
    \item \textbf{Participes passés en épithète trouvés :}
    \begin{itemize}
        \item \all{der bekannte Schriftsteller} (infinitif : \all{kennen}) — épithète de \all{Schriftsteller}.
        \item \all{die mitgebrachten Bücher} (infinitif : \all{mitbringen}) — épithète de \all{Bücher}.
        \item \all{das Gehörte} (infinitif : \all{hören}) — Partizip II substantivé (neutre, « ce qui avait été entendu »).
    \end{itemize}
    \item \textbf{Génitif :} \all{des Protagonisten} (dans \all{die Kindheit des Protagonisten}).
    \emph{Analyse :} nom \all{der Protagonist} (masculin, n-déclinaison) ; article au génitif singulier \all{des} ; terminaison \all{-en} (n-déclinaison : \all{der Protagonist, des Protagonisten, dem Protagonisten, den Protagonisten}).
    \item \textbf{Prétérit de \all{erscheinen} (fort : erscheinen – erschien – ist erschienen) :} 3\ieme{} pers. sg. \all{erschien} ; 3\ieme{} pers. pl. \all{erschienen}.
\end{enumerate}

\textbf{Barème indicatif (sur 20) :} A (4 pts) + B (8 pts) + C (8 pts).
\end{corrige}

\begin{corrige}[Exercice 10 — Thème et Version]
\textbf{A. Thème (propositions)}
\begin{enumerate}
    \item \all{Der in Wien lebende Dichter schreibt berühmte Gedichte.} (Partizip I \all{lebend} + complément de lieu \all{in Wien} placé avant.)
    \item \all{Die Gemälde des Malers sind im Museum.} (Génitif \all{des Malers}.)
    \item \all{Gestern ging der Schriftsteller in die Bibliothek und las ein altes Manuskript.} (Prétérits forts : \all{ging} [\all{gehen}], \all{las} [\all{lesen}].)
    \item \all{Wir besuchten eine die Besucher interessierende Ausstellung über moderne Kunst.} (Prétérit faible \all{besuchten} ; Partizip I \all{interessierend} + objet \all{die Besucher} préposé.)
    \item \all{Der von ihr geschriebene Roman ist ein Bestseller.} (Partizip II \all{geschrieben} + agent \all{von ihr}.)
\end{enumerate}

\textbf{B. Version}
\begin{quote}
« L'atelier du peintre célèbre se trouve dans la vieille ville. Les tableaux peints là-bas montrent des paysages du Rhin. Un visiteur qui regarde l'exposition ressent le calme de la nature. L'artiste, qui travaille ici depuis vingt ans, a gagné de nombreux prix. Ses œuvres sont connues dans le monde entier. »
\end{quote}
\emph{Notes de traduction :} \all{des berühmten Malers} = du peintre célèbre (génitif). \all{dort gemalten} = peints là-bas (Partizip II épithète). \all{der die Ausstellung betrachtet} = qui regarde l'exposition (relatif sujet). \all{der seit zwanzig Jahren hier arbeitet} = qui travaille ici depuis vingt ans (\all{seit} + présent pour une action commencée dans le passé et qui continue). \all{Seine Werke} = ses œuvres.
\end{corrige}

\begin{corrige}[Exercice 11 — Production écrite : Mon artiste préféré]
\textbf{Proposition de rédaction modèle (environ 95 mots)}

\all{Mein Lieblingsschriftsteller ist Erich Kästner. Der in Dresden geborene Autor lebte lange in Berlin und schrieb viele bekannte Bücher für Kinder und Erwachsene. Sein berühmtestes Werk ist „Emil und die Detektive“. Der von ihm geschriebene Roman handelt von einem Jungen, der sein gestohlenes Geld mit Hilfe anderer Kinder wiederfindet. Ich las dieses Buch als Kind und mochte den Humor des Autors sehr. Kästners Sprache ist klar, direkt und nie langweilig. Der Witz seiner Figuren gefiel mir schon immer. Ich lese seine Werke bis heute gern wieder.}

\textbf{Vérification des consignes :}
\begin{itemize}
    \item \textbf{Au moins 3 verbes forts au prétérit :} \all{schrieb} (\all{schreiben}), \all{las} (\all{lesen}), \all{gefiel} (\all{gefallen}) ; s'y ajoutent les irréguliers \all{mochte} (\all{mögen}) et \all{lebte} (\all{leben}, faible).
    \item \textbf{2 participes en épithète :} \all{der in Dresden geborene Autor} (Partizip II \all{geboren} + complément de lieu), \all{der von ihm geschriebene Roman} (Partizip II \all{geschrieben} + agent).
    \item \textbf{2 génitifs :} \all{den Humor des Autors} (génitif de possession), \all{Kästners Sprache} (génitif de nom propre en \all{-s}, sans article).
    \item \textbf{Structure :} introduction (nom, origine), présentation de l'œuvre (\all{Emil und die Detektive}), avis personnel (humour, langue), conclusion (relecture).
    \item \textbf{Orthographe :} majuscules aux noms, Umlaute (\all{Kästner, Bücher, gefiel}), \all{ß} respecté (\all{gestohlenes}).
\end{itemize}

\textbf{Points de vigilance pour l'élève :}
\begin{itemize}
    \item Ne pas confondre \all{leben} (faible : \all{lebte}) et \all{lesen} (fort : \all{las}).
    \item Partizip I : terminaison \all{-d} (\all{lebend, schreibend}) ; déclinaison adjectivale obligatoire.
    \item Partizip II des verbes forts : \all{ge-} + racine (souvent avec Ablaut) + \all{-en} ; l'agent est introduit par \all{von} + datif.
    \item Génitif : article \all{des} (masculin/neutre) ou \all{der} (féminin/pluriel) ; \all{-s} ou \all{-es} au nom masculin et neutre singulier, jamais au féminin ni au pluriel.
    \item Ordre des mots : les compléments (lieu, agent, objet) se placent avant le participe épithète (\all{in Dresden geborene}, \all{von ihm geschriebene}).
\end{itemize}
\end{corrige}

\newpage

\coursec{Fiche bilan}

\courssub{Carte mentale de la leçon}
\begin{popfigure}
\centering
\begin{tikzpicture}[
    mindmap,
    grow cyclic,
    every node/.style={concept, execute at begin node=\hskip0pt},
    concept color=popblue,
    level 1/.append style={level distance=4.5cm, sibling angle=360/7},
    level 2/.append style={level distance=3cm, sibling angle=45},
    text=white,
    font=\footnotesize\bfseries
]
\node[concept, scale=1.2, align=center]{AL22\\Arts \& Littérature} [clockwise from=30]
    child[concept color=popblue] { node[concept, align=center]{Texte support\\„Kunst und Literatur“} }
    child[concept color=poporange] { node[concept, align=center]{Lexique\\die Kunst, der Roman,\\das Gedicht\ldots} }
    child[concept color=popgreen] { node[concept, align=center]{Grammaire 1\\Participe présent/passé\\en épithète} }
    child[concept color=poppurple] { node[concept, align=center]{Grammaire 2\\Génitif\\des/der/des} }
    child[concept color=poppink] { node[concept, align=center]{Conjugaison\\Präteritum verbes forts\\schrieb, las, malte} }
    child[concept color=popgold] { node[concept, align=center]{Expression\\Résumé, questions,\\production guidée} }
    child[concept color=popdark] { node[concept, align=center]{Méthode\\Commentaire structuré\\Introduction/Corps/Conclusion} };
\end{tikzpicture}
\legende{Vue d'ensemble de la leçon AL22 : texte, lexique, deux points de grammaire, conjugaison, expression et méthode.}
\end{popfigure}

\courssub{Bilan synthétique}
\begin{bilan}

\textbf{1. Lexique clé (die Kunst und die Literatur)}
\begin{center}
\begin{tabularx}{\linewidth}{|>{\bfseries}l|X|X|}
\hline
\rowcolor{popblueL} \textbf{Allemand (Art. + Pl.)} & \textbf{Français} & \textbf{Famille / Note} \\ \hline
\all{die Kunst, die Künste} & l'art, les arts & \all{Künstler} (artiste), \all{künstlerisch} (artistique) \\ \hline
\all{der Maler, die Maler} & le peintre & \all{malen} (peindre), \all{das Gemälde, die Gemälde} (le tableau) \\ \hline
\all{der Schriftsteller, die Schriftsteller} & l'écrivain & \all{schreiben} (écrire), \all{die Schrift} (l'écriture) \\ \hline
\all{der Roman, die Romane} & le roman & \all{der Romanautor} (romancier) \\ \hline
\all{das Gedicht, die Gedichte} & le poème & \all{dichten} (composer des vers), \all{der Dichter} (le poète) \\ \hline
\all{die Literatur} & la littérature & \all{literarisch} (littéraire), \all{der Literaturkritiker} \\ \hline
\all{die Ausstellung, die Ausstellungen} & l'exposition & \all{ausstellen} (exposer), \all{besuchen} (visiter) \\ \hline
\all{das Theater, die Theater} & le théâtre & \all{das Stück, die Stücke} (la pièce), \all{aufführen} (jouer/représenter) \\ \hline
\all{die Musik} & la musique & \all{der Musiker}, \all{das Konzert, die Konzerte} \\ \hline
\all{lesen, las, gelesen} & lire & verbe fort ; \all{der Leser}, \all{die Lesung} (lecture publique) \\ \hline
\end{tabularx}
\end{center}

\textbf{2. Règles de grammaire à connaître par cœur}
\begin{center}
\begin{tabularx}{\linewidth}{|>{\bfseries}p{3.5cm}|X|X|}
\hline
\rowcolor{poporangeL} \textbf{Notion} & \textbf{Règle / Formation} & \textbf{Exemple traduit} \\ \hline
\textbf{Participe I (Présent) en épithète} & Infinitif + \cle{-d} $\rightarrow$ déclinaison comme adj. & \all{ein lesender Junge} « un garçon qui lit / lisant » \\ \hline
\textbf{Participe II (Passé) en épithète} & \cle{ge-} + radical + \cle{-t/en} (+ \cle{ge-} si préfixe inséparable) $\rightarrow$ déclinaison & \all{das gelesene Buch} « le livre lu » ; \all{ein geschriebenes Gedicht} « un poème écrit » \\ \hline
\textbf{Génitif (cas de l'appartenance)} & \textbf{Masc./Neutre fort :} \cle{des} + N \cle{-(e)s} \quad \textbf{Fém./Pluriel :} \cle{der} / \cle{der} (pas de -s) & \all{der Roman des Schriftstellers} (le roman de l'écrivain) ; \all{die Farbe des Gemäldes} ; \all{der Name der Malerin} \\ \hline
\textbf{Präteritum verbes forts} & Radical modifié (Ablaut) + terminaisons \cle{-, -st, -, -en, -t, -en} \textbf{SANS -te-} & \all{schreiben $\rightarrow$ schrieb} ; \all{lesen $\rightarrow$ las} ; \all{malen $\rightarrow$ malte} (faible) \\ \hline
\end{tabularx}
\end{center}

\textbf{3. Méthodes express}
\begin{itemize}
    \item \textbf{Commentaire de texte :} 1) \textit{Introduction} (accroche, présentation source/sujet, annonce du plan). 2) \textit{Corps} (2--3 parties : résumé structuré $\rightarrow$ analyse thèmes/vocabulaire/grammaire $\rightarrow$ avis personnel/ouverture). 3) \textit{Conclusion} (bilan + ouverture culturelle).
    \item \textbf{Traduction Thème (Fr $\rightarrow$ De) :} Identifier le verbe $\rightarrow$ Cas (sujet=Nom, COD=Acc, COI=Dat, possession=Gén) $\rightarrow$ Ordre (V2 phrase principale, verbe final subordonnée) $\rightarrow$ Accords (adj, participes).
    \item \textbf{Production orale/écrite :} Utiliser \cle{Weil} (subordonnée) pour justifier, \cle{Obwohl} pour nuancer, participes en épithète pour densifier (\all{Das von mir gelesene Buch war spannend}).
\end{itemize}

\end{bilan}

\courssub{Checklist avant le Bac}
\begin{aretenir}[Checklist avant le Bac]
$\square$ Je connais le vocabulaire de base (Kunst, Maler, Roman, Gedicht, Schriftsteller) avec articles et pluriels. \\
$\square$ Je sais former le \textbf{Participe I} (\all{lesend}) et l'accorder en épithète (\all{der lesende Schüler}). \\
$\square$ Je sais former le \textbf{Participe II} des verbes forts/faibles (\all{gelesen, gemalt, geschrieben}) et l'accorder (\all{das geschriebene Buch}). \\
$\square$ Je décline correctement le \textbf{Génitif} : \all{des Vaters, der Mutter, des Kindes, der Eltern} (articles + terminaisons noms M/N). \\
$\square$ Je conjugue au \textbf{Präteritum} les verbes forts du thème : \all{schreiben (schrieb)}, \all{lesen (las)}, \all{finden (fand)}, \all{geben (gab)}, \all{sehen (sah)}. \\
$\square$ Je distingue \textbf{verbes forts} (voyelle change, pas de -te-) et \textbf{verbes faibles} (malte, machte). \\
$\square$ Je sais structurer un \textbf{commentaire de texte} en 3 parties (Introduction, Corps, Conclusion) avec connecteurs logiques (\all{Zunächst, Außerdem, Abschließend}). \\
$\square$ Je peux rédiger un \textbf{résumé} en allemand (Présent ou Prétérit) en 5--6 phrases max. \\
$\square$ Je sais exprimer mon avis sur une œuvre : \all{Meiner Meinung nach ist der Roman... weil...}. \\
$\square$ Je repère les \textbf{faux amis} : \all{Roman} $\neq$ romain (c'est un roman), \all{Kunst} $\neq$ kunst (art), \all{lesen} $\neq$ leçon (lire). \\
$\square$ Je gère l'\textbf{ordre des mots} : V2 (phrase principale), verbe final (subordonnée \all{weil, dass, obwohl}), place du participe II en fin de groupe verbal. \\
$\square$ Je connais au moins un auteur germanophone classique (Goethe, Kafka, Hesse) et un contemporain (Kehlmann, Erpenbeck) pour l'ouverture culturelle.
\end{aretenir}

\findecours

\end{document}
